\documentclass[11pt,a4paper]{article}

\usepackage[T1]{fontenc}
\usepackage{lmodern}
\usepackage[a4paper,margin=23mm,includeheadfoot,headheight=15pt,headsep=15pt,footskip=22pt]{geometry}
\usepackage{setspace,titlesec,etoolbox,tocloft,bookmark}
\usepackage{microtype}
\usepackage{amsmath,amssymb,amsthm,mathtools}
\usepackage{booktabs,longtable,array,tabularx}
\usepackage{graphicx}
\usepackage{float}
\usepackage{xcolor}
\usepackage{enumitem}
\usepackage{fancyhdr}
\usepackage{hyperref}
\usepackage[nameinlink,noabbrev]{cleveref}
\usepackage{listings}
\usepackage{caption}
\usepackage{bm}
\usepackage{needspace}
\usepackage{placeins}

\graphicspath{{core_figures/}}
\allowdisplaybreaks

\hypersetup{
  pdftitle={Projector-Resolved Seed Calculus VI},
  pdfauthor={Komorebi P.},
  pdfsubject={Transfer semigroups, stochastic and Bayesian packets, fractional powers, Green operators, and PDE realizations},
  colorlinks=true,
  linkcolor=black,
  citecolor=black,
  urlcolor=black
}

\pagestyle{fancy}
\fancyhf{}
\fancyhead[L]{\sffamily\fontsize{8}{10}\selectfont\color{muted}KOMOREBI MATHEMATICS / PAPER VI}
\fancyhead[R]{\sffamily\fontsize{8}{10}\selectfont\color{muted}TRADITIONAL PAPER}
\fancyfoot[L]{\sffamily\fontsize{8}{10}\selectfont\color{muted}Transfer semigroups and memory}
\fancyfoot[R]{\sffamily\small\color{accent}\thepage}

\newtheorem{theorem}{Theorem}[section]
\newtheorem{proposition}[theorem]{Proposition}
\newtheorem{lemma}[theorem]{Lemma}
\newtheorem{corollary}[theorem]{Corollary}
\theoremstyle{definition}
\newtheorem{definition}[theorem]{Definition}
\newtheorem{example}[theorem]{Example}
\newtheorem{construction}[theorem]{Construction}
\newtheorem{conjecture}[theorem]{Conjecture}
\theoremstyle{remark}
\newtheorem{remark}[theorem]{Remark}
\newtheorem{boundary}[theorem]{Boundary}
\newtheorem{certificate}[theorem]{Numerical Certificate}
\newtheorem{classical}[theorem]{Classical Input}

\newcommand{\cA}{\mathcal A}
\newcommand{\cB}{\mathcal B}
\newcommand{\cC}{\mathcal C}
\newcommand{\cE}{\mathcal E}
\newcommand{\cF}{\mathcal F}
\newcommand{\cG}{\mathcal G}
\newcommand{\cH}{\mathcal H}
\newcommand{\cK}{\mathcal K}
\newcommand{\cL}{\mathcal L}
\newcommand{\cP}{\mathcal P}
\newcommand{\cR}{\mathcal R}
\newcommand{\cS}{\mathcal S}
\newcommand{\cT}{\mathcal T}
\newcommand{\bbC}{\mathbb C}
\newcommand{\bbR}{\mathbb R}
\newcommand{\bbZ}{\mathbb Z}
\newcommand{\bbN}{\mathbb N}
\newcommand{\bbP}{\mathbb P}
\newcommand{\Exp}{\mathbb E}
\newcommand{\diag}{\operatorname{diag}}
\newcommand{\rank}{\operatorname{rank}}
\newcommand{\Ran}{\operatorname{Ran}}
\newcommand{\Ker}{\operatorname{Ker}}
\newcommand{\Var}{\operatorname{Var}}
\newcommand{\Cov}{\operatorname{Cov}}
\newcommand{\Spec}{\operatorname{Spec}}
\newcommand{\Id}{\operatorname{Id}}
\newcommand{\Tr}{\operatorname{Tr}}
\newcommand{\norm}[1]{\left\lVert#1\right\rVert}
\newcolumntype{L}[1]{>{\raggedright\arraybackslash}p{#1}}
\newcolumntype{Y}{>{\raggedright\arraybackslash}X}

\lstset{
  language=Python,
  basicstyle=\ttfamily\footnotesize,
  breaklines=true,
  frame=single,
  showstringspaces=false,
  columns=fullflexible,
  keepspaces=true,
  numbers=left,
  numberstyle=\tiny,
  numbersep=7pt
}

\title{\textbf{Projector-Resolved Seed Calculus VI}\\[0.45em]
\large Transfer Semigroups: Stochastic, Bayesian, Fractional, and PDE Realizations}
\author{Komorebi P.}
\date{Reviewed edition: 6 September 2026\\\small Fifth Lean-checked formalization block; all 33 results covered}


\definecolor{accent}{HTML}{137668}
\definecolor{muted}{HTML}{536568}
\hypersetup{linkcolor=accent,citecolor=accent,urlcolor=accent,pdftitle={Komorebi Paper VI: Transfer Semigroups and Memory},pdfauthor={Komorebi P.}}
\frenchspacing\clubpenalty=10000\widowpenalty=10000\raggedbottom
\titleformat{\section}{\sffamily\bfseries\large\color{accent}}{\thesection}{.7em}{}
\titleformat{\subsection}{\sffamily\bfseries\normalsize}{\thesubsection}{.7em}{}
\titlespacing*{\section}{0pt}{1.4\baselineskip}{.55\baselineskip}
\pretocmd{\section}{\Needspace{7\baselineskip}}{}{}
\pretocmd{\subsection}{\Needspace{5\baselineskip}}{}{}
\setstretch{1.045}\setlength{\parskip}{.3em}\setlength{\parindent}{1.35em}
\setlength{\emergencystretch}{3em}
\setlength{\cftsecnumwidth}{2.5em}
\setlength{\cftsubsecindent}{2.5em}
\setlength{\cftsubsecnumwidth}{3.4em}
\newtheorem{standardfact}[theorem]{Standard Fact}
\crefname{standardfact}{standard fact}{standard facts}
\crefname{classical}{classical input}{classical inputs}
\crefname{boundary}{boundary}{boundaries}
\crefname{certificate}{numerical certificate}{numerical certificates}
\newcommand{\Dfun}{\mathsf D}
\newcommand{\Sfun}{\mathsf S}
\newcommand{\CC}{\mathbb C}
\newcommand{\RR}{\mathbb R}
\newcommand{\NN}{\mathbb N}
\newcommand{\Pcal}{\mathcal P}
\newcommand{\Tcal}{\mathcal T}
\newcommand{\Kcal}{\mathcal K}
\newcommand{\Bcal}{\mathcal B}
\newcommand{\Ecal}{\mathcal E}
\newcommand{\Fcal}{\mathcal F}
\newcommand{\Hcal}{\mathcal H}
\newcommand{\Lcal}{\mathcal L}
\newcommand{\Rcal}{\mathcal R}
\newcommand{\Scal}{\mathcal S}
\newcommand{\Ocal}{\mathcal O}
\newcommand{\Ob}{\operatorname{Ob}}
\newcommand{\tr}{\operatorname{tr}}
\newcommand{\supp}{\operatorname{supp}}
\newcommand{\osc}{\operatorname{osc}}
\newcommand{\id}{\mathrm{id}}
\newcommand{\e}{\mathrm e}
\newcommand{\ii}{\mathrm i}
\newcommand{\one}{\bm 1}
\newcommand{\dd}{\,\mathrm d}
\newcommand{\abs}[1]{\left\lvert#1\right\rvert}
\newcommand{\ip}[2]{\left\langle#1,#2\right\rangle}
\newcommand{\simplex}{\mathop{\rm int}\Delta}
\newcommand{\EE}{\mathbb E}
\newcommand{\PP}{\mathbb P}
\newcommand{\Acal}{\mathcal A}
\newcommand{\Ccal}{\mathcal C}
\newcommand{\Dcal}{\mathcal D}
\newcommand{\Mcal}{\mathcal M}
\newcommand{\Qcal}{\mathcal Q}
\newcommand{\Vcal}{\mathcal V}
\newcommand{\TV}{\operatorname{TV}}

% Operation Komorebi: shared publication profile.
\usepackage{komorebi-publication}
\renewcommand{\KomorebiTitle}{Transfer Semigroups and Memory}
\renewcommand{\KomorebiSubtitle}{Probabilistic refinement, fractional operators, and exact hidden-state elimination}
\renewcommand{\KomorebiSeries}{Core series VI}
\renewcommand{\KomorebiKind}{Traditional paper}
\renewcommand{\KomorebiEdition}{Standardized publication edition / 8 October 2026}
\renewcommand{\KomorebiShortTitle}{Paper VI}
\renewcommand{\KomorebiTitleSize}{38}
\renewcommand{\KomorebiTitleLeading}{43}
\renewcommand{\KomorebiProfile}{classical}
\renewcommand{\KomorebiOrientation}{portrait}
\KomorebiSetup
\begin{document}
\KomorebiCover
\KomorebiOpeningNote{\textbf{Part I.} The complete released 33-result formal core, including nonnormal memory and canonical finite holomorphic naturality.

\medskip
\textbf{Part II.} The complete strengthened mathematical development: bounded refinement towers, Bayesian information, Hermite transfer, quantitative memory compression, and competing risks.

\medskip
\textbf{Part III.} Distinct retained results from the earlier human edition, exact examples, and the bounded-operator reciprocal calculus received from Paper IV.}

\hypersetup{pageanchor=false}

\hypersetup{pageanchor=true}
\section*{Reading this edition}\addcontentsline{toc}{section}{Reading this edition}
Paper VI has several substantial source branches. This edition preserves the latest formal-core manuscript in Part I and the complete strengthened standalone manuscript in Part II. Part III retains twenty-four earlier results with distinct statements or hypotheses, five exact examples, and the holomorphic operator-substitution section migrated from Paper IV. Overlap is intentional and is not a claim of independent novelty for each proof environment.

The formal denominator is exactly the thirty-three scoped results of Part I. The inherited September 6, 2026 release records 341 explicit Lean theorem declarations across 47 proof modules and an audit of 782 namespace declarations. These are different units. No fresh Lean compilation was available for this October edition, and the formal count does not automatically include Parts II or III. The Lean guide gives the exact targets and specialization boundaries.

Part I includes three explicitly open conjectures. Part II supplies bounded or finite theorems for several related questions; it does not settle arbitrary continuum limits, unbounded-generator convergence, general statistical sufficiency, or globally optimal nonnormal model reduction. The reader should keep each result's declared space, topology, positivity, and stability hypotheses. Part I's mass-only Bayesian commutation criterion assumes at least two coarse blocks; Part II's stronger sufficiency convention also preserves within-block conditionals, so its one-block case still requires a constant likelihood.

The current checks, historical compiler evidence, and numerical regressions are recorded separately in the accompanying package. Finite checks and illustrations have bounded scopes. They do not replace proofs. The September 12 worked visual appendix, its exact models, and additional guides to the strengthened and retained material are supplied in the visual companion.

Read Part I for the common operator language and the exact released formal contract. Read Part II for bounded Banach-space development and the broader probabilistic and analytical realizations. Read Part III for product-topology limits, reversible coarse dynamics, prior-restricted and coarse filtering, killed-transfer perturbations, rational realization, and the transferred operator reciprocal family. Section numbers in Part II have prefix E; those in Part III have prefix R. Original labels and statement locations are mapped in the editable sources.

The final mathematical audit makes explicit the hypotheses needed for contour winding, positive memory decay, finite moment closure, and eventual absorption. It also clarifies summability across varying spaces. Every substantive change is recorded against preserved sources. Historical source test counts retain their historical dates.
\clearpage
\tableofcontents
\clearpage
\listoffigures
\listoftables
\clearpage
\part*{Part I\quad The released formal core}\addcontentsline{toc}{part}{Part I: The released formal core}
\begin{center}\large Projector-Resolved Seed Calculus VI\par\normalsize Transfer Semigroups: Stochastic, Bayesian, Fractional, and PDE Realizations\end{center}



\begin{abstract}
This paper develops the transfer-semigroup realization of the projector-resolved seed calculus. A seed is propagated by a one-parameter semigroup, a Markov kernel, a normalized Bayesian transfer, a fractional spectral power, a Green operator, or a potential-weighted Feynman--Kac operator. Exact packet projectors resolve every transfer into input-output blocks, and the semigroup law becomes block convolution over intermediate packets. Coarse and fine generators are compared by a refinement defect. A Duhamel identity proves that this generator-level defect is exactly the source of the corresponding semigroup defect and gives a quantitative stability bound. In the stochastic realization, point-mass seeds generate random walks and continuous-time Markov evolution; exact lumpability is the vanishing of the same refinement defect. Conditional expectation is treated as a weighted orthogonal packet projector, with nested partitions producing martingale-difference packets. Bayesian conditioning is a normalized nonlinear operator. Its independent-data composition law, packet-evidence decomposition, stability bound, conjugate updates, and continuous likelihood flow are proved explicitly.  The Bayesian refinement conjecture is resolved in finite dimensions: for strictly positive likelihoods and at least two coarse packets, exact commutation for every prior is equivalent to blockwise-constant likelihoods up to global scale, and every failure factors through an explicit information-loss packet. Fractional powers are constructed packetwise through the spectral projectors of a positive operator, and the square-root semigroup is recovered by subordination. The fractional-naturality conjecture is also resolved for bounded nonnegative self-adjoint refinement systems: the fractional intertwining defect has an exact resolvent integral, a gap-free Hölder certificate, a sharper common-gap estimate, and a Duhamel semigroup bound. Green kernels, finite-difference Poisson problems, and second-order low-mode convergence provide the PDE realization. A finite-state Feynman--Kac formula is proved by path enumeration. Exact Schur--Feshbach elimination shows that discarded latent packets return as a frequency-dependent self-energy and, in time coordinates, as a memory kernel. Survival and hazard functions are generated by transient semigroups. Nineteen code-generated figures, machine-readable results, the original thirty-seven automated tests, and seventeen added boundary-case checks accompany this reviewed edition. The inherited September 6 release records five checked Lean blocks covering all 33 scoped theorem-like results in Appendix~\ref{core:sec:lean-status}. Finite-dimensional theorems, classical infinite-dimensional extensions, numerical certificates, and open continuum claims are kept separate.
\end{abstract}

\begin{center}
\fbox{\begin{minipage}{0.93\textwidth}
\textbf{Scope and claim boundary.} The paper proves the stated finite-dimensional semigroup, packet, Markov, Bayesian, conditional-expectation, fractional spectral, Green-kernel, Feynman--Kac, Schur--Feshbach, and survival identities. It also proves quantitative refinement-defect formulas under explicit norm hypotheses, including exact transfer of a bounded generator defect to every fractional power $0<\alpha<1$ for bounded nonnegative self-adjoint generators. It does not claim that an arbitrary infinite-dimensional operator generates a semigroup, that every stochastic process admits a finite packet representation, that every posterior family is conjugate, that finite Bayesian sufficiency automatically extends to infinite disintegrations, or that every PDE continuum limit follows from a finite test. Domain, closability, sectoriality, measurability, tightness, and compactness hypotheses remain necessary in their respective infinite-dimensional settings. Classical sources are retained in the separate end-of-series reference ledger rather than reproduced in this paper.
\end{minipage}}
\end{center}



\section{The transfer-semigroup realization}\label{core:section:01}

Let $X$ be a finite-dimensional real or complex vector space, let $A\in\operatorname{End}(X)$, and let $s\in X$ be a primitive seed. The basic continuous transfer is
\[
T(t)=e^{tA},\qquad t\ge 0,
\]
while a positive operator $L$ is more naturally written
\[
H(t)=e^{-tL}.
\]
Both satisfy the semigroup law. The sign convention changes with the domain: Markov generators and evolution equations are usually written with $e^{tA}$; positive diffusion and elliptic operators are usually written with $e^{-tL}$.

The common operator chain is
\[
\boxed{
\text{seed}\longrightarrow\text{transfer}\longrightarrow\text{project}
\longrightarrow\text{condition}\longrightarrow\text{refine}
\longrightarrow\text{complete}\longrightarrow\text{reconstruct}.
}
\]

\begin{figure}[htbp]
\centering
\includegraphics[width=0.98\textwidth]{01_transfer_semigroup_architecture.pdf}
\caption{The common operator language used throughout the paper. Semigroups, resolvents, normalized Bayesian maps, and fractional powers are distinct realizations of the same transfer stage.}
\label{core:fig:architecture}
\end{figure}

\begin{definition}[Transfer-semigroup seed system]\label{core:record:001}
A finite transfer-semigroup seed system is a tuple
\[
\mathfrak S=(X,s,A,\{P_\alpha\}_{\alpha\in I},\tau),
\]
where $X$ is the state space, $s$ is the seed, $A$ is a transfer generator, the $P_\alpha$ are pairwise orthogonal idempotents summing to the identity, and $\tau$ specifies the topology in which a completion or limit is interpreted.
\end{definition}

\begin{theorem}[Finite semigroup law]\label{core:record:002}
For every matrix $A$ and all $s,t\in\bbR$,
\[
\boxed{e^{(s+t)A}=e^{sA}e^{tA}.}
\]
Moreover, $T(0)=I$ and $T'(t)=AT(t)=T(t)A$.
\end{theorem}

\begin{proof}
Both exponentials are absolutely convergent matrix power series. Since every power of $A$ commutes with every other power, the Cauchy product gives
\[
e^{sA}e^{tA}
=\sum_{m,n\ge0}\frac{s^mt^n}{m!n!}A^{m+n}
=\sum_{r\ge0}\frac{(s+t)^r}{r!}A^r.
\]
Termwise differentiation is valid in finite dimension and gives the derivative identities.
\end{proof}

\begin{boundary}[Infinite-dimensional generators]\label{core:record:003}
For an unbounded operator, writing $e^{tA}$ is not a definition unless a semigroup-generation theorem applies. The finite matrix identities in this paper remain algebraically correct models, but a continuum realization must state the operator domain, closure, growth bound, and topology of convergence.
\end{boundary}

\section{Packet blocks of a semigroup}\label{core:section:02}

Let
\[
P_\alpha P_\beta=\delta_{\alpha\beta}P_\alpha,
\qquad
\sum_{\alpha\in I}P_\alpha=I.
\]
For any operator $B$, define
\[
B_{\beta\alpha}=P_\beta BP_\alpha.
\]
Then
\[
B=\sum_{\alpha,\beta}B_{\beta\alpha}.
\]
For the semigroup, write
\[
T_{\beta\alpha}(t)=P_\beta e^{tA}P_\alpha.
\]

\begin{theorem}[Packet semigroup convolution]\label{core:record:004}
For every $s,t$ and every pair of packets $\alpha,\beta$,
\[
\boxed{
T_{\beta\alpha}(s+t)
=\sum_{\gamma\in I}T_{\beta\gamma}(s)T_{\gamma\alpha}(t).
}
\]
Thus the semigroup law becomes ordinary matrix multiplication in the packet indices.
\end{theorem}

\begin{proof}
Insert the resolution of the identity between the two semigroup factors:
\[
\begin{aligned}
P_\beta e^{(s+t)A}P_\alpha
&=P_\beta e^{sA}\left(\sum_\gamma P_\gamma\right)e^{tA}P_\alpha\\
&=\sum_\gamma
(P_\beta e^{sA}P_\gamma)(P_\gamma e^{tA}P_\alpha).
\end{aligned}
\]
\end{proof}

\begin{corollary}[Reducing packets]\label{core:record:005}
If $P_\beta AP_\alpha=0$ for every $\beta\ne\alpha$, then
\[
P_\beta e^{tA}P_\alpha=0
\quad(\beta\ne\alpha)
\]
for every $t$. Each packet evolves independently under $P_\alpha AP_\alpha$.
\end{corollary}

\begin{proof}
The hypothesis makes $A$ block diagonal. Every power $A^m$ is block diagonal with respect to the same projectors, and therefore so is the exponential series.
\end{proof}

\begin{figure}[htbp]
\centering
\includegraphics[width=0.66\textwidth]{02_packet_semigroup_blocks.pdf}
\caption{Norms of the packet blocks $P_\beta e^{tA}P_\alpha$. Off-diagonal entries quantify transfer between packets.}
\label{core:fig:packet-blocks}
\end{figure}

\begin{remark}[Path expansion]\label{core:record:006}
Expanding the exponential and inserting $I=\sum_\alpha P_\alpha$ between every pair of generators gives
\[
P_{\alpha_m}A^mP_{\alpha_0}
=
\sum_{\alpha_1,\ldots,\alpha_{m-1}}
P_{\alpha_m}AP_{\alpha_{m-1}}\cdots P_{\alpha_1}AP_{\alpha_0}.
\]
Every coefficient is therefore a sum over packet-transfer paths. The numerical library verifies that truncating this path series converges to the ordinary matrix exponential.
\end{remark}

\section{Refinement defects and the Duhamel identity}\label{core:section:03}

Let $X_f$ and $X_c$ be fine and coarse spaces, with generators $A_f$ and $A_c$, and let
\[
C:X_f\to X_c
\]
be a coarsening map. Define the generator-level refinement defect
\[
\boxed{\Delta=CA_f-A_cC.}
\]
The corresponding semigroup defect is
\[
\mathfrak D(t)=Ce^{tA_f}-e^{tA_c}C.
\]

\begin{theorem}[Exact Duhamel refinement identity]\label{core:record:007}
For every $t\ge0$,
\[
\boxed{
Ce^{tA_f}-e^{tA_c}C
=
\int_0^t e^{(t-r)A_c}\Delta e^{rA_f}\,dr.
}
\]
In particular, $\Delta=0$ implies exact semigroup intertwining at every time.
\end{theorem}

\begin{proof}
Set
\[
F(r)=e^{(t-r)A_c}Ce^{rA_f}.
\]
Then
\[
F'(r)=e^{(t-r)A_c}(CA_f-A_cC)e^{rA_f}.
\]
Integrating from $0$ to $t$ gives $F(t)-F(0)=Ce^{tA_f}-e^{tA_c}C$ and proves the formula.
\end{proof}

\begin{corollary}[Norm bound]\label{core:record:008}
Suppose
\[
\norm{e^{tA_f}}\le M_f e^{\omega_f t},
\qquad
\norm{e^{tA_c}}\le M_c e^{\omega_c t}.
\]
Then
\[
\norm{\mathfrak D(t)}
\le
M_fM_c\norm{\Delta}
\int_0^t e^{\omega_c(t-r)}e^{\omega_fr}\,dr.
\]
For contraction semigroups,
\[
\boxed{\norm{\mathfrak D(t)}\le t\norm{\Delta}.}
\]
\end{corollary}

\begin{figure}[htbp]
\centering
\includegraphics[width=0.80\textwidth]{14_semigroup_refinement_defect.pdf}
\caption{Exact generator intertwining gives a vanishing semigroup defect. A controlled perturbation produces a nonzero defect propagated by the Duhamel identity.}
\label{core:fig:semigroup-defect}
\end{figure}

\subsection{Resolvent form}

For $z$ in the common resolvent set, let
\[
R_f(z)=(zI-A_f)^{-1},
\qquad
R_c(z)=(zI-A_c)^{-1}.
\]

\begin{theorem}[Resolvent refinement identity]\label{core:record:009}
\[
\boxed{
CR_f(z)-R_c(z)C
=R_c(z)\Delta R_f(z).
}
\]
\end{theorem}

\begin{proof}
Multiply the left side by $zI-A_c$ on the left and $zI-A_f$ on the right. The result is
\[
(zI-A_c)C-C(zI-A_f)=CA_f-A_cC=\Delta.
\]
Multiplication by the two resolvents recovers the stated identity.
\end{proof}

\section{Markov seeds and stochastic transfer}\label{core:section:04}

Use column probability vectors. A finite continuous-time Markov generator $Q$ has nonnegative off-diagonal entries and column sums zero. The stochastic semigroup is
\[
K(t)=e^{tQ}.
\]

\begin{theorem}[Finite Markov semigroup]\label{core:record:010}
If $Q$ has nonnegative off-diagonal entries and column sums zero, then $e^{tQ}$ has nonnegative entries and column sums one for every $t\ge0$.
\end{theorem}

\begin{proof}
Choose $\lambda\ge\max_i(-Q_{ii})$. Then $B=Q+\lambda I$ has nonnegative entries, so $e^{tB}$ is entrywise nonnegative by its power series. Hence
\[
e^{tQ}=e^{-\lambda t}e^{tB}
\]
is nonnegative. Since $\bm 1^TQ=0$, differentiation gives
\[
\frac{d}{dt}\bm 1^Te^{tQ}=\bm 1^TQe^{tQ}=0,
\]
so $\bm 1^Te^{tQ}=\bm 1^T$.
\end{proof}

\begin{figure}[htbp]
\centering
\includegraphics[width=0.78\textwidth]{03_random_walk_transfer.pdf}
\caption{The point mass $\delta_0$ propagated by repeated symmetric random-walk transfer.}
\label{core:fig:random-walk}
\end{figure}

\begin{example}[Symmetric random walk]\label{core:record:011}
Starting from $\delta_0$ and applying the transfer
\[
T=\frac12(S+S^{-1})
\]
for $n$ steps gives
\[
\bbP(X_n=k)=
2^{-n}\binom{n}{(n+k)/2}
\]
when $k\equiv n\pmod2$, and zero otherwise. Hence
\[
\Exp[X_n]=0,
\qquad
\Var(X_n)=n.
\]
The parity packet is exact: all mass remains in the one residue class permitted by $n$.
\end{example}

\section{Markov lumping as exact refinement}\label{core:section:05}

Let $K$ be a column-stochastic fine transition matrix, let $C$ aggregate fine states into coarse packets, and let $\bar K$ be a coarse transition matrix. Define
\[
\Delta=CK-\bar KC.
\]

\begin{theorem}[Exact lumpability]\label{core:record:012}
If $\Delta=0$, then
\[
\boxed{CK^n=\bar K^nC}
\]
for every $n\ge0$. Thus fine evolution followed by coarse reconstruction equals coarse evolution from the beginning.
\end{theorem}

\begin{proof}
The result is immediate for $n=0$. If $CK^n=\bar K^nC$, then
\[
CK^{n+1}=\bar K^nCK=\bar K^{n+1}C.
\]
\end{proof}

\begin{theorem}[Finite-horizon defect formula]\label{core:record:013}
For arbitrary $\Delta=CK-\bar KC$,
\[
\boxed{
CK^n-\bar K^nC
=
\sum_{j=0}^{n-1}\bar K^{n-1-j}\Delta K^j.
}
\]
If $K$ and $\bar K$ are contractions in the chosen norm, then
\[
\norm{CK^n-\bar K^nC}\le n\norm{\Delta}.
\]
\end{theorem}

\begin{proof}
The identity follows by telescoping:
\[
\bar K^{n-1-j}CK^{j+1}-\bar K^{n-j}CK^j
=\bar K^{n-1-j}(CK-\bar KC)K^j.
\]
Summing from $j=0$ to $n-1$ leaves $CK^n-\bar K^nC$. The norm bound follows termwise.
\end{proof}

\begin{figure}[htbp]
\centering
\includegraphics[width=0.82\textwidth]{04_markov_lumping.pdf}
\caption{Exact lumping agrees at every horizon. A small generator-level or one-step defect creates a controlled coarse discrepancy.}
\label{core:fig:lumping}
\end{figure}

\section{Conditional expectation as a stochastic projector}\label{core:section:06}

Let $\Omega=\{1,\ldots,n\}$ carry strictly positive probabilities $\pi_i$, and let $\mathcal Q=\{B_1,\ldots,B_m\}$ be a partition. Define
\[
(P_{\mathcal Q}f)_i
=
\frac{\sum_{j\in B(i)}\pi_jf_j}{\sum_{j\in B(i)}\pi_j},
\]
where $B(i)$ is the block containing $i$.

\begin{theorem}[Finite conditional-expectation projector]\label{core:record:014}
With the weighted inner product
\[
\langle f,g\rangle_\pi=\sum_i\pi_if_i\overline{g_i},
\]
the operator $P_{\mathcal Q}$ satisfies
\[
P_{\mathcal Q}^2=P_{\mathcal Q},
\qquad
\langle P_{\mathcal Q}f,g\rangle_\pi
=
\langle f,P_{\mathcal Q}g\rangle_\pi.
\]
Its range is the space of functions constant on each partition block.
\end{theorem}

\begin{proof}
After one application, the result is constant on every block, so applying the same block average again changes nothing. Weighted self-adjointness follows by expanding both inner products block by block; both equal
\[
\sum_{B\in\mathcal Q}
\frac{1}{\pi(B)}
\left(\sum_{i\in B}\pi_i f_i\right)
\overline{\left(\sum_{j\in B}\pi_j g_j\right)}.
\]
\end{proof}

\begin{theorem}[Tower law and martingale packets]\label{core:record:015}
If $\mathcal Q_c$ is coarser than $\mathcal Q_f$, then
\[
\boxed{
P_cP_f=P_fP_c=P_c.
}
\]
For a nested tower $P_0\le P_1\le\cdots\le P_r$, define
\[
D_0=P_0,
\qquad
D_j=P_j-P_{j-1}.
\]
Then
\[
D_iD_j=\delta_{ij}D_j,
\qquad
P_r=\sum_{j=0}^rD_j.
\]
\end{theorem}

\begin{proof}
A coarse-block average of a fine-block average is the coarse-block average, and a fine-block average of a coarse-constant function leaves it unchanged. The difference identities follow from $P_iP_j=P_{\min(i,j)}$.
\end{proof}

\begin{figure}[htbp]
\centering
\includegraphics[width=0.86\textwidth]{08_hermite_packets.pdf}
\caption{Hermite--Gaussian packets, which diagonalize the Ornstein--Uhlenbeck semigroup. The finite numerical release verifies their Gaussian orthogonality through quadrature.}
\label{core:fig:hermite}
\end{figure}

\begin{classical}[Gaussian chaos realization]\label{core:record:016}
For the standard Gaussian measure, the normalized probabilists' Hermite polynomials form orthonormal chaos packets. The Ornstein--Uhlenbeck semigroup acts diagonally:
\[
T_tH_n=e^{-nt}H_n.
\]
Thus conditional expectation, martingale differences, and Gaussian chaos are all projector realizations of stochastic refinement.
\end{classical}

\section{Bayesian conditioning as a normalized nonlinear operator}\label{core:section:07}

Let $p=(p_i)$ be a probability vector and let $\lambda=(\lambda_i)$ be a nonnegative likelihood with positive evidence
\[
Z_\lambda(p)=\sum_i\lambda_ip_i.
\]
Define
\[
\boxed{
\cB_\lambda(p)
=
\frac{\lambda\odot p}{Z_\lambda(p)}.
}
\]
This map is nonlinear solely because of its scalar reconstruction factor.

\begin{theorem}[Independent-likelihood composition]\label{core:record:017}
For nonnegative likelihoods $\lambda$ and $\mu$, assume $Z_\lambda(p)>0$ and $Z_{\lambda\odot\mu}(p)>0$. Equivalently, both successive updates have positive evidence. Then
\[
\boxed{
\cB_\mu(\cB_\lambda(p))
=
\cB_{\lambda\odot\mu}(p).
}
\]
\end{theorem}

\begin{proof}
The first normalization appears in both numerator and denominator of the second update and cancels:
\[
\frac{\mu_i(\lambda_ip_i/Z_\lambda)}{\sum_j\mu_j\lambda_jp_j/Z_\lambda}
=
\frac{\mu_i\lambda_ip_i}{\sum_j\mu_j\lambda_jp_j}.
\]
\end{proof}

\subsection{Packet evidence and reconstruction}

Let the state indices be partitioned into packets $I_\alpha$. Define
\[
Z_\alpha=\sum_{i\in I_\alpha}\lambda_ip_i,
\qquad
m_\alpha=\frac{Z_\alpha}{\sum_\beta Z_\beta},
\qquad
q_{\alpha,i}=\begin{cases}
\lambda_ip_i/Z_\alpha,&Z_\alpha>0,\ i\in I_\alpha,\\
0,&\text{otherwise}.
\end{cases}
\]

\begin{theorem}[Bayesian packet reconstruction]\label{core:record:018}
\[
\boxed{
\cB_\lambda(p)
=
\sum_\alpha m_\alpha q_\alpha.
}
\]
Only packets with positive evidence carry a conditional probability vector. A zero-evidence packet has zero posterior mass and contributes the zero vector. The nonlinear coupling between packets occurs only through the total evidence used to normalize the packet evidences.
\end{theorem}

\begin{proof}
If $Z_\alpha=0$, nonnegativity implies $\lambda_i p_i=0$ for every $i\in I_\alpha$, so the reconstruction identity holds with both sides zero. For $Z_\alpha>0$ and $i\in I_\alpha$,
\[
m_\alpha q_{\alpha,i}
=
\frac{Z_\alpha}{Z}\frac{\lambda_ip_i}{Z_\alpha}
=
\frac{\lambda_ip_i}{Z}.
\]
The packet supports are disjoint, so summing reconstructs the full posterior.
\end{proof}

\begin{figure}[htbp]
\centering
\includegraphics[width=0.82\textwidth]{05_bayesian_packet_update.pdf}
\caption{A finite prior and posterior split into two state packets. Evidence changes both the local packet shape and the global packet mass.}
\label{core:fig:bayes-packet}
\end{figure}

\begin{theorem}[Likelihood perturbation bound]\label{core:record:019}
Let $p$ be fixed and let $\lambda,\mu\ge0$ have positive evidences $Z_\lambda,Z_\mu$. Then
\[
\boxed{
\norm{\cB_\lambda(p)-\cB_\mu(p)}_1
\le
\frac{2\sum_i|\lambda_i-\mu_i|p_i}{\min(Z_\lambda,Z_\mu)}.
}
\]
\end{theorem}

\begin{proof}
Add and subtract $\mu\odot p/Z_\lambda$ and use the triangle inequality:
\[
\norm{\frac{\lambda\odot p}{Z_\lambda}-\frac{\mu\odot p}{Z_\mu}}_1
\le
\frac{\norm{(\lambda-\mu)\odot p}_1}{Z_\lambda}
+
\norm{\mu\odot p}_1\left|\frac1{Z_\lambda}-\frac1{Z_\mu}\right|.
\]
The second term equals $|Z_\lambda-Z_\mu|/Z_\lambda$ and is bounded by the same weighted $L^1$ difference divided by $Z_\lambda$. Interchanging $\lambda$ and $\mu$ gives the minimum-evidence form.
\end{proof}

\section{Exact Bayesian refinement and information-loss packets}\label{core:section:08}

Let the fine state set $I$ be partitioned into nonempty coarse packets $I_a$, $a\in A$.  The coarsening map on probability vectors is
\[
(C_\pi p)_a=P_a=\sum_{i\in I_a}p_i.
\]
For a fine likelihood $\lambda=(\lambda_i)_{i\in I}$ and a prior $p$, define the effective coarse likelihood
\[
r_a(p)
=
\frac{\sum_{i\in I_a}\lambda_i p_i}{P_a}
\qquad(P_a>0).
\]
Its value on a zero-mass packet is irrelevant and may be set to zero.

\begin{theorem}[Exact effective-likelihood factorization]\label{core:record:020}
Whenever the fine evidence is positive,
\[
\boxed{
C_\pi\cB_\lambda(p)
=
\cB_{r(p)}(C_\pi p).}
\]
Thus the only fine information needed for the coarse posterior, beyond the coarse prior masses, is the vector of conditional packet evidences $r(p)$.
\end{theorem}

\begin{proof}
For $P=C_\pi p$ and $Z=\sum_i\lambda_ip_i=\sum_aP_ar_a(p)$,
\[
\bigl(C_\pi\cB_\lambda(p)\bigr)_a
=
\frac{\sum_{i\in I_a}\lambda_ip_i}{Z}
=
\frac{P_ar_a(p)}{\sum_bP_br_b(p)}
=
\bigl(\cB_{r(p)}(P)\bigr)_a.
\]
\end{proof}

The next result gives the exact necessary and sufficient condition for a coarse likelihood to work for every prior.

\begin{theorem}[Bayesian refinement criterion]\label{core:record:021}
Assume that the partition contains at least two coarse packets and that the fine and coarse likelihoods are strictly positive.  The identity
\[
\boxed{
C_\pi\cB_\lambda(p)
=
\cB_{\bar\lambda}(C_\pi p)}
\]
holds for every fine prior $p$ if and only if there exists one scalar $c>0$ such that
\[
\boxed{
\lambda_i=c\,\bar\lambda_a
\qquad(i\in I_a).}
\]
Equivalently, the likelihood is constant on each coarse packet, with the coarse values determined up to the globally invisible Bayesian scaling $c$.
\end{theorem}

\begin{proof}
If $\lambda_i=c\bar\lambda_a$ on $I_a$, then $P_a r_a(p)=c\bar\lambda_aP_a$ for every prior, including zero-mass packets, and the preceding theorem gives the desired identity because Bayesian normalization removes the common scalar $c$.

Conversely, fix two fine states $i,i'\in I_a$ and one state $j\in I_b$ in another coarse packet.  Vary the distribution of a fixed coarse mass inside $I_a$ between $i$ and $i'$ while leaving every coarse packet mass unchanged.  The right side depends only on those coarse masses and is therefore unchanged.  The left-side posterior mass of packet $a$ depends on the conditional average of $\lambda_i$ and $\lambda_{i'}$.  It can be independent of that variation only if $\lambda_i=\lambda_{i'}$.  Hence $\lambda$ is constant on every packet; write the value on $I_a$ as $\lambda_a$.

Now use priors supported on one state of packet $a$ and one state of packet $b$.  Equality of the two coarse posterior odds gives
\[
\frac{\lambda_a}{\lambda_b}
=
\frac{\bar\lambda_a}{\bar\lambda_b}.
\]
This holds for every pair of coarse packets, so $\lambda_a=c\bar\lambda_a$ for one common $c>0$.
\end{proof}

\begin{remark}[The one-packet boundary]\label{core:record:022}
If the coarse partition has only one packet, every posterior coarsens to the scalar $1$, so commutation is automatic and no within-packet likelihood information can be identified.  The nontrivial criterion begins with at least two coarse packets.
\end{remark}

When sufficiency fails, the defect still has an exact finite factorization.  Fix a declared coarse likelihood $\bar\lambda$, put
\[
P=C_\pi p,
\qquad
\delta(p)=r(p)-\bar\lambda,
\qquad
\bar b=\cB_{\bar\lambda}(P),
\qquad
Z_r=\langle r(p),P\rangle.
\]

\begin{theorem}[Bayesian information-loss packet]\label{core:record:023}
Assume nonnegative fine and declared coarse likelihoods, with both $Z_r>0$ and $Z_{\bar\lambda}>0$. The coarsening defect is
\[
\boxed{
C_\pi\cB_\lambda(p)
-
\cB_{\bar\lambda}(P)
=
\frac{1}{Z_r}
\left(\diag(P)-\bar bP^{\mathsf T}\right)
\delta(p).}
\]
The factor matrix annihilates $\bar\lambda$.  This records the common-scale null direction at the declared coarse likelihood; the full exact defect remains homogeneous under simultaneous likelihood rescaling because the numerator and evidence denominator scale together.  Moreover,
\[
\boxed{
\left\|C_\pi\cB_\lambda(p)-\cB_{\bar\lambda}(P)\right\|_1
\le
\frac{2\sum_aP_a|\delta_a(p)|}
{\min\{Z_r,Z_{\bar\lambda}\}}}
\]
and
\[
\sum_aP_a|\delta_a(p)|
\le
\sum_i p_i\left|\lambda_i-\bar\lambda_{\pi(i)}\right|.
\]
\end{theorem}

\begin{proof}
The effective-likelihood theorem reduces the defect to
\[
\cB_{\bar\lambda+\delta}(P)-\cB_{\bar\lambda}(P).
\]
Since $Z_r=Z_{\bar\lambda}+P^{\mathsf T}\delta$, direct subtraction gives
\[
\frac{1}{Z_r}
\left(\diag(P)-\bar bP^{\mathsf T}\right)\delta.
\]
Also,
\[
\left(\diag(P)-\bar bP^{\mathsf T}\right)\bar\lambda
=P\odot\bar\lambda-\bar bZ_{\bar\lambda}=0.
\]
The first norm bound is the likelihood perturbation theorem applied on the coarse simplex. For a zero-mass packet, both sides of the following estimate vanish. For $P_a>0$,
\[
P_a|\delta_a|
=
P_a\left|\sum_{i\in I_a}\frac{p_i}{P_a}
(\lambda_i-\bar\lambda_a)\right|
\le
\sum_{i\in I_a}p_i|\lambda_i-\bar\lambda_a|,
\]
and summing over $a$ proves the second inequality.
\end{proof}

\begin{corollary}[Nested Bayesian packet tower]\label{core:record:024}
Consider a nested finite packet tower with strictly positive likelihoods at every level. At each adjacent coarsening with at least two coarse packets, commutation for every prior is equivalent to the fine likelihood being blockwise proportional to the declared coarse likelihood, with one common positive proportionality factor. An adjacent coarsening to a single packet commutes automatically and imposes no such condition. Under these conditions, all composite coarsening squares commute.
\end{corollary}

\begin{proof}
At each adjacent coarsening with at least two coarse packets, apply the refinement criterion. For a single coarse packet, both normalized posteriors equal 1. Composition of these commuting squares gives the full tower statement.
\end{proof}

\begin{corollary}[Canonical coarsest likelihood-resolving partition]\label{core:record:025}
For a finite family of likelihoods $\{\lambda^{(y)}\}_{y\in Y}$, define
\[
i\sim j
\quad\Longleftrightarrow\quad
\lambda_i^{(y)}=\lambda_j^{(y)}
\text{ for every }y\in Y.
\]
The equivalence classes form the unique coarsest likelihood-resolving partition: every likelihood is constant on each packet. Every update therefore closes on this partition whenever its evidence is positive. If all likelihoods are strictly positive and there are at least two profile classes, this is also the unique coarsest nontrivial partition on which every update closes for every prior. The one-packet partition always closes, independently of the profiles, and is excluded from this last assertion.
\end{corollary}

\begin{proof}
Equality of likelihood profiles is an equivalence relation. Its classes are likelihood-resolving by definition, and every likelihood-resolving partition refines them. Exact closure follows by summing the weighted masses in each fiber, including zero-mass fibers. Under strict positivity, the refinement criterion implies that every nontrivial all-prior closed partition is likelihood-resolving. When every profile is identical, all partitions close; there need not be a unique coarsest nontrivial partition.
\end{proof}

\begin{figure}[htbp]
\centering
\includegraphics[width=0.92\textwidth]{17_bayesian_refinement_square.pdf}
\caption{For strictly positive likelihoods and at least two coarse packets, the Bayesian refinement square commutes for every prior exactly when the likelihood is constant on each coarsening fiber, up to one global scale.}
\label{core:fig:bayes-refinement}
\end{figure}

\begin{figure}[htbp]
\centering
\includegraphics[width=0.94\textwidth]{18_bayesian_information_loss.pdf}
\caption{When packet sufficiency fails, the exact defect factors through the difference between the conditional effective likelihood and the declared coarse likelihood.}
\label{core:fig:bayes-loss}
\end{figure}

\begin{certificate}[Bayesian refinement]\label{core:record:026}
For the finite test partition $\{0,1\}\sqcup\{2,3\}$, a blockwise-constant likelihood gave commutation residual $1.11\times10^{-16}$.  A nonconstant likelihood produced an exact defect-factorization residual $7.29\times10^{-17}$.  Its measured $L^1$ defect was $4.8787\times10^{-2}$, below the certified bound $1.1016\times10^{-1}$.  Two priors with identical coarse masses but different hidden within-packet shapes produced coarse posteriors differing by $1.2400\times10^{-1}$, which is a direct necessity witness.
\end{certificate}

\begin{boundary}[Finite exact theorem]\label{core:record:027}
The conjectured Bayesian refinement law is therefore resolved for finite packet systems and arbitrary priors.  Infinite-dimensional disintegration, posterior consistency, and stability of data-dependent partitions require their own measure-theoretic and asymptotic hypotheses.
\end{boundary}


\section{Bayesian flows, conjugacy, and hidden states}\label{core:section:09}

Fix a score vector $g$. The likelihood flow is
\[
\lambda_i(t)=e^{tg_i},
\]
and the Bayesian flow is
\[
\Phi_t(p)_i
=
\frac{e^{tg_i}p_i}{\sum_je^{tg_j}p_j}.
\]

\begin{theorem}[Bayesian nonlinear semigroup]\label{core:record:028}
\[
\boxed{
\Phi_{s+t}=\Phi_s\circ\Phi_t.
}
\]
Its generator is the centered multiplicative field
\[
\boxed{
\frac{d}{dt}p_i(t)
=
\bigl(g_i-\langle g,p(t)\rangle\bigr)p_i(t).
}
\]
\end{theorem}

\begin{proof}
The semigroup law is the likelihood-composition theorem with $e^{sg}e^{tg}=e^{(s+t)g}$. Differentiating the normalized exponential formula gives
\[
\dot p_i=g_ip_i-p_i\frac{\sum_jg_je^{tg_j}p_j(0)}{\sum_je^{tg_j}p_j(0)}.
\]
The quotient is $\langle g,p(t)\rangle$.
\end{proof}

\begin{figure}[htbp]
\centering
\includegraphics[width=0.80\textwidth]{07_bayesian_flow.pdf}
\caption{Continuous likelihood accumulation produces a nonlinear semigroup on the probability simplex.}
\label{core:fig:bayesian-flow}
\end{figure}

\begin{example}[Beta--Binomial conjugate packet]\label{core:record:029}
A $\operatorname{Beta}(a,b)$ prior and $s$ successes in $n$ trials produce
\[
\operatorname{Beta}(a+s,b+n-s).
\]
The prior exponents are additive packets: the likelihood increments them by $(s,n-s)$. In the numerical example,
\[
(a,b)=(2,3),\quad (s,n)=(7,10),
\]
so the posterior is $\operatorname{Beta}(9,6)$ with mean $0.6$.
\end{example}

\begin{figure}[htbp]
\centering
\includegraphics[width=0.80\textwidth]{06_beta_binomial_conjugacy.pdf}
\caption{A conjugate Bayesian update viewed as exact transfer of exponent packets.}
\label{core:fig:beta-conjugacy}
\end{figure}

\subsection{Hidden Markov filtering}

Let $K$ be column stochastic and let $\lambda^{(r)}$ be the likelihood associated with the $r$th observation. The filter recursion is
\[
\widehat p_r=Kp_{r-1},
\qquad
p_r=\cB_{\lambda^{(r)}}(\widehat p_r).
\]

\begin{proposition}[Prediction-update reconstruction]\label{core:record:030}
Every hidden Markov filtering step is the composition of a linear Markov transfer and a normalized nonlinear likelihood transfer. Assume a probability initial vector, a column-stochastic transition matrix, and nonnegative likelihoods with positive evidence at every step. Then every $p_r$ is a probability vector.
\end{proposition}

\begin{proof}
Column stochasticity gives $\widehat p_r\ge0$ and $\bm1^T\widehat p_r=1$. Bayesian normalization then gives $p_r\ge0$ and $\bm1^Tp_r=1$.
\end{proof}

\begin{figure}[htbp]
\centering
\includegraphics[width=0.80\textwidth]{15_hidden_markov_filter.pdf}
\caption{A hidden Markov filter alternates a linear prediction semigroup with nonlinear Bayesian normalization.}
\label{core:fig:hmm}
\end{figure}

\section{Fractional powers as spectral packet transfer}\label{core:section:10}

Let $L=L^*\ge0$ on a finite-dimensional Hilbert space. Its spectral resolution is
\[
L=\sum_j\lambda_jP_j.
\]
For $\alpha>0$, define
\[
\boxed{L^\alpha=\sum_j\lambda_j^\alpha P_j.}
\]
The associated fractional semigroup is
\[
T_\alpha(t)=e^{-tL^\alpha}
=\sum_je^{-t\lambda_j^\alpha}P_j.
\]

\begin{theorem}[Fractional packet calculus]\label{core:record:031}
For $\alpha,\beta>0$,
\[
L^\alpha L^\beta=L^{\alpha+\beta},
\]
and for $s,t\ge0$,
\[
\boxed{T_\alpha(s)T_\alpha(t)=T_\alpha(s+t).}
\]
Every spectral packet $P_j$ reduces $L^\alpha$ and $T_\alpha(t)$.
\end{theorem}

\begin{proof}
All statements hold coefficientwise on the mutually orthogonal spectral projectors. For example,
\[
L^\alpha L^\beta
=\sum_{j,k}\lambda_j^\alpha\lambda_k^\beta P_jP_k
=\sum_j\lambda_j^{\alpha+\beta}P_j.
\]
\end{proof}

\begin{figure}[htbp]
\centering
\includegraphics[width=0.80\textwidth]{09_fractional_semigroup.pdf}
\caption{Different spectral exponents produce different transfer profiles from the same point seed.}
\label{core:fig:fractional}
\end{figure}

\subsection{Square-root subordination}

\begin{theorem}[Poisson subordination]\label{core:record:032}
For $t>0$ and $L\ge0$,
\[
\boxed{
e^{-t\sqrt L}
=
\frac{t}{2\sqrt\pi}
\int_0^\infty
s^{-3/2}e^{-t^2/(4s)}e^{-sL}\,ds.
}
\]
The integral converges in operator norm in finite dimension.
\end{theorem}

\begin{proof}
By the spectral theorem it suffices to prove the scalar identity
\[
e^{-t\sqrt\lambda}
=
\frac{t}{2\sqrt\pi}
\int_0^\infty
s^{-3/2}e^{-t^2/(4s)-\lambda s}\,ds
\]
for $\lambda\ge0$. This is the classical Gaussian integral after the substitution $u=t/(2\sqrt s)$. Applying the identity to every eigenvalue and recombining the spectral projectors proves the operator formula.
\end{proof}

\subsection{Fractional naturality under a bounded refinement defect}

Let $A\ge0$ act on a fine Hilbert space, let $B\ge0$ act on a coarse Hilbert space, and let $C$ be a bounded coarsening map.  Set
\[
D=CA-BC,
\qquad
D_\alpha=CA^\alpha-B^\alpha C,
\qquad 0<\alpha<1.
\]

\begin{theorem}[Fractional intertwining theorem]\label{core:record:033}\label{core:thm:fractional-naturality}
For bounded nonnegative self-adjoint $A$ and $B$,
\[
\boxed{
D_\alpha
=
\frac{\sin(\pi\alpha)}{\pi}
\int_0^\infty
 t^\alpha(tI+B)^{-1}D(tI+A)^{-1}\,dt.
}
\]
The improper integral converges in operator norm.  Moreover,
\[
\boxed{
\norm{D_\alpha}
\le
K_\alpha\norm{C}^{1-\alpha}\norm{D}^{\alpha},
\qquad
K_\alpha=
\frac{2^{1-\alpha}\sin(\pi\alpha)}
{\pi\alpha(1-\alpha)}.
}
\]
If $A,B\ge mI$ for some $m>0$, then the stronger linear estimate holds:
\[
\boxed{
\norm{D_\alpha}
\le
\alpha m^{\alpha-1}\norm{D}.
}
\]
In particular, exact generator intertwining $CA=BC$ implies exact fractional intertwining $CA^\alpha=B^\alpha C$.
\end{theorem}

\begin{proof}
For $0<\alpha<1$, the Balakrishnan formula gives
\[
A^\alpha
=
\frac{\sin(\pi\alpha)}{\pi}
\int_0^\infty t^{\alpha-1}A(tI+A)^{-1}\,dt,
\]
and likewise for $B$.  Using
\[
A(tI+A)^{-1}=I-t(tI+A)^{-1}
\]
and the rectangular resolvent identity
\[
(tI+B)^{-1}C-C(tI+A)^{-1}
=
(tI+B)^{-1}(CA-BC)(tI+A)^{-1},
\]
the difference of the two Balakrishnan integrals becomes the displayed resolvent-defect formula.

For small $t$,
\[
\norm{CA(tI+A)^{-1}-B(tI+B)^{-1}C}
\le 2\norm C,
\]
because the two positive resolvent factors are contractions.  For large $t$, the resolvent identity gives the bound $\norm D/t$.  Splitting the integral at
\[
r=\frac{\norm D}{2\norm C}
\]
when $C$ and $D$ are nonzero yields
\[
\begin{aligned}
\norm{D_\alpha}
&\le
\frac{\sin(\pi\alpha)}{\pi}
\left(
2\norm C\int_0^r t^{\alpha-1}\,dt
+
\norm D\int_r^\infty t^{\alpha-2}\,dt
\right)\\
&=
K_\alpha\norm C^{1-\alpha}\norm D^\alpha.
\end{aligned}
\]
The zero cases are immediate.

If $A,B\ge mI$, then both resolvents have norm at most $(t+m)^{-1}$, so
\[
\norm{D_\alpha}
\le
\frac{\sin(\pi\alpha)}{\pi}\norm D
\int_0^\infty\frac{t^\alpha}{(t+m)^2}\,dt.
\]
After $t=mu$, the beta integral equals
\[
m^{\alpha-1}\frac{\alpha\pi}{\sin(\pi\alpha)},
\]
which gives the linear estimate.
\end{proof}

\begin{corollary}[Fractional semigroup refinement]\label{core:record:034}\label{core:cor:fractional-semigroup-refinement}
For $s\ge0$,
\[
\boxed{
Ce^{-sA^\alpha}-e^{-sB^\alpha}C
=
-\int_0^s
 e^{-(s-r)B^\alpha}D_\alpha e^{-rA^\alpha}\,dr.
}
\]
Hence
\[
\norm{Ce^{-sA^\alpha}-e^{-sB^\alpha}C}
\le s\norm{D_\alpha}.
\]
Under the common gap $A,B\ge mI$,
\[
\boxed{
\norm{Ce^{-sA^\alpha}-e^{-sB^\alpha}C}
\le
s e^{-sm^\alpha}\alpha m^{\alpha-1}\norm D.
}
\]
\end{corollary}

\begin{proof}
Differentiate
\[
r\longmapsto e^{-(s-r)B^\alpha}Ce^{-rA^\alpha}
\]
and integrate from $0$ to $s$.  Positive fractional semigroups are contractions; under the gap, the product of their two norm bounds is $e^{-sm^\alpha}$.
\end{proof}

\begin{corollary}[Fractional refinement towers]\label{core:record:035}
Let $D_\ell=C_\ell A_{\ell+1}-A_\ell C_\ell$ in a bounded nonnegative self-adjoint tower, with $\sup_\ell\norm{C_\ell}\le M$.  If
\[
\sum_\ell\norm{D_\ell}^{\alpha}<\infty,
\]
then the fractional defects
\[
C_\ell A_{\ell+1}^{\alpha}-A_\ell^{\alpha}C_\ell
\]
have summable operator norms. If the tower has a uniform gap $m>0$, it is enough that $\sum_\ell\norm{D_\ell}<\infty$. On every bounded time interval, the series of corresponding fractional-semigroup defect norms converges uniformly. The maps may have different domain and codomain spaces; no operator sum across unrelated spaces is asserted.
\end{corollary}

\begin{figure}[htbp]
\centering
\includegraphics[width=0.94\textwidth]{19_fractional_naturality.pdf}
\caption{A generator-level refinement defect transfers quantitatively to fractional powers and then to fractional semigroups.  The two upper curves are theorem certificates, not fitted error bars.}
\label{core:fig:fractional-naturality}
\end{figure}

\begin{boundary}[Unbounded and nonnormal generators]\label{core:record:036}
The theorem resolves fractional naturality for bounded nonnegative self-adjoint refinement systems and gives the exact estimate needed for finite discretization towers.  For unbounded positive generators, the same resolvent argument requires a declared common-domain condition and a bounded extension of $CA-BC$.  Sectorial or nonnormal generators require their own contour calculus and resolvent bounds.  Those domain questions are not inferred from the finite certificates.
\end{boundary}

\begin{boundary}[Non-self-adjoint fractional powers]\label{core:record:037}
The spectral definition above is complete for finite self-adjoint positive operators. Sectorial and nonnormal operators require a resolvent or contour functional calculus; the same packet language can be used only after those analytic hypotheses are established.
\end{boundary}

\section{Green operators and the Poisson seed}\label{core:section:11}

On the interval $(0,1)$ with Dirichlet boundary conditions, the Green kernel for $-u''=f$ is
\[
\boxed{G(x,y)=\min(x,y)-xy.}
\]
Thus the point source $\delta_y$ generates the solution kernel, and a general source is reconstructed by
\[
u(x)=\int_0^1G(x,y)f(y)\,dy.
\]

For $n$ interior points with spacing $h=1/(n+1)$, define
\[
L_h=\frac1{h^2}
\begin{pmatrix}
2&-1&&\\
-1&2&\ddots&\\
&\ddots&\ddots&-1\\
&&-1&2
\end{pmatrix}.
\]

\begin{theorem}[Exact discrete Green reconstruction at grid nodes]\label{core:record:038}
Let $x_i=ih$. Then
\[
\boxed{
\frac1h(L_h^{-1})_{ij}
=G(x_i,x_j).
}
\]
Thus the scaled discrete inverse reproduces the continuum Dirichlet Green kernel exactly at the grid nodes.
\end{theorem}

\begin{proof}
The inverse of the unscaled tridiagonal matrix has entries
\[
\frac{\min(i,j)(n+1-\max(i,j))}{n+1}.
\]
Multiplying by $h^2$ and then dividing by $h$ gives
\[
h\min(i,j)\left(1-\frac{\max(i,j)}{n+1}\right)
=\min(x_i,x_j)-x_ix_j.
\]
\end{proof}

\begin{figure}[htbp]
\centering
\includegraphics[width=0.80\textwidth]{10_poisson_green_kernel.pdf}
\caption{A section of the continuum Green kernel and the scaled discrete inverse. They agree at the grid nodes.}
\label{core:fig:green}
\end{figure}

\subsection{Low-mode convergence}

The $k$th discrete eigenvalue is
\[
\lambda_{k,h}=\frac4{h^2}\sin^2\left(\frac{k\pi h}{2}\right),
\]
whereas the continuum eigenvalue is $(k\pi)^2$.

\begin{theorem}[Second-order low-mode convergence]\label{core:record:039}
For fixed $k$,
\[
\lambda_{k,h}=(k\pi)^2+O(h^2).
\]
More precisely, for sufficiently small $h$,
\[
0\le (k\pi)^2-\lambda_{k,h}
\le \frac{(k\pi)^4}{12}h^2.
\]
\end{theorem}

\begin{proof}
For $x\ge0$, $0\le x^2-\sin^2x\le x^4/3$. Set $x=k\pi h/2$ and multiply by $4/h^2$.
\end{proof}

\begin{figure}[htbp]
\centering
\includegraphics[width=0.72\textwidth]{11_pde_refinement_convergence.pdf}
\caption{The first finite-difference eigenvalue converges with observed order approximately two.}
\label{core:fig:pde-convergence}
\end{figure}

\section{Finite-state Feynman--Kac transfer}\label{core:section:12}

Let $K$ be a column-stochastic matrix, let $V$ be a potential on the states, and let $f$ be terminal data. Define
\[
M_V=\diag(e^{-V})K^T.
\]
The backward Feynman--Kac value after $n$ steps is
\[
u_n=M_V^nf.
\]

\begin{theorem}[Discrete Feynman--Kac path formula]\label{core:record:040}
For an initial state $i$,
\[
\boxed{
(u_n)_i
=
\Exp_i\left[
\exp\left(-\sum_{r=0}^{n-1}V(X_r)\right)f(X_n)
\right].
}
\]
\end{theorem}

\begin{proof}
For $n=1$,
\[
(M_Vf)_i=e^{-V(i)}\sum_jK_{j i}f_j,
\]
which is the stated expectation. Iterating and expanding the matrix product sums over every path $(i=i_0,i_1,\ldots,i_n)$ with probability $\prod_rK_{i_{r+1},i_r}$ and weight $\prod_re^{-V(i_r)}$.
\end{proof}

\begin{figure}[htbp]
\centering
\includegraphics[width=0.78\textwidth]{12_feynman_kac_transfer.pdf}
\caption{Potential-weighted transfer values computed by the finite Feynman--Kac operator. Exhaustive path enumeration agrees with the matrix formula in the supplied tests.}
\label{core:fig:feynman-kac}
\end{figure}

\section{Exact elimination of latent packets}\label{core:section:13}

Let $X=P X\oplus Q X$ and write
\[
A=
\begin{pmatrix}
A_{PP}&A_{PQ}\\
A_{QP}&A_{QQ}
\end{pmatrix}.
\]

\begin{theorem}[Schur--Feshbach projected resolvent]\label{core:record:041}
Assume both $zI-A$ and $zI-A_{QQ}$ are invertible. Define
\[
\boxed{
F_P(z)
=zI-A_{PP}-A_{PQ}(zI-A_{QQ})^{-1}A_{QP}.
}
\]
Then
\[
\boxed{
\left.P(zI-A)^{-1}\right|_{\Ran P}=F_P(z)^{-1}.
}
\]
\end{theorem}

\begin{proof}
Solve the hidden block equation
\[
(zI-A_{QQ})x_Q=A_{QP}x_P+y_Q
\]
for $x_Q$ and substitute it into the retained equation. The resulting retained coefficient is $F_P(z)$. Invertibility of the full operator and of the hidden block implies invertibility of this Schur complement. Setting $y_Q=0$ gives the stated identity as an operator on $\Ran P$. Equivalently, this is the block inverse formula.
\end{proof}

The self-energy
\[
\Sigma_P(z)=A_{PQ}(zI-A_{QQ})^{-1}A_{QP}
\]
contains the exact effect of the eliminated packets. A bare truncation $A_{PP}$ discards this term and is generally not exact.

\subsection{Time-domain memory}

For the homogeneous system $\dot x=Ax$, variation of constants gives
\[
x_Q(t)=e^{tA_{QQ}}x_Q(0)
+\int_0^te^{(t-r)A_{QQ}}A_{QP}x_P(r)\,dr.
\]
Substitution yields
\[
\dot x_P(t)
=A_{PP}x_P(t)+A_{PQ}e^{tA_{QQ}}x_Q(0)
+\int_0^t\cK(t-r)x_P(r)\,dr,
\]
where
\[
\boxed{
\cK(t)=A_{PQ}e^{tA_{QQ}}A_{QP}.
}
\]

\begin{theorem}[Memory/self-energy correspondence]\label{core:record:042}
For finite complex block matrices, assume $\operatorname{Re}z>\operatorname{Re}\lambda$ for every $\lambda\in\sigma(A_{QQ})$. Then the following improper integral is absolutely convergent and
\[
\boxed{
\int_0^\infty e^{-zt}\cK(t)\,dt
=A_{PQ}(zI-A_{QQ})^{-1}A_{QP}
=\Sigma_P(z).
}
\]
\end{theorem}

\begin{proof}
Put $B=A_{QQ}$ and $T_z(t)=e^{-zt}e^{tB}=e^{t(B-zI)}$.
The generalized-eigenspace decomposition in finite dimensions writes $T_z$ as a
finite sum of constant matrix coefficients times $t^k e^{(\lambda-z)t}$.
This follows by expanding the nilpotent part on each generalized eigenspace;
equivalently, apply the decomposition to left multiplication by $B$ on the matrix
algebra and evaluate the exponential at $I$.
The strict spectral inequality makes every such term absolutely integrable on
$(0,\infty)$ and forces it to tend to zero. Hence $T_z$ is Bochner integrable
and $T_z(t)\to0$, including when $B$ is nonnormal.

Let $J=\int_0^\infty T_z(t)\,dt$. Integrating both derivative identities
$T_z'=(B-zI)T_z=T_z(B-zI)$ and using $T_z(0)=I$ gives
\[
(zI-B)J=J(zI-B)=I.
\]
Thus $J=(zI-B)^{-1}$. Continuous left and right multiplication carry the
Bochner integral through the two rectangular transfer blocks, yielding the claim.
\end{proof}

\begin{figure}[htbp]
\centering
\includegraphics[width=0.78\textwidth]{13_latent_memory_kernel.pdf}
\caption{The norm of the exact memory kernel produced by eliminated packets.}
\label{core:fig:memory}
\end{figure}

\section{Survival and hazard semigroups}\label{core:section:14}

Let $A$ be a finite subgenerator: its off-diagonal entries are nonnegative and its column sums are nonpositive. These sign conditions alone allow nonabsorbing states; eventual absorption is a separate hypothesis. Let $p_0$ be an initial transient-state distribution. Define
\[
p(t)=e^{tA}p_0,
\qquad
S(t)=\bm1^Tp(t).
\]
The exit-rate row vector is
\[
a^T=-\bm1^TA\ge0.
\]

\begin{theorem}[Semigroup survival identities]\label{core:record:043}
For $t\ge0$, the finite-time exit density and, wherever $S(t)>0$, the hazard are
\[
\boxed{
f(t)=a^Te^{tA}p_0=-S'(t),
\qquad
h(t)=\frac{f(t)}{S(t)}.
}
\]
Consequently, $S$ is nonincreasing and
\[
S(t)=\exp\left(-\int_0^th(r)\,dr\right)
\]
whenever $S(t)>0$. The density may be defective: its total mass is $1-\lim_{t\to\infty}S(t)$. It is a probability density when absorption is eventual, so that this limit is zero.
\end{theorem}

\begin{proof}
Differentiate:
\[
S'(t)=\bm1^TAe^{tA}p_0=-a^Te^{tA}p_0.
\]
The hazard identity is the definition of conditional instantaneous exit rate. Integrating $(\log S)'=-h$ gives the final formula.
\end{proof}

\begin{figure}[htbp]
\centering
\includegraphics[width=0.78\textwidth]{16_survival_hazard.pdf}
\caption{Survival and instantaneous hazard produced by one transient transfer semigroup.}
\label{core:fig:survival}
\end{figure}

\section{One common operator language}\label{core:section:15}

The preceding constructions are not identical, but they use the same packet operations. The domain dictionary is summarized in \cref{core:tab:dictionary}.

\begin{table}[htbp]
\centering
\small
\begin{tabularx}{\textwidth}{L{0.17\textwidth} L{0.19\textwidth} Y Y}
\toprule
Realization & Seed & Transfer & Packet/reconstruction law\\
\midrule
Markov & point mass $\delta_i$ & $K^n$ or $e^{tQ}$ & state lumping and $CK^n=\bar K^nC$\\
Conditional & random variable & $P_{\mathcal Q}$ & nested partitions and tower law\\
Bayesian & prior $p$ & $\cB_\lambda$ & packet evidences and normalized synthesis\\
Fractional & spectral source & $e^{-tL^\alpha}$ & spectral projectors $P_j$\\
PDE & source $f$ or $\delta_y$ & resolvent/Green operator & eigenmode or spatial packet synthesis\\
Feynman--Kac & terminal data $f$ & $(\diag(e^{-V})K^T)^n$ & weighted path reconstruction\\
Latent system & retained state & Feshbach/self-energy & hidden packets return as memory\\
Survival & initial transient law & $e^{tA}$ & sum over transient packets\\
\bottomrule
\end{tabularx}
\caption{The transfer-semigroup dictionary in the common operator language.}
\label{core:tab:dictionary}
\end{table}

For clarity, the finite-dimensional holomorphic calculus used below has the
following canonical definition. For $A$ acting on $\mathbb C^d$, its generalized
eigenspaces $E_\mu=\ker(A-\mu I)^d$ form a direct-sum decomposition over the
eigenvalues. Define the unique linear operator $F(A)$ by
\[
\left.F(A)\right|_{E_\mu}
=\left.\sum_{k=0}^{d-1}\frac{F^{(k)}(\mu)}{k!}(A-\mu I)^k\right|_{E_\mu}.
\]
The formula is independent of a basis and of increasing the truncation index:
the extra nilpotent powers vanish. It depends only on the finite Taylor jets of
$F$ at the spectrum. If $F$ is holomorphic on an open neighborhood $U$ of the
spectrum, each eigenvalue admits a radius $r_\mu>0$ with
$\overline{D(\mu,r_\mu)}\subset U$, and Cauchy's formula gives
\[
\frac{F^{(k)}(\mu)}{k!}
=\frac{1}{2\pi i}\oint_{|z-\mu|=r_\mu}
  \frac{F(z)}{(z-\mu)^{k+1}}\,dz.
\]
Consequently the same canonical operator is obtained from any such local
Cauchy circles, even when $U$ is disconnected. This is the standard finite
Taylor-jet definition of the holomorphic matrix function.

\begin{theorem}[Naturality under exact intertwining]\label{core:record:044}
For finite complex matrices, let $F$ be holomorphic on a neighborhood of the union of the spectra of $A_f$ and $A_c$. Real matrices are treated by complexification. If
\[
CA_f=A_cC,
\]
then
\[
\boxed{CF(A_f)=F(A_c)C.}
\]
In particular, exact generator refinement propagates simultaneously to semigroups, resolvents, integer powers, and analytic transfer functions.
\end{theorem}

\begin{proof}
The intertwining identity implies
$C(A_f-\mu I)^k=(A_c-\mu I)^kC$ for every $k\ge0$.
It therefore sends the generalized $\mu$-eigenspace of $A_f$ into that of $A_c$.
Use a common truncation index at least as large as both dimensions in the
canonical formula above. Multiplying its finite sum by $C$ now gives the desired
identity on every generalized eigenspace. The direct-sum decomposition gives
it on the full space. The argument covers nilpotent Jordan parts and uses only
local spectral Taylor jets; it does not require one Taylor series to converge
on both spectra. The zero-dimensional case is immediate.
\end{proof}

\begin{remark}[Where nonlinearity enters]\label{core:record:045}
Bayesian conditioning does not fit the preceding linear functional calculus because normalization depends on the current state. It nevertheless obeys an exact composition law and has a closed packet reconstruction. Paper III supplies the general nonlinear coordinate theorem; the Bayesian map is a particularly transparent probabilistic realization.
\end{remark}

\section{Numerical certificates}\label{core:section:16}

All calculations in this release were run on CPU. No GPU timing or performance claim is made. Floating-point residuals are implementation checks and are not substitutes for the symbolic proofs.

\begin{table}[htbp]
\centering
\small
\begin{tabular}{L{0.56\textwidth} r}
\toprule
Certificate & Value\\
\midrule
Packet-block semigroup composition residual & $2.554\times10^{-15}$\\
Operator reconstruction residual & $0$\\
Random-walk mass error & $0$\\
Exact Markov generator-lumping defect & $5.551\times10^{-17}$\\
Exact finite-horizon lumping error & $2.220\times10^{-16}$\\
Perturbed ten-step lumping error & $6.523\times10^{-3}$\\
Linear defect bound for that example & $2.000\times10^{-1}$\\
Bayesian composition residual & $5.551\times10^{-17}$\\
Bayesian packet reconstruction residual & $0$\\
Bayesian refinement commutation residual & $1.110\times10^{-16}$\\
Bayesian defect-factorization residual & $7.286\times10^{-17}$\\
Bayesian flow generator residual & $1.959\times10^{-11}$\\
Conditional-expectation tower residual & $0$\\
Hermite Gram residual & $9.992\times10^{-15}$\\
Fractional semigroup residual & $2.082\times10^{-17}$\\
Maximum square-root subordination error & $1.665\times10^{-16}$\\
Exact fractional-naturality residual & $0$\\
Fractional-power defect / gap certificate & $0.51059 / 0.95079$\\
Fractional-semigroup defect / certificate & $0.12969 / 0.44480$\\
Observed low-mode PDE convergence order & $1.999854$\\
Feynman--Kac enumeration residual & $1.665\times10^{-16}$\\
Feshbach projected-resolvent residual & $5.551\times10^{-17}$\\
Memory Laplace-transform residual & $6.939\times10^{-18}$\\
Duhamel refinement identity residual & $2.108\times10^{-16}$\\
Survival monotonicity violation & $0$\\
Hidden-Markov normalization error & $1.110\times10^{-16}$\\
\bottomrule
\end{tabular}
\caption{Selected results from the machine-readable numerical record.}
\label{core:tab:numerics}
\end{table}

\begin{certificate}[Test suite]\label{core:record:046}
The September 6 release records the following automated-test result:
\[
\boxed{54/54\text{ automated tests}.}
\]
The tests cover projector reconstruction, packet-block semigroup composition, exact and approximate Markov lumping, Bayesian composition and flow, the exact Bayesian refinement criterion, information-loss factorization, likelihood-profile partitions, conditional-expectation towers, Hermite orthogonality, fractional semigroups, square-root subordination, fractional intertwining and semigroup-defect bounds, discrete Green kernels, PDE convergence, Feynman--Kac enumeration, Schur--Feshbach elimination, memory kernels, refinement defects, hidden Markov filtering, and survival curves.
\end{certificate}

\section{Remaining conjectural extensions}\label{core:section:17}

The finite theory closes several operator identities but does not settle the corresponding universal continuum questions.

\begin{conjecture}[Transfer-semigroup realization]\label{core:record:047}
Every sufficiently regular realization of the seed calculus whose transfer stage is local in time admits a closed generator on a completed state space, and the packet analysis maps form a compatible system for its semigroup or evolution family.
\end{conjecture}

\begin{conjecture}[Uniform refinement control]\label{core:record:048}
For a nested family of physically or analytically meaningful discretizations, there exist norms and constants independent of resolution such that
\[
\norm{C_\ell e^{tA_{\ell+1}}-e^{tA_\ell}C_\ell}
\le
M\,t\,\norm{C_\ell A_{\ell+1}-A_\ell C_\ell}
\]
on every bounded time interval. Summable generator defects would then produce a coherent inverse-limit evolution.
\end{conjecture}

\begin{conjecture}[Latent-packet closure]\label{core:record:049}
For a stable infinite packet system, exact elimination of unresolved packets produces a causal memory kernel whose Laplace transform is the limiting packet self-energy. A Markovian closure exists precisely when this self-energy is constant, or can be represented by finitely many additional state packets.
\end{conjecture}

\begin{boundary}[What remains outside Paper VI]\label{core:record:050}
The following issues are not proved here:
\begin{enumerate}[leftmargin=2em]
\item existence and uniqueness for arbitrary stochastic differential equations;
\item posterior consistency or asymptotic normality in infinite statistical models;
\item domain invariance for unbounded fractional operators;
\item tightness and compactness required for stochastic continuum limits;
\item nonlinear PDE well-posedness beyond the stated finite and linear models;
\item an optimal packet choice for a general Markov, Bayesian, or PDE problem.
\end{enumerate}
These belong to later realization-specific work or to external analytic theory.
\end{boundary}

\section{Computational implementation}\label{core:section:18}

The accompanying Python package implements the formulas and finite models used in the paper.
\begin{center}
\texttt{projector\_transfer\_semigroups}
\end{center}
The computational organization mirrors the paper:
\[
\begin{array}{c|l}
\texttt{packets.py} & \text{projectors, blocks, reconstruction, block composition}\\
\texttt{semigroup.py} & \text{matrix semigroups, defects, Feshbach operators, memory}\\
\texttt{markov.py} & \text{random walks, exact lumping, survival}\\
\texttt{bayes.py} & \text{normalized updates, packet evidence, conjugacy, flows, HMM}\\
\texttt{conditional.py} & \text{conditional expectation and Gaussian chaos}\\
\texttt{fractional.py} & \text{spectral powers, subordination, and fractional refinement certificates}\\
\texttt{pde.py} & \text{Dirichlet Laplacians, Green kernels, eigenvalue error}\\
\texttt{feynman\_kac.py} & \text{weighted transfer and exhaustive path enumeration}.
\end{array}
\]

\subsection{Representative algorithms}

\begin{lstlisting}[caption={Normalized finite Bayesian update.}]
def bayes_update(prior, likelihood):
    p = normalize(prior)
    evidence = likelihood @ p
    if evidence <= 0:
        raise ValueError("positive evidence required")
    return likelihood * p / evidence
\end{lstlisting}

\begin{lstlisting}[caption={Packet block composition.}]
def block_composition(left, right):
    result = []
    for beta in range(len(left)):
        row = []
        for alpha in range(len(left)):
            value = sum(
                left[beta][gamma] @ right[gamma][alpha]
                for gamma in range(len(left))
            )
            row.append(value)
        result.append(row)
    return result
\end{lstlisting}

\begin{lstlisting}[caption={Finite Schur-Feshbach effective operator.}]
def feshbach_effective_operator(A, retained, z):
    App, Apq, Aqp, Aqq = split_blocks(A, retained)
    return z*I - App - Apq @ inv(z*I - Aqq) @ Aqp
\end{lstlisting}

The full implementation, tests, tables, figure-generation code, and exact numerical record are included with the release rather than embedded in the typeset paper.

\section{Conclusion}\label{core:section:19}

The transfer-semigroup realization gives one common operator language to several subjects that are often presented separately. Markov propagation is linear transfer of a probability seed. Conditional expectation is stochastic projection. Bayesian conditioning is normalized nonlinear transfer, and its exact finite coarsening law is now characterized by packetwise likelihood sufficiency. Fractional calculus is spectral packet reweighting, and bounded positive refinement defects now transfer to fractional powers and fractional semigroups by explicit resolvent and Duhamel estimates. A Green operator reconstructs a field from source packets. Feynman--Kac propagation adds path weights to stochastic transfer. Schur--Feshbach elimination shows that hidden packets are not destroyed; they return through self-energy and memory. Survival and hazard are scalar reconstructions of a transient semigroup.

The central structural identity is
\[
\boxed{
\text{generator defect}
\longrightarrow
\text{Duhamel semigroup defect}
\longrightarrow
\text{controlled refinement error}.
}
\]
Together with exact packet reconstruction, this makes the semigroup paper the natural bridge from the foundational calculus to the quantum, computational, and arithmetic realizations that follow.

\clearpage
\appendix
\section{Lean verification and scope}\label{core:section:20}\label{core:sec:lean-status}
\noindent\textbf{Fifth checked block, 6 September 2026.}
The release contains 341 explicit Lean theorem declarations across 47 proof modules,
including 31 additional declarations since the fourth checked block. The September 6 release records a successful full
build and namespace audit with Lean 4.34.0-rc2 and pinned mathlib. The audit
inspects 782 declarations, including definitions and generated auxiliaries; only
\texttt{propext}, \texttt{Classical.choice} and \texttt{Quot.sound} are allowed.
These are different counting units from the paper's 33 theorem-like results.

The ledger has \textbf{33 complete, zero partial and zero pending results}.
The last two results, 13.2 and 15.1, are now complete. The three separately stated
conjectures remain open and are excluded from this count.

``Complete'' means the stated scope follows from the named checked targets
and the ordinary specializations recorded in the ledger. Prose-to-Lean correspondence
is manually reviewed. The fractional targets use complex Hilbert spaces;
complexification for real operators is an ordinary mathematical specialization,
not a separately checked Lean construction. Completion of the theorem ledger does
not certify every prose sentence, numerical figure or open conjecture.

\small
\begin{longtable}{@{}p{0.08\textwidth}p{0.68\textwidth}p{0.17\textwidth}@{}}
\toprule Result & Manuscript statement & Lean scope\\
\midrule\endfirsthead
\toprule Result & Manuscript statement & Lean scope\\\midrule\endhead
\bottomrule\endfoot
1.2 & Finite semigroup law & Complete\\
2.1 & Packet semigroup convolution & Complete\\
2.2 & Reducing packets & Complete\\
3.1 & Exact Duhamel refinement identity & Complete\\
3.2 & Norm bound & Complete\\
3.3 & Resolvent refinement identity & Complete\\
4.1 & Finite Markov semigroup & Complete\\
5.1 & Exact lumpability & Complete\\
5.2 & Finite-horizon defect formula & Complete\\
6.1 & Finite conditional-expectation projector & Complete\\
6.2 & Tower law and martingale packets & Complete\\
7.1 & Independent-likelihood composition & Complete\\
7.2 & Bayesian packet reconstruction & Complete\\
7.3 & Likelihood perturbation bound & Complete\\
8.1 & Exact effective-likelihood factorization & Complete\\
8.2 & Bayesian refinement criterion & Complete\\
8.4 & Bayesian information-loss packet & Complete\\
8.5 & Nested Bayesian packet tower & Complete\\
8.6 & Canonical coarsest likelihood-resolving partition & Complete\\
9.1 & Bayesian nonlinear semigroup & Complete\\
9.3 & Prediction-update reconstruction & Complete\\
10.1 & Fractional packet calculus & Complete\\
10.2 & Poisson subordination & Complete\\
10.3 & Fractional intertwining theorem & Complete\\
10.4 & Fractional semigroup refinement & Complete\\
10.5 & Fractional refinement towers & Complete\\
11.1 & Exact discrete Green reconstruction at grid nodes & Complete\\
11.2 & Second-order low-mode convergence & Complete\\
12.1 & Discrete Feynman--Kac path formula & Complete\\
13.1 & Schur--Feshbach projected resolvent & Complete\\
13.2 & Memory/self-energy correspondence & Complete\\
14.1 & Semigroup survival identities & Complete\\
15.1 & Naturality under exact intertwining & Complete\\
\end{longtable}
\normalsize

\paragraph{General nonnormal memory (13.2).}
The actual damped exponential is decomposed into finitely many terms of the form
$t^k e^{(\mu-z)t}$. Integrability and convergence to zero follow from
$\operatorname{Re}\mu<\operatorname{Re}z$. The formal proof uses generalized
eigenspaces of left multiplication on the matrix algebra, so nilpotent Jordan
parts are included. Both inverse identities follow by integrating the two
exponential derivative formulas. A continuous linear sandwich map gives the
literal rectangular matrix memory integral and the actual inverse. No normality
or externally supplied growth bound is required.

\paragraph{Canonical holomorphic naturality (15.1).}
The standard finite Taylor-jet calculus defined before the theorem is constructed
using the direct sum of generalized eigenspaces. Uniqueness, cutoff stability and
exact rectangular naturality are proved. Cauchy's formula identifies the local
coefficients with actual circle integrals; admissible radii exist in every stated
open holomorphic neighborhood, and changing them leaves the canonical operator
unchanged. The argument permits disconnected neighborhoods and different matrix
dimensions. It makes no assertion about general Banach-algebra contour deformation.

\paragraph{Fractional results (10.3--10.5).}
The inherited fourth block proves the scalar beta normalization, the Lebesgue
integral representation of actual positive fractional powers, the literal
resolvent-defect formula, both printed norm constants and exact intertwining.
The rectangular complex-Hilbert-space semigroup proof substitutes those constants
into Duhamel's identity. The varying-space tower proof gives summability of
fractional-defect norms and uniform convergence of time-defect norm series on
bounded time intervals. Here ``summable'' means summable in norm. Operators with
unrelated codomains are not added together; a separate fixed-space theorem proves
operator-series convergence when such a sum is defined.

\paragraph{Verification and numerical evidence.}
The September 6 release records that the full pinned build, project-wide axiom audit, all 54 companion tests and six
verifier controls passed. No fresh Lean build is asserted for this October edition. The 54 tests comprise the original 37 and 17 added boundary
and validation cases. The 19 numerical figure pairs are unchanged. The earlier
79-field numerical comparison remains historical evidence, not a rerun in that September 6
block. The September 12 visual release subsequently records a numerical replay; current October checks are reported separately in the final package. Machine-readable source contracts, exact theorem targets, full logs and
previous checked ledgers accompany the source archive. Numerical tests are finite
regressions and are not substitutes for the Lean proofs.


\clearpage
\part*{Part II\quad Extended mathematical development}\addcontentsline{toc}{part}{Part II: Extended mathematical development}
\setcounter{section}{0}\setcounter{figure}{0}\setcounter{table}{0}
\renewcommand{\thesection}{E\arabic{section}}
\renewcommand{\theHsection}{E.\arabic{section}}
\renewcommand{\thefigure}{E\arabic{figure}}
\renewcommand{\theHfigure}{E.\arabic{figure}}
\renewcommand{\thetable}{E\arabic{table}}
\renewcommand{\theHtable}{E.\arabic{table}}
\graphicspath{{extension_figures/}}
\begin{abstract}
A bounded generator and a finite packet system turn a semigroup into a matrix of transfer blocks. We prove convolution, differential, path, Laplace, resolvent, and defect formulas. Exact intertwining yields an inverse limit semigroup. Weighted summable adjacent defects give convergent reconstructed evolutions, while range leakage determines whether the limit remains a semigroup. For finite Markov chains we identify a canonical lumping obstruction and quantify its failure. Conditional expectations give orthogonal martingale packets, and the Ornstein Uhlenbeck semigroup gives a countable Hermite model. Bayesian updating is treated as normalized transfer. We characterize exact coarse sufficiency, determine the coarsest nontrivial fixed sufficient partition for a likelihood family when the common level sets are nontrivial, and split relative information into coarse and local terms. Positive fractional powers are handled through spectral, Balakrishnan, and subordination formulas. Green operators, finite differences, and finite state Feynman Kac formulas provide analytic realizations. Schur Feshbach elimination produces a rational self energy and an exact memory kernel. Every stable finite hidden system yields a polynomial exponential expansion, including Jordan blocks. A certified optimal primary block selection and a Volterra estimate control reduced trajectories. State dependent killing, survival laws, moments, and competing risks complete the finite theory.
\end{abstract}


\section{Introduction}\label{extended:section:01}\label{extended:sec:introduction}

One parameter semigroups provide a common language for evolution equations, Markov processes, diffusion, and resolvent methods. Standard analytic foundations appear in \cite{Pazy1983,EngelNagel2000}. Finite Markov chains and their coarse descriptions are treated in \cite{KemenySnell1960,Norris1997,RogersPitman1981}. Conditional expectation and martingale difference decompositions are basic parts of probability theory \cite{Williams1991}. Gaussian transfer and Hermite expansions provide a spectral model for the Ornstein Uhlenbeck semigroup \cite{Janson1997}. Fractional powers of generators and subordinate semigroups belong to the functional calculus of positive and sectorial operators \cite{Balakrishnan1960,SchillingSongVondracek2012}. The Feynman Kac formula, projection methods, and phase type laws add three further views of transfer \cite{Kac1949,Feshbach1958,Zwanzig1961,Neuts1981}.

These subjects share a simple operation. A state is propagated, resolved into declared components, compared across levels, and reconstructed. The same operation can be linear, as in a Markov semigroup, or normalized and nonlinear, as in Bayesian conditioning. Hidden components can be discarded from the state vector without disappearing from the dynamics. Their exact influence returns through a frequency dependent self energy or a time dependent memory kernel.

The first part of the article treats bounded semigroups on Banach spaces. A finite packet system resolves every operator into input and output blocks. The semigroup law becomes block convolution. Splitting the generator into its packet diagonal and transfer part gives a convergent Dyson expansion over directed packet paths. For a coarse map \(C\), the generator defect
\[
\Delta=CA_f-A_cC
\]
produces the entire semigroup defect through a Duhamel integral. The same defect controls resolvents and holomorphic functions of the generators.

The next question concerns refinement towers. Exact intertwining gives a coordinatewise semigroup on the bounded inverse limit. Approximate refinement is more delicate. Even if every adjacent defect is small, reconstruction from a fine level can introduce a second error because a coarse state need not remain inside the reconstructed fine subspace. We separate these two effects. Summability of adjacent defects gives a uniform limit of reconstructed coarse evolutions. An asymptotic closure condition is then sufficient for the limiting family to satisfy the semigroup law.

The probabilistic part begins with finite Markov chains. Exact lumpability is the vanishing of a canonical obstruction on the kernel of aggregation. In coordinates, this is equality of the transition mass sent from any two states in the same source block into every target block. The usual finite horizon defect formula follows from telescoping. Conditional expectation is then treated as an orthogonal projector. A nested family of partitions produces mutually orthogonal martingale packets and an exact energy decomposition.

Bayesian updating requires more care because normalization depends on the current prior. Given a partition, every fine update has an exact two stage description. The block masses are updated by prior dependent average likelihoods, while the conditional distribution inside each block is updated by its centered likelihood variation. A fixed coarse Bayesian rule works for every prior precisely when the likelihood is constant inside each block. This gives a finite and exact form of Bayesian sufficiency. For a family of likelihoods with at least two common level sets, those level sets form the unique coarsest nontrivial fixed sufficient partition. Relative entropy then separates exactly into information retained by the coarse packet masses and information carried by the local conditional changes. The same part derives perturbation bounds, the replicator flow, a variance identity for accumulated evidence, and stability estimates for hidden Markov filtering.

For a positive operator \(L\), spectral packets define \(L^\alpha\) and \(\e^{-tL^\alpha}\). Exact intertwining passes to every fractional power. When both fine and coarse operators have a common positive lower bound \(m\), the Balakrishnan formula gives the quantitative estimate
\[
\norm{CL_f^\alpha-L_c^\alpha C}
\leq
\alpha m^{\alpha-1}
\norm{CL_f-L_cC},
\qquad 0<\alpha<1.
\]
A Duhamel argument then controls the fractional semigroup defect. This closes a finite version of fractional naturality under explicit spectral hypotheses.

The final sections move from generators to source reconstruction and hidden variables. The one dimensional Dirichlet Green kernel is reproduced exactly by the scaled inverse of the second difference matrix at grid nodes. A discrete maximum principle gives a complete second order error estimate for smooth Poisson solutions. Finite state Feynman Kac transfer is proved in discrete and continuous time. Schur Feshbach elimination yields the exact projected resolvent, the retained memory equation, and the Laplace correspondence between memory and self energy. Every stable finite hidden system gives a polynomial exponential memory expansion. Primary blocks can be ranked optimally for the certified \(L^1\) bound, and a Volterra estimate carries kernel error to the retained solution. Survival, event density, moments, and competing risks are then read directly from a transient semigroup.

Several claims are intentionally absent. No arbitrary unbounded operator is asserted to generate a semigroup. The inverse limit theorem uses bounded generators or an explicit uniform continuity condition. Bayesian sufficiency is proved for a finite partition and a fixed likelihood. The fractional defect estimate uses a common spectral gap. The Poisson estimate concerns the stated one dimensional scheme. Memory compression is controlled for the declared finite or integrable kernels. General stochastic differential equations, posterior consistency, nonlinear partial differential equations, and optimal model reduction remain outside the article.

The same words have different meanings in the realizations below. The table records the object that is transferred and the condition that makes its coarse description exact.

\begin{table}[t]
\centering
\small
\begin{tabular}{>{\raggedright\arraybackslash}p{.18\textwidth}>{\raggedright\arraybackslash}p{.29\textwidth}>{\raggedright\arraybackslash}p{.43\textwidth}}
\toprule
Realization & Transfer & Exact coarse condition \\
\midrule
Operator evolution & $\e^{tA}$ & generator intertwining \\
Markov chain & $K^n$ or $\e^{tQ}$ & equality of aggregated transition masses \\
Conditional expectation & $P_{\mathcal Q}$ & nesting of the underlying partitions \\
Bayesian update & $\Bcal_\lambda$ & likelihood constant inside each coarse packet \\
Fractional diffusion & $\e^{-tL^\alpha}$ & spectral or resolvent intertwining \\
Hidden dynamics & projected resolvent & vanishing self energy, or an explicit memory state \\
Survival law & transient semigroup & summation over all transient states \\
\bottomrule
\end{tabular}
\caption{The transfer operations studied in the article and the condition that makes each coarse description exact.}
\label{extended:tab:transfer_dictionary}
\end{table}

\section{Packet resolved semigroups}\label{extended:section:02}\label{extended:sec:packet_semigroups}

Let \(X\) be a complex Banach space. A finite packet system is a family of bounded idempotents
\[
\Pcal=(P_\alpha)_{\alpha\in I}
\]
such that
\[
P_\alpha P_\beta=\delta_{\alpha\beta}P_\alpha,
\qquad
\sum_{\alpha\in I}P_\alpha=I_X.
\]
No orthogonality is needed in this section.

\begin{definition}[Packet resolved transfer system]\label{extended:record:001}
A packet resolved transfer system is a triple
\[
(X,A,\Pcal)
\]
consisting of a Banach state space, a bounded generator, and a finite packet system. A marked transfer system also includes an initial state or source $s\in X$. The topology and all limiting operations are supplied by the ambient Banach space structure.
\end{definition}

\begin{definition}[Packet blocks]\label{extended:record:002}
For \(B\in\mathcal B(X)\), define
\[
B_{\beta\alpha}=P_\beta BP_\alpha.
\]
If \(A\in\mathcal B(X)\) and \(T(t)=\e^{tA}\), write
\[
T_{\beta\alpha}(t)=P_\beta T(t)P_\alpha.
\]
\end{definition}

\begin{theorem}[Semigroup and packet convolution]\label{extended:record:003}\label{extended:thm:packet_convolution}
For every bounded operator \(A\), the family \(T(t)=\e^{tA}\) is a uniformly continuous group and
\[
T(s+t)=T(s)T(t).
\]
For every \(\alpha,\beta\in I\),
\[
T_{\beta\alpha}(s+t)
=
\sum_{\gamma\in I}
T_{\beta\gamma}(s)T_{\gamma\alpha}(t).
\]
\end{theorem}

\begin{proof}
The exponential series converges in operator norm. Since every power of \(A\) commutes with every other power, the Cauchy product gives the group law. Inserting \(I_X=\sum_\gamma P_\gamma\) between the two factors gives the block formula.
\end{proof}

\begin{proposition}[Block differential system]\label{extended:record:004}\label{extended:prop:block_differential}
The packet blocks are differentiable in operator norm and satisfy
\[
\frac{\dd}{\dd t}T_{\beta\alpha}(t)
=
\sum_{\gamma\in I}
A_{\beta\gamma}T_{\gamma\alpha}(t)
=
\sum_{\gamma\in I}
T_{\beta\gamma}(t)A_{\gamma\alpha}.
\]
Their initial values are
\[
T_{\beta\alpha}(0)=\delta_{\beta\alpha}P_\alpha.
\]
\end{proposition}

\begin{proof}
Differentiate \(T(t)\) in operator norm and insert the packet resolution between \(A\) and \(T(t)\), then between \(T(t)\) and \(A\).
\end{proof}

\begin{theorem}[Laplace resolvent and packet transform]\label{extended:record:005}\label{extended:thm:laplace_packet_resolvent}
Let \(A\in\mathcal B(X)\), let \(T(t)=\e^{tA}\), and let
\[
\operatorname{Re}z>\norm A.
\]
Then the Bochner integral converges in operator norm and
\[
(zI-A)^{-1}
=
\int_0^\infty \e^{-zt}T(t)\dd t.
\]
For every pair of packets,
\[
P_\beta(zI-A)^{-1}P_\alpha
=
\int_0^\infty
\e^{-zt}T_{\beta\alpha}(t)\dd t.
\]
\end{theorem}

\begin{proof}
The estimate \(\norm{T(t)}\leq\e^{t\norm A}\) gives absolute convergence of the operator valued integral. Let
\[
R(z)=\int_0^\infty\e^{-zt}T(t)\dd t.
\]
Since \(T'(t)=AT(t)=T(t)A\), integration by parts gives
\[
(zI-A)R(z)
=
-\int_0^\infty
\frac{\dd}{\dd t}
\bigl(\e^{-zt}T(t)\bigr)
\dd t
=I.
\]
The same calculation on the other side gives
\(
R(z)(zI-A)=I
\).
Multiplication by \(P_\beta\) and \(P_\alpha\) may be passed under the norm convergent integral, which gives the packet formula.
\end{proof}

\begin{corollary}[Resolvent expansion over packet paths]\label{extended:record:006}\label{extended:cor:resolvent_packet_paths}
With \(A=D+N\) as in the directed path decomposition, let
\[
R_D(z)=(zI-D)^{-1}.
\]
If
\[
\norm{R_D(z)N}<1,
\]
then
\[
(zI-A)^{-1}
=
\sum_{m=0}^\infty
\bigl(R_D(z)N\bigr)^mR_D(z).
\]
After multiplication by \(P_\beta\) and \(P_\alpha\), the term of order \(m\) is a sum over directed packet paths of length \(m\) from \(\alpha\) to \(\beta\).
\end{corollary}

\begin{proof}
Factor
\[
zI-A=(zI-D)(I-R_D(z)N)
\]
and expand the second inverse by its norm convergent Neumann series. The diagonal resolvent preserves each packet and every factor \(N\) changes the packet along an active transfer edge.
\end{proof}

\begin{theorem}[Reducing packet criterion]\label{extended:record:007}\label{extended:thm:reducing_packets}
The following conditions are equivalent.
\begin{enumerate}
\item \(AP_\alpha=P_\alpha A\) for every \(\alpha\).
\item \(A_{\beta\alpha}=0\) whenever \(\beta\neq\alpha\).
\item \(T_{\beta\alpha}(t)=0\) whenever \(\beta\neq\alpha\) and \(t\in\RR\).
\item Every packet space \(P_\alpha X\) is invariant under \(T(t)\) for every \(t\).
\end{enumerate}
When these conditions hold,
\[
T(t)|_{P_\alpha X}
=
\e^{tA_{\alpha\alpha}}.
\]
\end{theorem}

\begin{proof}
The first two conditions are equivalent because the packet resolution gives
\[
AP_\alpha=\sum_\beta A_{\beta\alpha},
\qquad
P_\alpha A=\sum_\gamma A_{\alpha\gamma}.
\]
A block diagonal generator has block diagonal powers, so the second condition implies the third. Differentiating the off diagonal blocks at zero proves the converse. The equivalence with invariance follows from the definition of the blocks. The final identity follows by restricting the exponential series.
\end{proof}

\begin{figure}[t]
\centering
\includegraphics[width=.31\textwidth]{packet_blocks_t0_0.pdf}\hfill
\includegraphics[width=.31\textwidth]{packet_blocks_t0_8.pdf}\hfill
\includegraphics[width=.31\textwidth]{packet_blocks_t2_0.pdf}
\caption{Packet block norms at three times. Transfer appears in the off diagonal blocks as soon as the generator couples the two packet spaces.}
\label{extended:fig:packet_blocks}
\end{figure}

\subsection{Directed path expansion}

Set
\[
D=\sum_{\alpha\in I}P_\alpha AP_\alpha,
\qquad
N=A-D.
\]
The operator \(D\) preserves every packet and \(N\) contains only transfers between distinct packets.

\begin{theorem}[Dyson expansion over packet paths]\label{extended:record:008}\label{extended:thm:dyson_paths}
For \(t\geq0\),
\[
\begin{aligned}
\e^{tA}
={}&\e^{tD}
+
\sum_{m=1}^{\infty}
\int_{0<s_1<\cdots<s_m<t}
\e^{(t-s_m)D}N\e^{(s_m-s_{m-1})D}
\cdots
N\e^{s_1D}
\dd s_1\cdots\dd s_m.
\end{aligned}
\]
The series converges absolutely in operator norm. After multiplication by \(P_\beta\) and \(P_\alpha\), the term of order \(m\) is a sum over directed packet paths of length \(m\) from \(\alpha\) to \(\beta\).
\end{theorem}

\begin{proof}
Write the variation of constants identity
\[
\e^{tA}
=
\e^{tD}
+
\int_0^t\e^{(t-s)D}N\e^{sA}\dd s.
\]
Iterating the identity gives the displayed series. The norm of the term of order \(m\) is at most
\[
\e^{t\norm D}\frac{(t\norm N)^m}{m!}.
\]
The majorant is summable. Inserting the packet resolution around every factor \(N\) leaves only products whose consecutive packet labels are joined by nonzero transfer blocks.
\end{proof}

\begin{corollary}[Termination on an acyclic transfer graph]\label{extended:record:009}\label{extended:cor:acyclic_termination}
Suppose the directed graph with an edge \(\alpha\to\beta\) whenever \(P_\beta NP_\alpha\neq0\) is acyclic. If its longest directed path has length \(q\), then every Dyson term of order greater than \(q\) vanishes.
\end{corollary}

\begin{proof}
Every nonzero term of order \(m\) contains a directed transfer path of length \(m\). An acyclic graph has no path longer than \(q\).
\end{proof}

\section{Generator defects and functional calculus}\label{extended:section:03}\label{extended:sec:defects}

Let \(X_f\) and \(X_c\) be complex Banach spaces. Let \(A_f\in\mathcal B(X_f)\), \(A_c\in\mathcal B(X_c)\), and \(C\in\mathcal B(X_f,X_c)\). Define
\[
\Delta=CA_f-A_cC
\]
and
\[
\mathfrak D(t)=C\e^{tA_f}-\e^{tA_c}C.
\]

\begin{theorem}[Duhamel defect identity]\label{extended:record:010}\label{extended:thm:duhamel_defect}
For every \(t\geq0\),
\[
\mathfrak D(t)
=
\int_0^t
\e^{(t-s)A_c}\Delta\e^{sA_f}
\dd s.
\]
If
\[
\norm{\e^{tA_f}}\leq M_f\e^{\omega_f t},
\qquad
\norm{\e^{tA_c}}\leq M_c\e^{\omega_c t},
\]
then
\[
\norm{\mathfrak D(t)}
\leq
M_fM_c\norm\Delta
\int_0^t
\e^{\omega_c(t-s)+\omega_fs}
\dd s.
\]
For contraction semigroups,
\[
\norm{\mathfrak D(t)}\leq t\norm\Delta.
\]
\end{theorem}

\begin{proof}
For fixed \(t\), set
\[
F(s)=\e^{(t-s)A_c}C\e^{sA_f}.
\]
Then
\[
F'(s)=\e^{(t-s)A_c}\Delta\e^{sA_f}.
\]
Integrating from zero to \(t\) gives the identity. The estimates follow by taking norms under the integral.
\end{proof}

\begin{proposition}[Infinitesimal recovery of the defect]\label{extended:record:011}\label{extended:prop:infinitesimal_defect}
The generator defect is the first order semigroup defect:
\[
\lim_{t\to0}
\frac{\mathfrak D(t)}{t}
=
\Delta
\]
in operator norm. Consequently,
\[
C\e^{tA_f}=\e^{tA_c}C
\quad\hbox{for every }t
\]
if and only if \(CA_f=A_cC\).
\end{proposition}

\begin{proof}
The operator norm expansions
\[
\e^{tA_f}=I+tA_f+o(t),
\qquad
\e^{tA_c}=I+tA_c+o(t)
\]
give the limit. The equivalence follows from this limit and \cref{extended:thm:duhamel_defect}.
\end{proof}

\begin{theorem}[Resolvent defect identity]\label{extended:record:012}\label{extended:thm:resolvent_defect}
Let \(z\) belong to the resolvent sets of \(A_f\) and \(A_c\). With
\[
R_f(z)=(zI-A_f)^{-1},
\qquad
R_c(z)=(zI-A_c)^{-1},
\]
one has
\[
CR_f(z)-R_c(z)C
=
R_c(z)\Delta R_f(z).
\]
\end{theorem}

\begin{proof}
Multiply the proposed identity on the left by \(zI-A_c\) and on the right by \(zI-A_f\). Both sides reduce to
\[
(zI-A_c)C-C(zI-A_f)=\Delta.
\]
\end{proof}

\begin{theorem}[Holomorphic defect formula]\label{extended:record:013}\label{extended:thm:holomorphic_defect}
Let \(\Gamma\) be the positively oriented rectifiable boundary of a finite union of pairwise disjoint bounded Jordan domains. Assume that both spectra lie inside this union and that \(\Gamma\) lies in the common resolvent set. Thus its winding number is one about every spectral point. If \(g\) is holomorphic on a neighborhood of the closure of the union, then
\[
Cg(A_f)-g(A_c)C
=
\frac{1}{2\pi\ii}
\int_\Gamma
 g(z)R_c(z)\Delta R_f(z)
\dd z.
\]
In particular,
\[
\begin{aligned}
\norm{Cg(A_f)-g(A_c)C}
\leq{}&
\frac{\operatorname{length}(\Gamma)}{2\pi}
\sup_{z\in\Gamma}\abs{g(z)}
\sup_{z\in\Gamma}\norm{R_c(z)}
\sup_{z\in\Gamma}\norm{R_f(z)}
\norm\Delta.
\end{aligned}
\]
\end{theorem}

\begin{proof}
Apply the Dunford integral formula to both operators and use \cref{extended:thm:resolvent_defect} under the contour integral. The norm estimate is immediate.
\end{proof}

\begin{corollary}[Exact naturality]\label{extended:record:014}\label{extended:cor:exact_naturality}
If \(CA_f=A_cC\), then
\[
Cg(A_f)=g(A_c)C
\]
for every holomorphic function covered by \cref{extended:thm:holomorphic_defect}. The conclusion includes exponentials, resolvents, and polynomials.
\end{corollary}

\begin{proof}
Set \(\Delta=0\) in \cref{extended:thm:holomorphic_defect}.
\end{proof}

\begin{figure}[t]
\centering
\includegraphics[width=.68\textwidth]{duhamel_defect.pdf}
\caption{A finite semigroup defect and a direct Duhamel bound. The bound is designed to be uniform rather than sharp in this example.}
\label{extended:fig:duhamel_defect}
\end{figure}

\section{Refinement towers and limiting evolution}\label{extended:section:04}\label{extended:sec:refinement_towers}

The exact inverse limit and the approximate completion problem require different hypotheses. Exact intertwining preserves the inverse limit coordinate by coordinate. Approximate intertwining gives a Cauchy family only after reconstruction maps are chosen.

\subsection{Exact inverse limits}

Let
\[
X_0\xleftarrow{C_0}X_1\xleftarrow{C_1}X_2\xleftarrow{}\cdots
\]
be Banach spaces with bounded connecting maps. Define the bounded inverse limit
\[
X_\infty
=
\left\{
(x_\ell)_{\ell\geq0}:
C_\ell x_{\ell+1}=x_\ell,
\quad
\sup_\ell\norm{x_\ell}<\infty
\right\}
\]
with the supremum norm.

\begin{lemma}[Completeness of the bounded inverse limit]\label{extended:record:015}\label{extended:lem:inverse_limit_banach}
The space \(X_\infty\) is Banach.
\end{lemma}

\begin{proof}
The product of the spaces \(X_\ell\), equipped with the bounded supremum norm, is Banach. The equations \(C_\ell x_{\ell+1}=x_\ell\) define a closed linear subspace.
\end{proof}

\begin{theorem}[Exact projective limit semigroup]\label{extended:record:016}\label{extended:thm:exact_projective_limit}
For each \(\ell\), let \(T_\ell(t)\) be a strongly continuous semigroup on \(X_\ell\). Assume
\[
C_\ell T_{\ell+1}(t)=T_\ell(t)C_\ell
\]
for every \(\ell\) and \(t\geq0\), and assume the uniform bound
\[
\sup_\ell\norm{T_\ell(t)}\leq M\e^{\omega t}.
\]
Define
\[
T_\infty(t)(x_\ell)_\ell
=
(T_\ell(t)x_\ell)_\ell.
\]
Then \(T_\infty(t)\) is a bounded semigroup on \(X_\infty\). It is strongly continuous if
\[
\lim_{t\downarrow0}
\sup_\ell
\norm{T_\ell(t)x_\ell-x_\ell}=0
\]
for every \((x_\ell)_\ell\in X_\infty\).

If \(T_\ell(t)=\e^{tA_\ell}\), the generators are bounded, and
\[
\sup_\ell\norm{A_\ell}<\infty,
\qquad
C_\ell A_{\ell+1}=A_\ell C_\ell,
\]
then the continuity condition holds and
\[
A_\infty(x_\ell)_\ell=(A_\ell x_\ell)_\ell
\]
is the bounded generator of \(T_\infty\).
\end{theorem}

\begin{proof}
Exact intertwining shows that \((T_\ell(t)x_\ell)_\ell\) remains compatible. The uniform growth bound makes the coordinatewise operator bounded on the supremum norm. The semigroup law holds in every coordinate. The stated continuity condition is exactly strong continuity in the supremum norm.

For bounded generators with common bound \(M_A\),
\[
\sup_\ell
\norm{(\e^{tA_\ell}-I)x_\ell}
\leq
(\e^{tM_A}-1)\sup_\ell\norm{x_\ell}.
\]
This tends to zero with \(t\). Exact generator intertwining makes \((A_\ell x_\ell)_\ell\) compatible, and coordinatewise differentiation identifies the generator.
\end{proof}

\begin{remark}\label{extended:record:017}
Strong continuity on each level does not by itself imply strong continuity in the supremum norm of the inverse limit. The additional uniform condition in \cref{extended:thm:exact_projective_limit} cannot be omitted without replacing the topology.
\end{remark}

\subsection{Summable adjacent defects}

Assume now that each \(C_\ell:X_{\ell+1}\to X_\ell\) has a bounded right inverse
\[
R_\ell:X_\ell\to X_{\ell+1},
\qquad
C_\ell R_\ell=I_{X_\ell}.
\]
For \(m<n\), set
\[
C_{m,n}=C_mC_{m+1}\cdots C_{n-1},
\qquad
R_{n,m}=R_{n-1}\cdots R_m.
\]
The empty products are
\[
C_{m,m}=I_{X_m},
\qquad
R_{m,m}=I_{X_m}.
\]

\begin{definition}[Reconstructed fine evolution]\label{extended:record:018}
For a semigroup \(T_n(t)\) on \(X_n\), define the evolution seen from level \(m\) by
\[
U_{m,n}(t)=C_{m,n}T_n(t)R_{n,m}.
\]
\end{definition}

\begin{theorem}[Weighted summable defect completion]\label{extended:record:019}\label{extended:thm:summable_defect_completion}
Define the adjacent semigroup defect
\[
E_n(t)=C_nT_{n+1}(t)-T_n(t)C_n.
\]
Fix a coarse level \(m\) and a time horizon \(T>0\). For \(n\geq m\), set
\[
w_{m,n}
=
\norm{C_{m,n}}\,
\norm{R_nR_{n,m}}.
\]
If
\[
\sum_{n=m}^{\infty}
w_{m,n}
\sup_{0\leq t\leq T}\norm{E_n(t)}<\infty,
\]
then \(U_{m,n}(t)\) converges in operator norm uniformly for \(0\leq t\leq T\). Its limit is denoted by \(U_m(t)\), and
\[
\sup_{0\leq t\leq T}
\norm{U_m(t)-U_{m,n}(t)}
\leq
\sum_{j=n}^{\infty}w_{m,j}
\sup_{0\leq t\leq T}\norm{E_j(t)}.
\]
\end{theorem}

\begin{proof}
Since \(C_nR_n=I\),
\[
\begin{aligned}
U_{m,n+1}(t)-U_{m,n}(t)
={}&
C_{m,n}
\bigl(C_nT_{n+1}(t)-T_n(t)C_n\bigr)
R_nR_{n,m}.
\end{aligned}
\]
Therefore
\[
\norm{U_{m,n+1}(t)-U_{m,n}(t)}
\leq
w_{m,n}\norm{E_n(t)}.
\]
The weighted summability hypothesis makes the sequence uniformly Cauchy. Summing the tail estimates gives the stated bound.
\end{proof}

\begin{corollary}[Contractive summable defect completion]\label{extended:record:020}\label{extended:cor:contractive_defect_completion}
Assume
\[
\norm{C_\ell}\leq1,
\qquad
\norm{R_\ell}\leq1,
\qquad
\norm{T_\ell(t)}\leq1
\]
for every \(\ell\) and \(t\geq0\). If
\[
\sum_{\ell=0}^{\infty}
\sup_{0\leq t\leq T}\norm{E_\ell(t)}<\infty,
\]
then the conclusion of \cref{extended:thm:summable_defect_completion} holds, and
\[
\sup_{0\leq t\leq T}
\norm{U_m(t)-U_{m,n}(t)}
\leq
\sum_{\ell=n}^{\infty}
\sup_{0\leq t\leq T}\norm{E_\ell(t)}.
\]
\end{corollary}

\begin{proof}
Every weight \(w_{m,n}\) is at most one.
\end{proof}

\begin{corollary}[Generator defect criterion]\label{extended:record:021}\label{extended:cor:generator_defect_completion}
Assume
\[
\norm{C_\ell}\leq1,
\qquad
\norm{R_\ell}\leq1,
\]
and suppose \(T_\ell(t)=\e^{tA_\ell}\) are contraction semigroups with bounded generators. Set
\[
\Delta_\ell=C_\ell A_{\ell+1}-A_\ell C_\ell.
\]
If
\[
\sum_{\ell=0}^{\infty}\norm{\Delta_\ell}<\infty,
\]
then the conclusion of \cref{extended:cor:contractive_defect_completion} holds on every bounded time interval, and
\[
\sup_{0\leq t\leq T}
\norm{U_m(t)-U_{m,n}(t)}
\leq
T\sum_{\ell=n}^{\infty}\norm{\Delta_\ell}.
\]
\end{corollary}

\begin{proof}
Apply the contraction estimate in \cref{extended:thm:duhamel_defect} to every adjacent pair.
\end{proof}

The limit in \cref{extended:thm:summable_defect_completion} need not satisfy the semigroup law. The obstruction comes from leaving the reconstructed fine subspace between two time steps.

\begin{theorem}[Asymptotic closure criterion]\label{extended:record:022}\label{extended:thm:asymptotic_closure}
Assume the weighted convergence in \cref{extended:thm:summable_defect_completion} holds for every bounded time interval, and define
\[
P_{m,n}=R_{n,m}C_{m,n}
\]
on \(X_n\). Suppose that for every \(T>0\),
\[
\rho_{m,n}(T)
=
\sup_{0\leq s,t\leq T}
\norm{
C_{m,n}T_n(s)(I-P_{m,n})T_n(t)R_{n,m}
}
\longrightarrow0.
\]
Then
\[
U_m(s+t)=U_m(s)U_m(t)
\]
for all \(s,t\geq0\). If the convergence in \cref{extended:thm:summable_defect_completion} is locally uniform and every \(T_n\) is norm continuous, then \(U_m\) is a norm continuous semigroup.
\end{theorem}

\begin{proof}
At a finite level,
\[
\begin{aligned}
U_{m,n}(s)U_{m,n}(t)-U_{m,n}(s+t)
={}&
C_{m,n}T_n(s)
(P_{m,n}-I)
T_n(t)R_{n,m}.
\end{aligned}
\]
Its norm is at most \(\rho_{m,n}(T)\) for \(s,t\leq T\). Pass to the limit using uniform convergence. Norm continuity follows from a uniform limit of norm continuous functions on compact intervals.
\end{proof}

\begin{remark}\label{extended:record:023}
Summable generator defects control consistency between adjacent levels. The quantity \(\rho_{m,n}(T)\) controls closure of the reconstructed fine range under the fine evolution. These are different requirements.
\end{remark}

\begin{example}[Summable defects do not force a semigroup limit]\label{extended:record:024}\label{extended:ex:summable_not_semigroup}
Let \(X_0=\RR\) and let \(X_n=\RR^2\) for every \(n\geq1\). Set
\[
C_0(x,y)=x,
\qquad
R_0x=(x,0),
\]
and use identity connecting maps and reconstructions above level one. Let \(T_0(t)=I\) and let every higher level carry the rotation group
\[
T_n(t)
=
\begin{pmatrix}
\cos t&-\sin t\\
\sin t&\cos t
\end{pmatrix}.
\]
Only the first adjacent defect is nonzero, so every defect sum is finite. Yet for \(n\geq1\),
\[
U_{0,n}(t)=C_0T_n(t)R_0=\cos t.
\]
The limiting family is not a semigroup because in general
\[
\cos(s+t)\neq\cos s\cos t.
\]
The failure is exactly the excursion from the reconstructed line into its orthogonal complement between two time steps.
\end{example}

\begin{proposition}[Range leakage criterion]\label{extended:record:025}\label{extended:prop:range_leakage}
Under the contractive assumptions of \cref{extended:cor:contractive_defect_completion}, define
\[
\eta_{m,n}(T)
=
\sup_{0\leq t\leq T}
\norm{(I-P_{m,n})T_n(t)R_{n,m}}.
\]
Then
\[
\rho_{m,n}(T)\leq\eta_{m,n}(T).
\]
Consequently, \(\eta_{m,n}(T)\to0\) for every \(T>0\) is sufficient for the limiting family \(U_m\) to satisfy the semigroup law.
\end{proposition}

\begin{proof}
For \(s,t\leq T\), the contraction assumptions give
\[
\begin{aligned}
&\norm{
C_{m,n}T_n(s)(I-P_{m,n})T_n(t)R_{n,m}
}\\
&\qquad\leq
\norm{(I-P_{m,n})T_n(t)R_{n,m}}.
\end{aligned}
\]
Take the supremum and apply \cref{extended:thm:asymptotic_closure}.
\end{proof}

\begin{theorem}[Generator control of range leakage]\label{extended:record:026}\label{extended:thm:generator_range_leakage}
Let \(T(t)=\e^{tA}\) on a Banach space \(X\), let \(P\) be a bounded projection, and set \(Q=I-P\). Assume
\[
\norm{T(t)}\leq M\e^{\omega t}
\]
and
\[
\norm{\e^{t(QAQ)|_{QX}}}
\leq M_Q\e^{\omega_Qt}.
\]
Then
\[
\norm{QT(t)P}
\leq
M M_Q\norm P^2\norm{QAP}
\int_0^t
\e^{\omega_Q(t-s)+\omega s}
\dd s.
\]
In particular, for orthogonal projections and contraction bounds,
\[
\sup_{0\leq t\leq T}
\norm{QT(t)P}
\leq
T\norm{QAP}.
\]
\end{theorem}

\begin{proof}
For \(x\in PX\), write
\[
p(t)=PT(t)x,
\qquad
q(t)=QT(t)x.
\]
The hidden component satisfies
\[
q'(t)=QAQq(t)+QAPp(t),
\qquad
q(0)=0.
\]
Variation of constants gives
\[
q(t)=
\int_0^t
\e^{(t-s)(QAQ)|_{QX}}
QAPp(s)
\dd s.
\]
Since
\(
\norm{p(s)}\leq\norm P\norm{T(s)}\norm x
\)
and \(x=Px\), the first estimate follows. The contraction case is immediate.
\end{proof}

\begin{corollary}[A usable closure condition for refinement towers]\label{extended:record:027}\label{extended:cor:generator_closure_tower}
Assume the contractive summability hypothesis of \cref{extended:cor:contractive_defect_completion} holds on every bounded time interval. Suppose every \(P_{m,n}\) is an orthogonal projection on a Hilbert space and \(T_n(t)=\e^{tA_n}\) is contractive. If, for every fixed \(m\),
\[
\norm{(I-P_{m,n})A_nP_{m,n}}
\longrightarrow0,
\]
then the limiting family \(U_m\) is a semigroup.
\end{corollary}

\begin{proof}
A bounded generator of a contraction semigroup on a Hilbert space is dissipative. If \(Q_{m,n}=I-P_{m,n}\) and \(y\in Q_{m,n}X_n\), then
\[
\operatorname{Re}\ip{Q_{m,n}A_nQ_{m,n}y}{y}
=
\operatorname{Re}\ip{A_ny}{y}
\leq0.
\]
Hence the compressed hidden semigroup is contractive. Apply \cref{extended:thm:generator_range_leakage} to \(P_{m,n}\). Since \(P_{m,n}R_{n,m}=R_{n,m}\),
\[
\eta_{m,n}(T)
\leq
T\norm{(I-P_{m,n})A_nP_{m,n}}
\norm{R_{n,m}}.
\]
The right side tends to zero. Use \cref{extended:prop:range_leakage}.
\end{proof}

\section{Finite Markov transfer and lumping}\label{extended:section:05}\label{extended:sec:markov}

Probability vectors are columns. A column stochastic matrix \(K\) has nonnegative entries and column sums one. A continuous time generator \(Q\) has nonnegative off diagonal entries and column sums zero.

\begin{theorem}[Finite Markov semigroup]\label{extended:record:028}\label{extended:thm:finite_markov}
If \(Q\) is a finite Markov generator, then \(\e^{tQ}\) is column stochastic for every \(t\geq0\).
\end{theorem}

\begin{proof}
Choose \(a\geq\max_i(-Q_{ii})\). The matrix \(Q+aI\) is nonnegative, so its exponential is nonnegative. Hence
\[
\e^{tQ}=\e^{-at}\e^{t(Q+aI)}
\]
is nonnegative. Since \(\one^TQ=0\), differentiation gives
\[
\one^T\e^{tQ}=\one^T.
\]
\end{proof}

Let \(\{B_1,\ldots,B_m\}\) be a partition of the fine state set. Its aggregation matrix is
\[
C_{\alpha i}=\one_{\{i\in B_\alpha\}}.
\]

\begin{definition}[Canonical lumping obstruction]\label{extended:record:029}
For a linear fine transition operator \(K\), define
\[
\Ob_C(K)
=
(CK)|_{\Ker C}
\in
\mathcal B(\Ker C,\RR^m).
\]
\end{definition}

\begin{theorem}[Descent through an aggregation map]\label{extended:record:030}\label{extended:thm:canonical_lumping_obstruction}
Let \(C:X_f\to X_c\) be a surjective linear map between finite dimensional spaces. A unique operator \(\overline K:X_c\to X_c\) satisfying
\[
CK=\overline KC
\]
exists if and only if
\[
\Ob_C(K)=0.
\]
For any right inverse \(R\) of \(C\), the only possible coarse operator is
\[
\overline K=CKR,
\]
and
\[
CK-\overline KC
=
CK(I-RC).
\]
\end{theorem}

\begin{proof}
If \(CK=\overline KC\), then \(CK\) vanishes on \(\Ker C\). Conversely, suppose the obstruction vanishes and define
\[
\overline K(Cx)=CKx.
\]
If \(Cx=Cy\), then \(x-y\in\Ker C\), so the value does not depend on the representative. Surjectivity defines \(\overline K\) on all of \(X_c\), and the intertwining relation follows from the definition. It also gives uniqueness. If \(R\) is a right inverse, then
\[
\overline K=\overline KCR=CKR.
\]
Substitution gives the defect identity.
\end{proof}

\begin{corollary}[Coordinate form of the lumping obstruction]\label{extended:record:031}\label{extended:cor:lumping_obstruction_coordinates}
For the partition aggregation above, \(\Ker C\) is spanned by differences \(e_i-e_j\) with \(i\) and \(j\) in one source block. Also,
\[
\Ob_C(K)(e_i-e_j)
=
\left(
\sum_{r\in B_\alpha}(K_{ri}-K_{rj})
\right)_{\alpha=1}^m.
\]
\end{corollary}

\begin{proof}
Inside each block, the vectors with coordinate sum zero are spanned by differences of coordinate vectors. Applying \(CK\) gives the displayed aggregated column differences.
\end{proof}

\begin{theorem}[Characterization of exact discrete lumping]\label{extended:record:032}\label{extended:thm:lumping_characterization}
Let \(K\) be column stochastic. A coarse column stochastic matrix \(\overline K\) satisfying
\[
CK=\overline KC
\]
exists if and only if, for every source block \(B_\beta\), every pair \(i,j\in B_\beta\), and every target block \(B_\alpha\),
\[
\sum_{r\in B_\alpha}K_{ri}
=
\sum_{r\in B_\alpha}K_{rj}.
\]
When it exists, \(\overline K\) is unique and
\[
\overline K_{\alpha\beta}
=
\sum_{r\in B_\alpha}K_{ri}
\]
for any \(i\in B_\beta\).
\end{theorem}

\begin{proof}
The \(j\)th column of \(CK\) records the transition mass from state \(j\) into every target block. The \(j\)th column of \(\overline KC\) is the column of \(\overline K\) indexed by the source block containing \(j\). Equality therefore holds exactly under the stated condition. The formula gives uniqueness. Summing over \(\alpha\) gives one, so the coarse matrix is column stochastic.
\end{proof}

\begin{corollary}[Continuous time lumping]\label{extended:record:033}\label{extended:cor:continuous_lumping}
A coarse Markov generator \(\overline Q\) satisfying
\[
CQ=\overline QC
\]
exists if and only if the aggregated rates
\[
\sum_{r\in B_\alpha}Q_{ri}
\]
depend only on the source block of \(i\). In that case
\[
C\e^{tQ}=\e^{t\overline Q}C
\]
for every \(t\geq0\).
\end{corollary}

\begin{proof}
The characterization is the same column argument as in \cref{extended:thm:lumping_characterization}. The off diagonal entries of \(\overline Q\) are nonnegative and its column sums are zero. Exact semigroup intertwining follows from \cref{extended:prop:infinitesimal_defect}.
\end{proof}

\begin{theorem}[Finite horizon lumping defect]\label{extended:record:034}\label{extended:thm:finite_horizon_lumping}
For arbitrary fine and coarse operators \(K\) and \(\overline K\), with the fixed aggregation map \(C\), set
\[
\Delta_K=CK-\overline KC.
\]
Then
\[
CK^n-\overline K^nC
=
\sum_{j=0}^{n-1}
\overline K^{n-1-j}\Delta_KK^j.
\]
If \(K\) and \(\overline K\) are column stochastic, then
\[
\norm{CK^n-\overline K^nC}_{1\to1}
\leq
n\norm{\Delta_K}_{1\to1}.
\]
For every probability vector \(p\), the total variation discrepancy is at most
\[
\frac n2\norm{\Delta_K}_{1\to1}.
\]
\end{theorem}

\begin{proof}
Sum the telescoping identities
\[
\overline K^{n-1-j}CK^{j+1}
-
\overline K^{n-j}CK^j
=
\overline K^{n-1-j}\Delta_KK^j.
\]
Stochastic matrices are contractions in the \(\ell^1\) norm. The last statement uses total variation equal to one half of the \(\ell^1\) distance.
\end{proof}

\begin{corollary}[Stationary mass descends]\label{extended:record:035}\label{extended:cor:stationary_descends}
If \(CK=\overline KC\) and \(Kp=p\), then
\[
\overline K(Cp)=Cp.
\]
\end{corollary}

\begin{proof}
Apply the intertwining identity to \(p\).
\end{proof}

\begin{figure}[t]
\centering
\includegraphics[width=.72\textwidth]{markov_lumping.pdf}
\caption{Exact lumping agrees at every horizon. A small violation of the block transition condition produces a visible but controlled discrepancy.}
\label{extended:fig:markov_lumping}
\end{figure}

\section{Conditional expectation and martingale packets}\label{extended:section:06}\label{extended:sec:conditional_expectation}

Let \(\Omega=\{1,\ldots,n\}\) have strictly positive probabilities \(\pi_i\). For a partition \(\mathcal Q\), let \(B(i)\) denote the block containing \(i\). Define
\[
(P_\mathcal Qf)_i
=
\frac{\sum_{j\in B(i)}\pi_jf_j}
{\sum_{j\in B(i)}\pi_j}.
\]
Equip \(\CC^n\) with
\[
\ip{f}{g}_\pi
=
\sum_i\pi_if_i\overline{g_i}.
\]

\begin{theorem}[Finite conditional expectation]\label{extended:record:036}\label{extended:thm:finite_conditional_expectation}
The operator \(P_\mathcal Q\) is the orthogonal projector onto the functions that are constant on every block of \(\mathcal Q\). In particular,
\[
P_\mathcal Q^2=P_\mathcal Q,
\qquad
P_\mathcal Q^*=P_\mathcal Q.
\]
\end{theorem}

\begin{proof}
One application produces a block constant function, so a second application changes nothing. Expanding both weighted inner products block by block gives
\[
\ip{P_\mathcal Qf}{g}_\pi
=
\sum_{B\in\mathcal Q}
\frac{
\left(\sum_{i\in B}\pi_if_i\right)
\overline{\left(\sum_{j\in B}\pi_jg_j\right)}
}{\pi(B)}
=
\ip{f}{P_\mathcal Qg}_\pi.
\]
The range is exactly the block constant subspace.
\end{proof}

\begin{theorem}[Tower law and martingale packets]\label{extended:record:037}\label{extended:thm:martingale_packets}
Let $P_0,P_1,\ldots,P_r$ come from increasingly fine partitions, so that
\[
\Ran P_0\subseteq\Ran P_1\subseteq\cdots\subseteq\Ran P_r.
\]
Define
\[
D_0=P_0,
\qquad
D_j=P_j-P_{j-1}
\quad(1\leq j\leq r).
\]
Then
\[
P_iP_j=P_{\min(i,j)},
\qquad
D_iD_j=\delta_{ij}D_j,
\qquad
P_r=\sum_{j=0}^rD_j.
\]
Every \(D_j\) is an orthogonal projector, and
\[
\norm{P_rf}_\pi^2
=
\sum_{j=0}^r\norm{D_jf}_\pi^2
\]
and
\[
\norm f_\pi^2
=
\sum_{j=0}^r\norm{D_jf}_\pi^2
+
\norm{f-P_rf}_\pi^2.
\]
\end{theorem}

\begin{proof}
A coarse average of a fine average is the coarse average. A fine average leaves a coarse constant function unchanged. Hence \(P_iP_j=P_{\min(i,j)}\). Expanding products of the differences gives the orthogonality relations. The sum telescopes to \(P_r\). Orthogonal Pythagorean identities give the norm formulas.
\end{proof}

\begin{corollary}[Law of total variance]\label{extended:record:038}\label{extended:cor:total_variance}
For a real valued \(f\), let \(P_0f=\Exp_\pi[f]\one\). Then
\[
\Var_\pi(f)
=
\Var_\pi(P_\mathcal Qf)
+
\norm{f-P_\mathcal Qf}_\pi^2.
\]
\end{corollary}

\begin{proof}
Subtract the mean and apply the orthogonal decomposition into the range of \(P_\mathcal Q\) and its orthogonal complement.
\end{proof}

\begin{figure}[t]
\centering
\includegraphics[width=.68\textwidth]{martingale_packet_energy.pdf}
\caption{Energy carried by the mean packet and three successive martingale difference packets in a finite weighted example.}
\label{extended:fig:martingale_energy}
\end{figure}

\begin{standardfact}[General conditional expectations]\label{extended:record:039}
On a probability space, conditional expectation onto a sub sigma algebra is the orthogonal projector in \(L^2\). Nested sigma algebras satisfy the same tower law. Under the usual hypotheses, martingale convergence gives a limit in \(L^2\) and almost surely. Standard treatments include \cite{Williams1991}.
\end{standardfact}

\section{Gaussian transfer and Hermite packets}\label{extended:section:07}\label{extended:sec:ou_hermite}

Let \(\gamma\) be the standard Gaussian probability measure on \(\RR\). For \(t\geq0\), define the Ornstein Uhlenbeck transfer
\[
(\Ocal_t f)(x)
=
\int_\RR
f\left(
\e^{-t}x+\sqrt{1-\e^{-2t}}\,y
\right)
\gamma(\dd y).
\]
Let \(H_n\) be the probabilists' Hermite polynomial defined by
\[
\exp\left(ux-\frac{u^2}{2}\right)
=
\sum_{n=0}^\infty H_n(x)\frac{u^n}{n!},
\]
and set \(h_n=H_n/\sqrt{n!}\).

\begin{theorem}[Ornstein Uhlenbeck semigroup]\label{extended:record:040}\label{extended:thm:ou_semigroup}
The operators \((\Ocal_t)_{t\geq0}\) form a positive Markov contraction semigroup on \(L^p(\gamma)\) for every \(1\leq p\leq\infty\). On \(L^2(\gamma)\), the semigroup is self adjoint and
\[
\Ocal_t h_n=\e^{-nt}h_n.
\]
The family \((h_n)_{n\geq0}\) is an orthonormal basis of \(L^2(\gamma)\).
\end{theorem}

\begin{proof}
Positivity and preservation of constants follow from the integral formula. If \(X\) and \(Y\) are independent standard Gaussian variables, then
\[
\e^{-s}
\left(
\e^{-t}x+\sqrt{1-\e^{-2t}}X
\right)
+
\sqrt{1-\e^{-2s}}Y
\]
has the same distribution as
\[
\e^{-(s+t)}x
+
\sqrt{1-\e^{-2(s+t)}}Z
\]
with \(Z\) standard Gaussian. This proves the semigroup law. Jensen's inequality and invariance of \(\gamma\) give the \(L^p\) contraction. The joint Gaussian pair
\[
\left(
X,
\e^{-t}X+\sqrt{1-\e^{-2t}}Y
\right)
\]
is symmetric under interchange, which proves self adjointness.

Apply \(\Ocal_t\) to the Hermite generating function. Gaussian integration gives
\[
\begin{aligned}
\Ocal_t
\exp\left(ux-\frac{u^2}{2}\right)
&=
\exp\left(
\e^{-t}ux-
\frac{\e^{-2t}u^2}{2}
\right)\\
&=
\sum_{n=0}^\infty
\e^{-nt}H_n(x)\frac{u^n}{n!}.
\end{aligned}
\]
Comparison of coefficients proves the eigenvalue formula. The same generating function shows
\[
\int H_mH_n\dd\gamma=n!\delta_{mn}.
\]
The Hermite polynomials span all polynomials, and polynomials are dense in \(L^2(\gamma)\). Hence the normalized family is an orthonormal basis. See \cite{Janson1997} for the general Gaussian Hilbert space setting.
\end{proof}

\begin{corollary}[Hermite packet reconstruction]\label{extended:record:041}\label{extended:cor:hermite_packets}
Let \(P_n\) be the orthogonal projection onto \(\CC h_n\). Then
\[
I=\sum_{n=0}^\infty P_n
\]
in the strong operator topology and
\[
\Ocal_t
=
\sum_{n=0}^\infty\e^{-nt}P_n.
\]
For \(f\in L^2(\gamma)\),
\[
\norm{\Ocal_t f}_2^2
=
\sum_{n=0}^\infty
\e^{-2nt}\norm{P_nf}_2^2.
\]
\end{corollary}

\begin{proof}
Apply Parseval's identity and the eigenvalue formula in \cref{extended:thm:ou_semigroup}.
\end{proof}

\begin{proposition}[Generator and Dirichlet form]\label{extended:record:042}\label{extended:prop:ou_generator}
On the polynomial subspace, the infinitesimal generator of \((\Ocal_t)\) is the differential operator
\[
\Lcal f=f''-xf'.
\]
It satisfies
\[
\Lcal h_n=-nh_n.
\]
For every polynomial \(f\),
\[
\ip{-\Lcal f}{f}_{L^2(\gamma)}
=
\int_\RR\abs{f'(x)}^2\gamma(\dd x).
\]
\end{proposition}

\begin{proof}
Write
\[
\rho_t=\e^{-t}=1-t+O(t^2),
\qquad
\sigma_t^2=1-\e^{-2t}=2t+O(t^2).
\]
For a polynomial \(f\), Taylor expansion at \(x\) and Gaussian integration give
\[
\begin{aligned}
\Ocal_tf(x)
={}&f(x)
+f'(x)(\rho_t-1)x
+\frac12f''(x)\sigma_t^2
+o(t)\\
={}&f(x)+t\bigl(f''(x)-xf'(x)\bigr)+o(t).
\end{aligned}
\]
Every omitted term is a polynomial in the Gaussian variable multiplied by a quantity that is \(o(t)\), so the expansion may be integrated term by term. Therefore
\[
\left.\frac{\dd}{\dd t}\Ocal_tf\right|_{t=0}
=f''-xf'.
\]
The Hermite eigenvalue identity also follows by differentiating
\(
\Ocal_th_n=\e^{-nt}h_n
\)
at zero. Gaussian integration by parts gives
\[
\int_\RR f'(x)\overline{g(x)}\gamma(\dd x)
=
\int_\RR f(x)\overline{xg(x)-g'(x)}\gamma(\dd x).
\]
Set \(g=f'\) and rearrange.
\end{proof}

\begin{corollary}[Gaussian spectral gap]\label{extended:record:043}\label{extended:cor:ou_spectral_gap}
If \(f\in L^2(\gamma)\) has mean zero, then
\[
\norm{\Ocal_tf}_2
\leq
\e^{-t}\norm f_2.
\]
For every polynomial \(f\),
\[
\Var_\gamma(f)
\leq
\int_\RR\abs{f'(x)}^2\gamma(\dd x).
\]
\end{corollary}

\begin{proof}
A mean zero function has no \(h_0\) component, so every surviving Hermite multiplier is at most \(\e^{-t}\). For a polynomial
\[
f-\Exp_\gamma f=\sum_{n\geq1}c_nh_n.
\]
Then
\[
\Var_\gamma(f)=\sum_{n\geq1}\abs{c_n}^2
\leq
\sum_{n\geq1}n\abs{c_n}^2
=
\ip{-\Lcal f}{f}.
\]
Apply \cref{extended:prop:ou_generator}.
\end{proof}

\begin{corollary}[Fractional Hermite transfer]\label{extended:record:044}\label{extended:cor:fractional_hermite}
Let \(-\Lcal\) be the nonnegative Ornstein Uhlenbeck number operator with
\[
-\Lcal h_n=nh_n.
\]
For every \(\alpha>0\),
\[
\e^{-t(-\Lcal)^\alpha}f
=
\sum_{n=0}^\infty
\e^{-tn^\alpha}P_nf.
\]
This is a contraction semigroup on \(L^2(\gamma)\).
\end{corollary}

\begin{proof}
The spectral theorem applies to the orthonormal Hermite basis. Every multiplier has modulus at most one and obeys the scalar semigroup law.
\end{proof}

\begin{figure}[t]
\centering
\includegraphics[width=.72\textwidth]{ou_hermite.pdf}
\caption{Hermite packet multipliers for the Ornstein Uhlenbeck semigroup and two fractional powers of its number operator.}
\label{extended:fig:ou_hermite}
\end{figure}

\section{Bayesian transfer and evidence packets}\label{extended:section:08}\label{extended:sec:bayesian_transfer}

Let
\[
\Delta_{n-1}
=
\left\{p\in\RR^n:
p_i\geq0,
\quad
\sum_ip_i=1
\right\}
\]
be the probability simplex. Let \(\lambda\in[0,\infty)^n\) have positive evidence
\[
Z_\lambda(p)=\sum_i\lambda_ip_i.
\]
The Bayesian transfer is
\[
\Bcal_\lambda(p)
=
\frac{\lambda\odot p}{Z_\lambda(p)}.
\]

\begin{proposition}[Scale invariance and composition]\label{extended:record:045}\label{extended:prop:bayes_composition}
For every \(c>0\),
\[
\Bcal_{c\lambda}=\Bcal_\lambda.
\]
If all relevant evidences are positive, then
\[
\Bcal_\mu(\Bcal_\lambda(p))
=
\Bcal_{\lambda\odot\mu}(p).
\]
\end{proposition}

\begin{proof}
A common likelihood scale cancels between the numerator and denominator. In the second identity, the first evidence appears in both the numerator and the denominator of the second update and cancels.
\end{proof}

Let \(I=I_1\sqcup\cdots\sqcup I_m\) be a partition. Define packet evidences
\[
Z_\alpha
=
\sum_{i\in I_\alpha}\lambda_ip_i.
\]
For \(Z_\alpha>0\), define
\[
m_\alpha=\frac{Z_\alpha}{\sum_\beta Z_\beta},
\qquad
q_{\alpha,i}=\frac{\lambda_ip_i}{Z_\alpha}
\quad(i\in I_\alpha).
\]
Packets with zero evidence receive posterior mass zero.

\begin{theorem}[Bayesian packet reconstruction]\label{extended:record:046}\label{extended:thm:bayesian_packet_reconstruction}
The posterior is the disjoint packet synthesis
\[
\Bcal_\lambda(p)
=
\sum_{\alpha:Z_\alpha>0}
m_\alpha q_\alpha.
\]
The local shape \(q_\alpha\) depends only on the prior and likelihood inside packet \(I_\alpha\). Coupling between different packets occurs through the normalization of the packet evidences.
\end{theorem}

\begin{proof}
For \(i\in I_\alpha\),
\[
m_\alpha q_{\alpha,i}
=
\frac{Z_\alpha}{\sum_\beta Z_\beta}
\frac{\lambda_ip_i}{Z_\alpha}
=
\frac{\lambda_ip_i}{Z_\lambda(p)}.
\]
The packet supports are disjoint.
\end{proof}

\begin{theorem}[Likelihood perturbation]\label{extended:record:047}\label{extended:thm:likelihood_perturbation}
Let \(\lambda,\mu\geq0\) have positive evidences under a fixed prior \(p\). Set
\[
d_p(\lambda,\mu)
=
\sum_i\abs{\lambda_i-\mu_i}p_i.
\]
Then
\[
\norm{\Bcal_\lambda(p)-\Bcal_\mu(p)}_1
\leq
\frac{2d_p(\lambda,\mu)}
{\max\{Z_\lambda(p),Z_\mu(p)\}}.
\]
\end{theorem}

\begin{proof}
Add and subtract \(\mu\odot p/Z_\lambda(p)\). This gives
\[
\norm{\Bcal_\lambda(p)-\Bcal_\mu(p)}_1
\leq
\frac{2d_p(\lambda,\mu)}{Z_\lambda(p)}.
\]
Interchanging \(\lambda\) and \(\mu\) gives the same estimate with \(Z_\mu(p)\). Take the smaller of the two upper bounds.
\end{proof}

\begin{theorem}[Stability with respect to the prior]\label{extended:record:048}\label{extended:thm:prior_stability}
Assume
\[
0<a\leq\lambda_i\leq b
\]
for every \(i\). Then
\[
\norm{\Bcal_\lambda(p)-\Bcal_\lambda(q)}_1
\leq
\frac{2b}{a}\norm{p-q}_1
\]
for all priors \(p,q\in\Delta_{n-1}\).
\end{theorem}

\begin{proof}
The evidences are at least \(a\). Add and subtract \(\lambda\odot q/Z_\lambda(p)\). The numerator difference is at most \(b\norm{p-q}_1\), and the evidence difference has the same bound. The two terms give the stated constant.
\end{proof}

\begin{figure}[t]
\centering
\includegraphics[width=.72\textwidth]{bayesian_packets.pdf}
\caption{A prior and posterior on two coarse packets. The likelihood is constant in the first packet and varies inside the second, so only the second local conditional distribution changes.}
\label{extended:fig:bayesian_packets}
\end{figure}

\section{Bayesian coarsening and finite sufficiency}\label{extended:section:09}\label{extended:sec:bayesian_sufficiency}

A fixed coarse description will be called sufficient when it reproduces the coarse posterior masses for every positive prior and leaves the conditional distribution inside every packet unchanged. This excludes a coarse rule that hides an unrecorded local update.

Let \(C\) send a fine probability vector to its block masses
\[
(Cp)_\alpha
=
\sum_{i\in I_\alpha}p_i.
\]
Assume first that every block mass is positive. Write
\[
r_{\alpha,i}
=
\frac{p_i}{(Cp)_\alpha}
\]
for the prior conditional distribution inside packet \(I_\alpha\). Define the prior dependent packet average likelihood
\[
\overline\lambda_\alpha(p)
=
\sum_{i\in I_\alpha}\lambda_ir_{\alpha,i}.
\]

\begin{theorem}[Exact two stage Bayesian update]\label{extended:record:049}\label{extended:thm:two_stage_bayes}
For every strictly positive prior,
\[
C\Bcal_\lambda(p)
=
\Bcal_{\overline\lambda(p)}(Cp).
\]
For every packet with \(\overline\lambda_\alpha(p)>0\), the posterior conditional distribution is
\[
q_{\alpha,i}
=
\frac{\lambda_ir_{\alpha,i}}
{\overline\lambda_\alpha(p)}.
\]
A packet with zero average likelihood has zero posterior mass. Thus a fine Bayesian update consists of a coarse update by packet average likelihoods followed by independent local updates inside the packets that survive.
\end{theorem}

\begin{proof}
The posterior mass of packet \(I_\alpha\) is
\[
\frac{\sum_{i\in I_\alpha}\lambda_ip_i}
{\sum_j\lambda_jp_j}
=
\frac{(Cp)_\alpha\overline\lambda_\alpha(p)}
{\sum_\beta(Cp)_\beta\overline\lambda_\beta(p)}.
\]
This is the coarse Bayesian formula. Dividing the fine posterior coordinates by the packet mass gives the local formula.
\end{proof}

\begin{theorem}[Characterization of fixed coarse sufficiency]\label{extended:record:050}\label{extended:thm:bayesian_sufficiency}
Assume \(\lambda_i>0\) for every state. The following conditions are equivalent.
\begin{enumerate}
\item There is a positive coarse likelihood \(\ell\in(0,\infty)^m\) such that
\[
C\Bcal_\lambda(p)
=
\Bcal_\ell(Cp)
\]
for every strictly positive prior \(p\), and the posterior conditional distribution inside every packet equals its prior conditional distribution.
\item The likelihood is constant on every packet.
\end{enumerate}
When these conditions hold, \(\ell_\alpha\) can be chosen as the common value of \(\lambda\) on \(I_\alpha\), and it is unique up to a common positive scale.
\end{theorem}

\begin{proof}
If \(\lambda\) is constant on each packet, \cref{extended:thm:two_stage_bayes} has packet averages independent of the prior and its local posterior formula gives \(q_\alpha=r_\alpha\). This proves the first condition.

Conversely, fix any strictly positive prior. Equality of the local posterior and prior conditionals gives
\[
\frac{\lambda_i r_{\alpha,i}}
{\overline\lambda_\alpha(p)}
=
r_{\alpha,i}
\]
for every \(i\in I_\alpha\). Positivity of the conditional prior implies
\[
\lambda_i=\overline\lambda_\alpha(p)
\]
for every state in the packet. Thus the likelihood is constant on every packet. If two positive coarse likelihoods define the same update for every positive coarse prior, comparison of any two posterior odds shows that their coordinate ratios agree. Hence they differ only by a common positive factor.
\end{proof}

\begin{theorem}[Coarsest fixed sufficient partition]\label{extended:record:051}\label{extended:thm:coarsest_bayesian_partition}
Let \(\mathcal L\) be a nonempty family of strictly positive likelihood vectors. Define
\[
i\sim_{\mathcal L}j
\quad\Longleftrightarrow\quad
\lambda_i=\lambda_j
\quad\hbox{for every }\lambda\in\mathcal L.
\]
Let \(\mathcal Q_{\mathcal L}\) be the partition into equivalence classes. A partition \(\mathcal Q\) gives a fixed coarse sufficient description in the sense of \cref{extended:thm:bayesian_sufficiency} for every \(\lambda\in\mathcal L\) if and only if every block of \(\mathcal Q\) is contained in a block of \(\mathcal Q_{\mathcal L}\). Hence \(\mathcal Q_{\mathcal L}\) is the unique coarsest partition with this property.
\end{theorem}

\begin{proof}
Apply \cref{extended:thm:bayesian_sufficiency} to every likelihood in the family. A fixed sufficient description exists exactly when each likelihood is constant on every block of \(\mathcal Q\). This is equivalent to every block of \(\mathcal Q\) being contained in one common level set of the family. The equivalence class partition is itself admissible, and every other admissible partition refines it.
\end{proof}

\begin{definition}[Within packet information loss]\label{extended:record:052}
For a prior \(p\), define
\[
\eta_i(p)
=
\lambda_i-\overline\lambda_\alpha(p)
\quad\hbox{when }i\in I_\alpha.
\]
Then
\[
\sum_{i\in I_\alpha}r_{\alpha,i}\eta_i(p)=0.
\]
\end{definition}

\begin{proposition}[Exact local distortion]\label{extended:record:053}\label{extended:prop:local_bayes_distortion}
For every packet with positive mass and evidence,
\[
q_{\alpha,i}-r_{\alpha,i}
=
\frac{r_{\alpha,i}\eta_i(p)}
{\overline\lambda_\alpha(p)}.
\]
Consequently,
\[
\norm{q_\alpha-r_\alpha}_1
=
\frac{
\sum_{i\in I_\alpha}r_{\alpha,i}\abs{\eta_i(p)}
}{\overline\lambda_\alpha(p)}.
\]
If \(\lambda\) is positive on the packet, then
\[
\norm{q_\alpha-r_\alpha}_1
\leq
\frac{\osc_{I_\alpha}(\lambda)}
{\min_{i\in I_\alpha}\lambda_i}.
\]
\end{proposition}

\begin{proof}
Subtract \(r_{\alpha,i}\) from the local posterior formula in \cref{extended:thm:two_stage_bayes}. The norm identity follows because the factor multiplying \(\eta_i\) is nonnegative. The last estimate uses
\[
\sum_i r_{\alpha,i}\abs{\eta_i}
\leq
\osc_{I_\alpha}(\lambda)
\]
and \(\overline\lambda_\alpha(p)\geq\min_{I_\alpha}\lambda\).
\end{proof}

\begin{theorem}[Approximate coarse Bayesian rule]\label{extended:record:054}\label{extended:thm:approximate_coarse_bayes}
Choose one positive number \(\ell_\alpha\) for every packet and define the block constant fine likelihood
\[
\widetilde\lambda_i=\ell_\alpha
\quad(i\in I_\alpha).
\]
Then
\[
\norm{
C\Bcal_\lambda(p)-\Bcal_\ell(Cp)
}_1
\leq
\frac{
2\sum_i p_i\abs{\lambda_i-\widetilde\lambda_i}
}{
\max\{Z_\lambda(p),Z_{\widetilde\lambda}(p)\}
}.
\]
\end{theorem}

\begin{proof}
Block constancy gives
\[
C\Bcal_{\widetilde\lambda}(p)
=
\Bcal_\ell(Cp).
\]
Aggregation is an \(\ell^1\) contraction. Apply \cref{extended:thm:likelihood_perturbation} to the two fine posteriors and then aggregate.
\end{proof}

\begin{definition}[Relative information]\label{extended:record:055}
For probability vectors \(q\) and \(p\), define
\[
D(q\Vert p)
=
\sum_iq_i\log\frac{q_i}{p_i}.
\]
The convention is \(0\log0=0\). A term with \(q_i>0\) and \(p_i=0\) has value \(+\infty\).
\end{definition}

\begin{theorem}[Exact information decomposition]\label{extended:record:056}\label{extended:thm:bayesian_information_chain}
Let \(p\) and \(q\) be probability vectors on a finite partition \((I_\alpha)_\alpha\), and assume \(q_i=0\) whenever \(p_i=0\). Let
\[
m=Cp,
\qquad
n=Cq,
\]
and let \(p_\alpha\) and \(q_\alpha\) be the conditional distributions inside packet \(I_\alpha\) whenever the corresponding mass is positive. Then
\[
D(q\Vert p)
=
D(n\Vert m)
+
\sum_{\alpha:n_\alpha>0}
 n_\alpha D(q_\alpha\Vert p_\alpha).
\]
The first term is the information retained by the coarse masses. The second term is the exact information carried by changes inside the packets.
\end{theorem}

\begin{proof}
For \(i\in I_\alpha\), write
\[
p_i=m_\alpha p_{\alpha,i},
\qquad
q_i=n_\alpha q_{\alpha,i}.
\]
Then
\[
\log\frac{q_i}{p_i}
=
\log\frac{n_\alpha}{m_\alpha}
+
\log\frac{q_{\alpha,i}}{p_{\alpha,i}}.
\]
Multiply by \(q_i\), sum first inside each packet, and then sum over the packets. This is the finite relative entropy chain rule. Standard information theoretic treatments include \cite{CoverThomas2006}.
\end{proof}

\begin{corollary}[Coarse data processing and equality]\label{extended:record:057}\label{extended:cor:coarse_data_processing}
Under the assumptions of \cref{extended:thm:bayesian_information_chain},
\[
D(Cq\Vert Cp)
\leq
D(q\Vert p).
\]
Equality holds if and only if
\[
q_\alpha=p_\alpha
\]
for every packet with positive posterior mass.
\end{corollary}

\begin{proof}
Every conditional relative information term in \cref{extended:thm:bayesian_information_chain} is nonnegative. It vanishes exactly when the two conditional distributions agree.
\end{proof}

\begin{corollary}[Information created by a Bayesian update]\label{extended:record:058}\label{extended:cor:bayesian_information_gain}
Let \(q=\Bcal_\lambda(p)\) with positive likelihood and positive prior. Then
\[
D(q\Vert p)
=
\sum_iq_i\log\lambda_i
-
\log Z_\lambda(p).
\]
Inside a packet,
\[
D(q_\alpha\Vert r_\alpha)
=
\sum_{i\in I_\alpha}
q_{\alpha,i}
\log\frac{\lambda_i}{\overline\lambda_\alpha(p)}.
\]
Hence
\[
D(q_\alpha\Vert r_\alpha)
\leq
\log
\frac{\max_{i\in I_\alpha}\lambda_i}
{\min_{i\in I_\alpha}\lambda_i}.
\]
\end{corollary}

\begin{proof}
The posterior ratio is
\[
\frac{q_i}{p_i}
=
\frac{\lambda_i}{Z_\lambda(p)}.
\]
The local ratio is
\[
\frac{q_{\alpha,i}}{r_{\alpha,i}}
=
\frac{\lambda_i}{\overline\lambda_\alpha(p)}.
\]
Insert these identities into the definition of relative entropy. The final estimate follows because both \(\lambda_i\) and the packet average lie between the packet minimum and maximum.
\end{proof}

\begin{corollary}[Information criterion for local sufficiency]\label{extended:record:059}\label{extended:cor:information_sufficiency}
For a fixed positive likelihood, the within packet term in \cref{extended:thm:bayesian_information_chain} vanishes for every positive prior if and only if the likelihood is constant on every packet.
\end{corollary}

\begin{proof}
If the likelihood is block constant, every local posterior equals its local prior. Conversely, choose any positive prior. Vanishing relative entropy gives \(q_\alpha=r_\alpha\). The local posterior formula then forces all likelihood values in that packet to equal their packet average. Repeat for every packet.
\end{proof}

\begin{figure}[t]
\centering
\includegraphics[width=.72\textwidth]{bayesian_information_split.pdf}
\caption{Relative information separated into a coarse packet contribution and conditional contributions inside the packets.}
\label{extended:fig:bayesian_information_split}
\end{figure}

\begin{remark}\label{extended:record:060}
The exact coarse likelihood in \cref{extended:thm:two_stage_bayes} depends on the conditional prior inside each packet. \Cref{extended:thm:bayesian_sufficiency} identifies when that dependence disappears. \Cref{extended:thm:bayesian_information_chain} measures what is lost when only the coarse posterior masses are retained.
\end{remark}

\section{Bayesian flows and hidden Markov recursion}\label{extended:section:10}\label{extended:sec:bayesian_flows}

Fix a score vector \(g\in\RR^n\). For \(t\geq0\), define
\[
\Phi_t(p)_i
=
\frac{\e^{tg_i}p_i}
{\sum_j\e^{tg_j}p_j}.
\]

\begin{theorem}[Bayesian flow and replicator equation]\label{extended:record:061}\label{extended:thm:bayesian_flow}
The maps \(\Phi_t\) form a nonlinear semigroup:
\[
\Phi_{s+t}=\Phi_s\circ\Phi_t.
\]
Every coordinate satisfies
\[
\frac{\dd}{\dd t}p_i(t)
=
\bigl(g_i-\ip{g}{p(t)}\bigr)p_i(t).
\]
The support of \(p(t)\) is independent of \(t\), and for any two coordinates in the support,
\[
\frac{p_i(t)}{p_j(t)}
=
\e^{t(g_i-g_j)}
\frac{p_i(0)}{p_j(0)}.
\]
\end{theorem}

\begin{proof}
The semigroup law follows from multiplication of the likelihoods \(\e^{sg}\) and \(\e^{tg}\). Differentiating the normalized formula gives the replicator equation. A coordinate that starts at zero remains zero, while a positive coordinate remains positive. Dividing two explicit coordinate formulas gives the odds identity.
\end{proof}

\begin{theorem}[Evidence cumulants and score monotonicity]\label{extended:record:062}\label{extended:thm:evidence_cumulants}
Set
\[
\psi(t)
=
\log\left(
\sum_i\e^{tg_i}p_i(0)
\right).
\]
Then
\[
\psi'(t)=\ip{g}{p(t)},
\qquad
\psi''(t)=\Var_{p(t)}(g)\geq0.
\]
Consequently,
\[
\frac{\dd}{\dd t}\ip{g}{p(t)}
=
\Var_{p(t)}(g).
\]
\end{theorem}

\begin{proof}
Differentiate the log normalizing constant once and twice. The second derivative is the centered second moment under the tilted distribution.
\end{proof}

\begin{corollary}[Concentration on the maximal score packet]\label{extended:record:063}\label{extended:cor:bayes_flow_concentration}
Let
\[
G=\max\{g_i:p_i(0)>0\}
\]
and let \(M\) be the set of supported indices with score \(G\). If \(M\) is the full support, then the conclusion below is immediate. Otherwise, suppose
\[
\delta
=
\min\{G-g_i:p_i(0)>0,
\quad i\notin M\}>0.
\]
Then
\[
1-\sum_{i\in M}p_i(t)
\leq
\frac{1-p_M(0)}{p_M(0)}\e^{-\delta t},
\]
where \(p_M(0)=\sum_{i\in M}p_i(0)\).
\end{corollary}

\begin{proof}
Factor \(\e^{tG}\) from the normalizing constant. The total weight outside \(M\) is at most \((1-p_M(0))\e^{-\delta t}\), while the weight on \(M\) is \(p_M(0)\).
\end{proof}

\begin{figure}[t]
\centering
\includegraphics[width=.72\textwidth]{bayesian_flow.pdf}
\caption{A finite Bayesian flow. Mass concentrates on the states with the largest score, while every relative odds ratio follows an exact exponential law.}
\label{extended:fig:bayesian_flow}
\end{figure}

\subsection{Prediction and update}

Let \(K\) be column stochastic and let \(\lambda^{(r)}\) be the likelihood of the \(r\)th observation. The hidden Markov filter is
\[
\widehat p_r=Kp_{r-1},
\qquad
p_r=\Bcal_{\lambda^{(r)}}(\widehat p_r).
\]
A standard account of finite hidden Markov models appears in \cite{Rabiner1989}.

\begin{proposition}[Simplex preservation]\label{extended:record:064}\label{extended:prop:hmm_preservation}
Suppose \(p_{r-1}\) is a probability vector and the new evidence is positive. Prediction produces the probability vector \(\widehat p_r\), and normalization produces the probability vector \(p_r\).
\end{proposition}

\begin{proof}
Column stochasticity preserves nonnegativity and total mass. Bayesian normalization then preserves the same properties.
\end{proof}

\begin{theorem}[One step filter stability]\label{extended:record:065}\label{extended:thm:hmm_stability}
Let \(K\) and \(\widetilde K\) be column stochastic. Assume
\[
0<a\leq\lambda_i\leq b.
\]
Then
\[
\norm{
\Bcal_\lambda(Kp)-\Bcal_\lambda(\widetilde Kq)
}_1
\leq
\frac{2b}{a}
\left(
\norm{p-q}_1
+
\norm{K-\widetilde K}_{1\to1}
\right).
\]
\end{theorem}

\begin{proof}
Apply \cref{extended:thm:prior_stability} to \(Kp\) and \(\widetilde Kq\). Then use
\[
\norm{Kp-\widetilde Kq}_1
\leq
\norm{K(p-q)}_1
+
\norm{(K-\widetilde K)q}_1
\]
and the stochastic contraction property.
\end{proof}

\begin{example}[Beta binomial transfer]\label{extended:record:066}
A beta prior with parameters \((a,b)\) and binomial data containing \(s\) successes in \(n\) trials gives the posterior parameters
\[
(a+s,b+n-s).
\]
The update is additive in the two exponent coordinates. This is a finite dimensional conjugate realization of likelihood composition. Standard Bayesian references include \cite{BernardoSmith1994}.
\end{example}

\section{Fractional functional calculus}\label{extended:section:11}\label{extended:sec:fractional}

Let \(H\) be a finite dimensional Hilbert space and let \(L=L^*\geq0\). Write its spectral resolution as
\[
L=\sum_j\lambda_jP_j.
\]
For \(\alpha>0\), define
\[
L^\alpha
=
\sum_j\lambda_j^\alpha P_j
\]
and
\[
T_\alpha(t)
=
\e^{-tL^\alpha}
=
\sum_j\e^{-t\lambda_j^\alpha}P_j.
\]

\begin{theorem}[Fractional packet calculus]\label{extended:record:067}\label{extended:thm:fractional_packet_calculus}
For \(\alpha,\beta>0\),
\[
L^\alpha L^\beta=L^{\alpha+\beta}.
\]
For \(s,t\geq0\),
\[
T_\alpha(s)T_\alpha(t)=T_\alpha(s+t).
\]
Every spectral packet \(P_jH\) reduces \(L^\alpha\) and \(T_\alpha(t)\).
\end{theorem}

\begin{proof}
The claims hold on each spectral packet and the projectors are mutually orthogonal.
\end{proof}

\begin{theorem}[Balakrishnan formula in finite dimension]\label{extended:record:068}\label{extended:thm:balakrishnan}
For \(0<\alpha<1\),
\[
L^\alpha
=
\frac{\sin(\pi\alpha)}{\pi}
\int_0^\infty
s^{\alpha-1}L(sI+L)^{-1}
\dd s.
\]
The integral converges in operator norm.
\end{theorem}

\begin{proof}
It is enough to prove the formula on an eigenvector with eigenvalue \(x\geq0\). For \(x>0\), the substitution \(s=xu\) gives
\[
\frac{\sin(\pi\alpha)}{\pi}
\int_0^\infty
s^{\alpha-1}\frac{x}{s+x}\dd s
=x^\alpha
\frac{\sin(\pi\alpha)}{\pi}
\int_0^\infty\frac{u^{\alpha-1}}{1+u}\dd u
=x^\alpha.
\]
The value at \(x=0\) is zero. There are finitely many eigenvalues, so the operator integral converges. This is the finite spectral form of the fractional power construction in \cite{Balakrishnan1960}.
\end{proof}

\begin{theorem}[Exact fractional naturality]\label{extended:record:069}\label{extended:thm:exact_fractional_naturality}
Let \(L_f\) and \(L_c\) be finite dimensional positive self adjoint operators and let \(C\) satisfy
\[
CL_f=L_cC.
\]
Then, for every \(\alpha>0\) and \(t\geq0\),
\[
CL_f^\alpha=L_c^\alpha C
\]
and
\[
C\e^{-tL_f^\alpha}
=
\e^{-tL_c^\alpha}C.
\]
\end{theorem}

\begin{proof}
The union of the two finite spectra is a finite set. There is a polynomial that agrees with \(x^\alpha\) on that set, and another polynomial that agrees with \(\e^{-tx^\alpha}\). Exact intertwining passes to every polynomial by induction.
\end{proof}

\begin{theorem}[Fractional defect bound]\label{extended:record:070}\label{extended:thm:fractional_defect}
Let \(0<\alpha<1\). Suppose
\[
L_f\geq mI,
\qquad
L_c\geq mI
\]
for some \(m>0\). Define
\[
\Delta_L=CL_f-L_cC.
\]
Then
\[
\norm{CL_f^\alpha-L_c^\alpha C}
\leq
\alpha m^{\alpha-1}\norm{\Delta_L}.
\]
Consequently,
\[
\norm{
C\e^{-tL_f^\alpha}
-
\e^{-tL_c^\alpha}C
}
\leq
\alpha t m^{\alpha-1}\norm{\Delta_L}.
\]
\end{theorem}

\begin{proof}
Set
\[
R_f(s)=(sI+L_f)^{-1},
\qquad
R_c(s)=(sI+L_c)^{-1}.
\]
The resolvent identity gives
\[
CR_f(s)-R_c(s)C
=
-R_c(s)\Delta_LR_f(s).
\]
Since \(LR=I-sR\),
\[
CL_fR_f(s)-L_cR_c(s)C
=
sR_c(s)\Delta_LR_f(s).
\]
Insert this into \cref{extended:thm:balakrishnan}. The spectral lower bound gives
\[
\norm{R_f(s)},\norm{R_c(s)}
\leq
\frac{1}{s+m}.
\]
Therefore
\[
\begin{aligned}
\norm{CL_f^\alpha-L_c^\alpha C}
&\leq
\frac{\sin(\pi\alpha)}{\pi}
\norm{\Delta_L}
\int_0^\infty
\frac{s^\alpha}{(s+m)^2}\dd s\\
&=
\alpha m^{\alpha-1}\norm{\Delta_L}.
\end{aligned}
\]
The semigroup estimate follows from \cref{extended:thm:duhamel_defect}, applied to the contraction generators \(-L_f^\alpha\) and \(-L_c^\alpha\).
\end{proof}

\begin{remark}\label{extended:record:071}
The factor \(m^{\alpha-1}\) records the singular behavior of \(x^\alpha\) near zero. Without a common positive lower bound, a uniform Lipschitz estimate of this form cannot be expected.
\end{remark}

\begin{standardfact}[Stable subordination]\label{extended:record:072}\label{extended:fact:stable_subordination}
Let \(0<\alpha<1\) and \(t>0\). There is a probability measure \(\eta_{\alpha,t}\) on \([0,\infty)\) such that
\[
\int_0^\infty \e^{-s\lambda}\eta_{\alpha,t}(\dd s)
=
\e^{-t\lambda^\alpha}
\]
for every \(\lambda\geq0\). Consequently, for a finite dimensional positive self adjoint operator,
\[
\e^{-tL^\alpha}
=
\int_0^\infty
\e^{-sL}\eta_{\alpha,t}(\dd s)
\]
in operator norm. The measures satisfy the convolution law
\[
\eta_{\alpha,s}*\eta_{\alpha,t}
=
\eta_{\alpha,s+t}.
\]
\end{standardfact}

The existence of \(\eta_{\alpha,t}\) follows from complete monotonicity of \(\e^{-t\lambda^\alpha}\). The operator identity is obtained by applying the scalar Laplace transform on every spectral packet. This is the finite spectral form of Bochner subordination. See \cite{SchillingSongVondracek2012}.

\begin{theorem}[Poisson subordination]\label{extended:record:073}\label{extended:thm:poisson_subordination}
For \(t>0\),
\[
\e^{-t\sqrt L}
=
\frac{t}{2\sqrt\pi}
\int_0^\infty
s^{-3/2}
\e^{-t^2/(4s)}
\e^{-sL}
\dd s.
\]
The integral converges in operator norm.
\end{theorem}

\begin{proof}
On a spectral packet with eigenvalue \(\lambda\), the formula reduces to
\[
\e^{-t\sqrt\lambda}
=
\frac{t}{2\sqrt\pi}
\int_0^\infty
s^{-3/2}\e^{-t^2/(4s)-\lambda s}\dd s.
\]
This scalar identity follows from the Gaussian integral. Recombine the finitely many spectral packets.
\end{proof}

\begin{figure}[t]
\centering
\includegraphics[width=.72\textwidth]{fractional_multipliers.pdf}
\caption{Fractional semigroup multipliers at time one. Smaller exponents preserve more of the high spectral range.}
\label{extended:fig:fractional_multipliers}
\end{figure}

\begin{remark}\label{extended:record:074}
For sectorial and nonnormal operators, fractional powers require a resolvent or contour calculus together with domain control. The finite self adjoint theory proved here does not remove those analytic requirements. A broader account appears in \cite{SchillingSongVondracek2012}.
\end{remark}

\section{Green operators and finite difference transfer}\label{extended:section:12}\label{extended:sec:green_pde}

Consider
\[
-u''=f
\quad\hbox{on }(0,1),
\qquad
u(0)=u(1)=0.
\]
The continuum Green kernel is
\[
G(x,y)=\min(x,y)-xy.
\]
Thus
\[
u(x)=\int_0^1G(x,y)f(y)\dd y.
\]

For \(n\) interior nodes, let
\[
h=\frac{1}{n+1},
\qquad
x_i=ih,
\]
and define
\[
L_h
=
\frac{1}{h^2}
\begin{pmatrix}
2&-1&&\\
-1&2&\ddots&\\
&\ddots&\ddots&-1\\
&&-1&2
\end{pmatrix}.
\]

\begin{theorem}[Exact nodal Green reconstruction]\label{extended:record:075}\label{extended:thm:nodal_green}
For \(1\leq i,j\leq n\),
\[
\frac{1}{h}(L_h^{-1})_{ij}
=
G(x_i,x_j).
\]
\end{theorem}

\begin{proof}
The inverse of the unscaled tridiagonal matrix has entry
\[
\frac{\min(i,j)(n+1-\max(i,j))}{n+1}.
\]
Multiplying by \(h^2\) and dividing by \(h\) gives
\[
h\min(i,j)
\left(1-\frac{\max(i,j)}{n+1}\right)
=
\min(x_i,x_j)-x_ix_j.
\]
\end{proof}

\begin{corollary}[Discrete Green quadrature]\label{extended:record:076}\label{extended:cor:discrete_green_quadrature}
If \(u^h\) solves
\[
L_hu^h=f^h,
\]
then every nodal value has the exact representation
\[
u_i^h
=
h\sum_{j=1}^nG(x_i,x_j)f_j^h.
\]
Thus the discrete Green operator is the nodal rectangle quadrature of the continuum Green kernel.
\end{corollary}

\begin{proof}
By \cref{extended:thm:nodal_green},
\[
(L_h^{-1})_{ij}=hG(x_i,x_j).
\]
Multiply the inverse matrix by the source vector.
\end{proof}

\begin{proposition}[Discrete maximum bound]\label{extended:record:077}\label{extended:prop:discrete_inverse_bound}
The inverse \(L_h^{-1}\) has nonnegative entries and
\[
\norm{L_h^{-1}}_{\infty\to\infty}
=
\max_{1\leq i\leq n}
\frac{x_i(1-x_i)}{2}
\leq
\frac18.
\]
\end{proposition}

\begin{proof}
Nonnegativity follows from the entry formula in \cref{extended:thm:nodal_green}. Let \(v_i=x_i(1-x_i)/2\). Since the second difference is exact on quadratics,
\[
L_hv=\one.
\]
Hence the \(i\)th row sum of \(L_h^{-1}\) is \(v_i\). For a nonnegative matrix, the infinity operator norm is the largest row sum.
\end{proof}

\begin{theorem}[Second order Poisson solution error]\label{extended:record:078}\label{extended:thm:poisson_error}
Suppose \(u\in C^4[0,1]\) solves the Dirichlet problem above. Let
\[
f^h_i=f(x_i)
\]
and let \(u^h\) solve
\[
L_hu^h=f^h.
\]
Then
\[
\max_i\abs{u^h_i-u(x_i)}
\leq
\frac{h^2}{96}
\norm{u^{(4)}}_\infty.
\]
\end{theorem}

\begin{proof}
Taylor expansion with remainder gives
\[
\norm{L_h(u(x_i))_i-f^h}_\infty
\leq
\frac{h^2}{12}\norm{u^{(4)}}_\infty.
\]
Apply \(L_h^{-1}\) and use \cref{extended:prop:discrete_inverse_bound}.
\end{proof}

\begin{figure}[t]
\centering
\includegraphics[width=.72\textwidth]{green_kernel.pdf}
\caption{The continuum Dirichlet Green section and the scaled discrete inverse. The nodal values agree exactly.}
\label{extended:fig:green_kernel}
\end{figure}

The discrete eigenpairs are
\[
\lambda_{k,h}
=
\frac{4}{h^2}
\sin^2\left(\frac{k\pi h}{2}\right),
\qquad
v_{k,h}(i)=\sin(k\pi x_i).
\]

\begin{theorem}[Low mode eigenvalue control]\label{extended:record:079}\label{extended:thm:low_mode}
For \(kh\leq1\),
\[
0\leq(k\pi)^2-\lambda_{k,h}
\leq
\frac{(k\pi)^4}{12}h^2.
\]
Consequently, for every fixed \(k\),
\[
\lambda_{k,h}=(k\pi)^2+O(h^2).
\]
\end{theorem}

\begin{proof}
Set \(x=k\pi h/2\). The condition \(kh\leq1\) gives \(0\leq x\leq\pi/2\). On this interval,
\[
\sin x\geq x-\frac{x^3}{6},
\]
so
\[
0\leq x^2-\sin^2x\leq\frac{x^4}{3}.
\]
Multiply by \(4/h^2\).
\end{proof}

\begin{corollary}[Heat and fractional mode transfer]\label{extended:record:080}\label{extended:cor:mode_transfer}
Let \(\lambda_k=(k\pi)^2\). For \(t\geq0\),
\[
0\leq
\e^{-t\lambda_{k,h}}-\e^{-t\lambda_k}
\leq
\frac{t(k\pi)^4}{12}h^2\e^{-t\lambda_{k,h}}.
\]
For \(0<\alpha\leq1\) and fixed \(k\),
\[
\lambda_{k,h}^\alpha
=
\lambda_k^\alpha+O(h^2)
\]
and
\[
\e^{-t\lambda_{k,h}^\alpha}
=
\e^{-t\lambda_k^\alpha}+O(th^2)
\]
uniformly for \(t\) in bounded intervals.
\end{corollary}

\begin{proof}
The ordinary heat estimate follows from the mean value theorem and \cref{extended:thm:low_mode}. For fixed \(k\), the eigenvalues stay uniformly away from zero once \(h\) is small. The function \(x^\alpha\) is then Lipschitz on a common compact interval containing both eigenvalues. Apply the mean value theorem twice.
\end{proof}

Finite difference schemes and their stability theory are treated more broadly in \cite{Strikwerda2004}. The present result uses the special one dimensional Green structure to obtain an explicit constant.

\section{Finite state Feynman Kac transfer}\label{extended:section:13}\label{extended:sec:feynman_kac}

Let \(K\) be column stochastic. Let \(V\) be a real potential on the states and let
\[
D_V=\diag(V).
\]
For functions written as columns, the backward Markov operator is \(K^T\).

\begin{theorem}[Discrete Feynman Kac formula]\label{extended:record:081}\label{extended:thm:discrete_feynman_kac}
Define
\[
M_V=\diag(\e^{-V})K^T.
\]
For terminal data \(f\),
\[
(M_V^nf)_i
=
\Exp_i\left[
\exp\left(
-\sum_{r=0}^{n-1}V(X_r)
\right)
f(X_n)
\right].
\]
\end{theorem}

\begin{proof}
For one step,
\[
(M_Vf)_i
=
\e^{-V(i)}\sum_jK_{ji}f_j.
\]
Iterating expands the matrix product over all paths. The product of transition entries is the path probability, and the diagonal factors give the exponential potential weight.
\end{proof}

\begin{theorem}[Continuous time finite state formula]\label{extended:record:082}\label{extended:thm:continuous_feynman_kac}
Let \(Q\) be a finite Markov generator and let \((X_t)_{t\geq0}\) be the corresponding chain. Then
\[
\left(
\e^{t(Q^T-D_V)}f
\right)_i
=
\Exp_i\left[
\exp\left(
-\int_0^tV(X_s)\dd s
\right)
f(X_t)
\right].
\]
If \(V\geq v_0\), then
\[
\norm{\e^{t(Q^T-D_V)}}_{\infty\to\infty}
\leq
\e^{-v_0t}.
\]
\end{theorem}

\begin{proof}
Conditioning over the first small time interval gives the backward equation
\[
\dot u=(Q^T-D_V)u,
\qquad
u(0)=f.
\]
The finite system has a unique solution, namely the displayed matrix exponential. The expectation bound follows from
\[
\exp\left(-\int_0^tV(X_s)\dd s\right)
\leq
\e^{-v_0t}.
\]
This finite argument is the matrix form of the Feynman Kac principle originating in \cite{Kac1949}.
\end{proof}

\begin{theorem}[Alternating propagation and killing]\label{extended:record:083}\label{extended:thm:feynman_kac_product}
For every \(t\geq0\),
\[
\e^{t(Q^T-D_V)}
=
\lim_{n\to\infty}
\left(
\e^{-tD_V/n}\e^{tQ^T/n}
\right)^n
\]
in operator norm. Each finite product alternates one Markov propagation step with one diagonal survival weight.
\end{theorem}

\begin{proof}
Set \(B=-D_V\) and \(C=Q^T\). For bounded matrices,
\[
\e^{tB/n}\e^{tC/n}
=
I+\frac{t}{n}(B+C)+O(n^{-2})
\]
uniformly in operator norm. The standard telescoping comparison with
\(
(I+t(B+C)/n)^n
\)
shows that the product converges to \(\e^{t(B+C)}\). This is the displayed factor order, so the limit follows.
\end{proof}

\begin{corollary}[Positivity and composition]\label{extended:record:084}\label{extended:cor:feynman_kac_semigroup}
The operators
\[
F_V(t)=\e^{t(Q^T-D_V)}
\]
form a positive semigroup. If \(f\geq0\), then \(F_V(t)f\geq0\).
\end{corollary}

\begin{proof}
Both claims follow from the expectation formula.
\end{proof}

\section{Schur Feshbach elimination}\label{extended:section:14}\label{extended:sec:feshbach}

Let \(X=X_P\oplus X_Q\) and write
\[
A=
\begin{pmatrix}
A_{PP}&A_{PQ}\\
A_{QP}&A_{QQ}
\end{pmatrix}.
\]
For \(zI-A_{QQ}\) invertible, define
\[
\Sigma_P(z)
=
A_{PQ}(zI-A_{QQ})^{-1}A_{QP}
\]
and
\[
F_P(z)
=
zI-A_{PP}-\Sigma_P(z).
\]

\begin{theorem}[Schur Feshbach resolvent]\label{extended:record:085}\label{extended:thm:feshbach_resolvent}
Assume \(zI-A_{QQ}\) is invertible. Then \(zI-A\) is invertible if and only if \(F_P(z)\) is invertible. In that case,
\[
P(zI-A)^{-1}P
=
F_P(z)^{-1}.
\]
In finite dimension,
\[
\det(zI-A)
=
\det(zI-A_{QQ})\det F_P(z).
\]
\end{theorem}

\begin{proof}
Set
\[
L_z=
\begin{pmatrix}
I&-A_{PQ}(zI-A_{QQ})^{-1}\\
0&I
\end{pmatrix}
\]
and
\[
S_z=
\begin{pmatrix}
F_P(z)&0\\
-A_{QP}&zI-A_{QQ}
\end{pmatrix}.
\]
Direct multiplication gives \(zI-A=L_zS_z\). The first factor is invertible. The claims follow from triangular inversion and the determinant of the second factor. This is the finite block form of the Feshbach reduction \cite{Feshbach1958}.
\end{proof}

\subsection{Time domain memory}

For the homogeneous equation \(\dot x=Ax\), variation of constants gives
\[
x_Q(t)
=
\e^{tA_{QQ}}x_Q(0)
+
\int_0^t
\e^{(t-s)A_{QQ}}A_{QP}x_P(s)
\dd s.
\]
Substitution into the retained equation gives
\[
\begin{aligned}
\dot x_P(t)
={}&
A_{PP}x_P(t)
+
A_{PQ}\e^{tA_{QQ}}x_Q(0)\\
&+
\int_0^t
\Kcal(t-s)x_P(s)
\dd s,
\end{aligned}
\]
where
\[
\Kcal(t)
=
A_{PQ}\e^{tA_{QQ}}A_{QP}.
\]

\begin{theorem}[Memory and self energy]\label{extended:record:086}\label{extended:thm:memory_self_energy}
If \(z\) lies to the right of the spectrum of \(A_{QQ}\), then
\[
\int_0^\infty
\e^{-zt}\Kcal(t)\dd t
=
\Sigma_P(z).
\]
If, for constants \(M\geq0\) and \(\gamma>0\),
\[
\norm{\e^{tA_{QQ}}}
\leq
M\e^{-\gamma t},
\]
then
\[
\norm{\Kcal(t)}
\leq
M\norm{A_{PQ}}\norm{A_{QP}}\e^{-\gamma t}
\]
and
\[
\norm{\Kcal}_{L^1(0,\infty)}
\leq
\frac{M}{\gamma}
\norm{A_{PQ}}\norm{A_{QP}}.
\]
\end{theorem}

\begin{proof}
If \(X_Q=\{0\}\), the memory kernel and self-energy are both zero, so the identity is immediate. Otherwise write \(B=A_{QQ}\), which is a bounded operator on the hidden Banach space, and put \(s(B)=\max\operatorname{Re}\Spec(B)\). Choose a real \(\beta\) with \(s(B)<\beta<\operatorname{Re}z\). Spectral mapping and the spectral-radius formula give \(r(\e^B)=\e^{s(B)}<\e^\beta\), hence \(\norm{\e^{nB}}\leq M_\beta\e^{\beta n}\) for every integer \(n\geq0\), after increasing \(M_\beta\) to cover the finitely many initial values. Write \(t=n+r\), where \(0\leq r<1\). Then
\[
\norm{\e^{tB}}\leq M_\beta\exp(\norm B+\max(0,-\beta))\e^{\beta t}.
\]
Thus \(\e^{-zt}\e^{tB}\) is Bochner integrable and tends to zero. Integrating its derivative gives both inverse identities for \(R=\int_0^\infty\e^{-zt}\e^{tB}\dd t\), so \(R=(zI-B)^{-1}\). This is the bounded-generator growth principle in \cite[Chapter IV, Corollary 3.12]{EngelNagel2000}. Multiplication by the coupling blocks gives the first identity. The remaining estimates follow directly from the assumed positive decay rate.
\end{proof}

\begin{theorem}[Criterion for exact local closure]\label{extended:record:087}\label{extended:thm:local_closure}
Assume \(x_Q(0)=0\) and \(d=\dim X_Q<\infty\). The following conditions are equivalent.
\begin{enumerate}
\item For every retained initial state,
\[
P\e^{tA}P
=
\e^{tA_{PP}}P
\]
for all \(t\geq0\).
\item \(\Sigma_P(z)=0\) for every sufficiently large \(z\).
\item \(\Kcal(t)=0\) for every \(t\geq0\).
\item If \(d=\dim X_Q\), then
\[
A_{PQ}A_{QQ}^kA_{QP}=0
\]
for \(0\leq k<d\).
\end{enumerate}
\end{theorem}

\begin{proof}
By \cref{extended:thm:laplace_packet_resolvent}, the Laplace transform of the projected propagator is the projected resolvent. By \cref{extended:thm:feshbach_resolvent}, the first condition is equivalent to
\[
(zI-A_{PP}-\Sigma_P(z))^{-1}
=
(zI-A_{PP})^{-1}
\]
for all sufficiently large \(z\), which is equivalent to the second condition. The second and third conditions are equivalent by \cref{extended:thm:memory_self_energy} and uniqueness of the Laplace transform.

For large \(z\),
\[
\Sigma_P(z)
=
\sum_{k=0}^{\infty}
\frac{A_{PQ}A_{QQ}^kA_{QP}}{z^{k+1}}.
\]
Thus the second condition implies the moment identities. Conversely, the first \(d\) identities imply all later ones by the Cayley Hamilton theorem, so the series vanishes.
\end{proof}

\begin{remark}\label{extended:record:088}
A finite hidden block produces a strictly proper rational self energy. It tends to zero as \(z\) tends to infinity. Hence a nonzero constant self energy cannot arise from hidden dynamics alone. A constant correction must already be included in the retained generator.
\end{remark}

\section{Memory modes and controlled compression}\label{extended:section:15}\label{extended:sec:memory_compression}

Finite hidden systems produce rational self energies and finite modal memory. This is the state space viewpoint familiar from linear systems theory \cite{Kailath1980}. The same formulas also underlie the Mori Zwanzig representation of reduced dynamics \cite{Zwanzig1961,ChorinHaldKupferman2000}.

\begin{theorem}[Polynomial exponential memory of a finite hidden system]\label{extended:record:089}\label{extended:thm:modal_memory}
Let \(X_Q\) be a finite dimensional complex space. Let \(\mu_1,\ldots,\mu_r\) be the distinct eigenvalues of \(A_{QQ}\), let \(Q_j\) be the corresponding primary spectral projections, and set
\[
N_j=(A_{QQ}-\mu_jI)Q_j.
\]
Let \(\nu_j\) be the least positive integer for which \(N_j^{\nu_j}=0\), and define
\[
B_{j,k}
=
A_{PQ}N_j^kQ_jA_{QP}
\qquad
(0\leq k<\nu_j).
\]
Then
\[
\Kcal(t)
=
\sum_{j=1}^r
\e^{\mu_jt}
\sum_{k=0}^{\nu_j-1}
\frac{t^k}{k!}B_{j,k}
\]
and
\[
\Sigma_P(z)
=
\sum_{j=1}^r
\sum_{k=0}^{\nu_j-1}
\frac{B_{j,k}}{(z-\mu_j)^{k+1}}.
\]
Suppose \(\operatorname{Re}\mu_j<0\) for every \(j\). If \(J\subseteq\{1,\ldots,r\}\) and \(\Kcal_J\) retains the complete primary blocks indexed by \(J\), then
\[
\norm{\Kcal-\Kcal_J}_{L^1(0,\infty)}
\leq
\sum_{j\notin J}
\sum_{k=0}^{\nu_j-1}
\frac{\norm{B_{j,k}}}
{(-\operatorname{Re}\mu_j)^{k+1}}.
\]
\end{theorem}

\begin{proof}
The primary decomposition gives
\[
A_{QQ}Q_j=\mu_jQ_j+N_j
\]
with commuting scalar and nilpotent parts on \(Q_jX_Q\). Therefore
\[
\e^{tA_{QQ}}Q_j
=
\e^{\mu_jt}
\sum_{k=0}^{\nu_j-1}
\frac{t^k}{k!}N_j^kQ_j.
\]
Multiplication by \(A_{PQ}\) and \(A_{QP}\) gives the memory formula. On the \(j\)th primary space, the nilpotent resolvent expansion is
\[
(zI-A_{QQ})^{-1}Q_j
=
\sum_{k=0}^{\nu_j-1}
\frac{N_j^kQ_j}{(z-\mu_j)^{k+1}}.
\]
Multiplication by the coupling blocks gives the self energy formula for every \(z\) outside the hidden spectrum. Equivalently, on a right half plane it is the Laplace transform identity
\[
\int_0^\infty
\e^{-zt}\e^{\mu_jt}
\frac{t^k}{k!}
\dd t
=
\frac{1}{(z-\mu_j)^{k+1}}.
\]
Under stability,
\[
\int_0^\infty
\e^{\operatorname{Re}\mu_jt}
\frac{t^k}{k!}
\dd t
=
\frac{1}{(-\operatorname{Re}\mu_j)^{k+1}}.
\]
Sum the omitted terms.
\end{proof}

\begin{corollary}[Optimal selection for the certified primary block bound]\label{extended:record:090}\label{extended:cor:optimal_primary_selection}
Under the stable assumptions of \cref{extended:thm:modal_memory}, define the primary block weights
\[
w_j
=
\sum_{k=0}^{\nu_j-1}
\frac{\norm{B_{j,k}}}
{(-\operatorname{Re}\mu_j)^{k+1}}.
\]
Among all sets \(J\) containing \(k\) primary blocks, the certified upper bound
\[
\sum_{j\notin J}w_j
\]
is minimized by retaining the \(k\) largest weights. If
\[
w_{(1)}\geq w_{(2)}\geq\cdots\geq w_{(r)},
\]
then the least number of retained primary blocks certified to give \(L^1\) error at most \(\varepsilon\) is the least \(k\) for which
\[
\sum_{j>k}w_{(j)}\leq\varepsilon.
\]
\end{corollary}

\begin{proof}
If a retained set contains a block with smaller weight than an omitted block, exchanging them reduces the omitted weight sum. Repeating the exchange gives the set of the \(k\) largest weights. The tolerance statement is the corresponding sorted tail condition.
\end{proof}

\begin{corollary}[Normal hidden dynamics]\label{extended:record:091}\label{extended:cor:normal_modal_memory}
Suppose \(X_Q\) is a finite dimensional Hilbert space. If \(A_{QQ}\) is normal with spectral resolution
\[
A_{QQ}=\sum_{j=1}^r\mu_jQ_j,
\]
then
\[
\Kcal(t)=\sum_{j=1}^r\e^{\mu_jt}B_j,
\qquad
\Sigma_P(z)=\sum_{j=1}^r\frac{B_j}{z-\mu_j},
\]
where
\[
B_j=A_{PQ}Q_jA_{QP}.
\]
If \(\operatorname{Re}\mu_j<0\), then
\[
\norm{\Kcal-\Kcal_J}_{L^1(0,\infty)}
\leq
\sum_{j\notin J}
\frac{\norm{B_j}}{-\operatorname{Re}\mu_j}.
\]
\end{corollary}

\begin{proof}
Normality removes every nilpotent part, so \(\nu_j=1\). Apply \cref{extended:thm:modal_memory}.
\end{proof}

\begin{corollary}[Pole order and hidden Jordan depth]\label{extended:record:092}\label{extended:cor:pole_jordan_depth}
Every finite hidden block produces a strictly proper rational self energy. At an eigenvalue \(\mu_j\), its pole order is at most the size of the largest Jordan block of \(A_{QQ}\) at \(\mu_j\). Equality holds when the corresponding highest coefficient \(B_{j,\nu_j-1}\) is nonzero.
\end{corollary}

\begin{proof}
Read the pole order from the partial fraction expansion in \cref{extended:thm:modal_memory}.
\end{proof}

\begin{figure}[t]
\centering
\includegraphics[width=.72\textwidth]{jordan_memory.pdf}
\caption{A stable Jordan block produces a polynomial exponential memory mode. A diagonal hidden block produces only pure exponentials.}
\label{extended:fig:jordan_memory}
\end{figure}

\begin{corollary}[Finite auxiliary state realization]\label{extended:record:093}\label{extended:cor:finite_auxiliary_realization}
Every kernel of the form
\[
\Kcal(t)=B\e^{tD}C
\]
can be represented exactly by adjoining the finite auxiliary state
\[
\dot y=Dy+Cx_P,
\qquad
\dot x_P=A_{PP}x_P+By,
\qquad
y(0)=0.
\]
A nonzero auxiliary initial state adds the explicit forcing \(B\e^{tD}y(0)\). The self energy of the augmented system is
\[
B(zI-D)^{-1}C.
\]
\end{corollary}

\begin{proof}
Eliminate \(y\) by variation of constants or apply \cref{extended:thm:feshbach_resolvent} to the augmented block matrix.
\end{proof}

\begin{theorem}[Volterra stability under kernel compression]\label{extended:record:094}\label{extended:thm:volterra_compression}
Let \(S(t)\) be a strongly continuous semigroup on a Banach space \(X\), fix \(T>0\), and assume
\[
\norm{S(t)}\leq M_T
\quad(0\leq t\leq T).
\]
Let \(K,\widetilde K\in L^1((0,T),\mathcal B(X))\), let \(g\in L^1((0,T),X)\), and let \(x,\widetilde x\in C([0,T],X)\). Suppose these are mild solutions with the same initial state and forcing, except that \(K\) is replaced by \(\widetilde K\):
\[
\begin{aligned}
x(t)
={}&S(t)x_0+
\int_0^tS(t-s)g(s)\dd s\\
&+
\int_0^tS(t-s)
\int_0^sK(s-r)x(r)\dd r\dd s.
\end{aligned}
\]
Assume
\[
\sup_{0\leq t\leq T}\norm{\widetilde x(t)}\leq B.
\]
Set
\[
L=\norm K_{L^1(0,T)},
\qquad
\varepsilon=\norm{K-\widetilde K}_{L^1(0,T)}.
\]
Then
\[
\sup_{0\leq t\leq T}
\norm{x(t)-\widetilde x(t)}
\leq
M_TBT\varepsilon
\exp(M_TLT).
\]
\end{theorem}

\begin{proof}
Let
\[
E(t)=\sup_{0\leq r\leq t}\norm{x(r)-\widetilde x(r)}.
\]
Subtract the two mild equations. Fubini's theorem and the triangle inequality give
\[
E(t)
\leq
M_TL\int_0^tE(r)\dd r
+
M_TBt\varepsilon.
\]
For \(t\leq T\), replace the last factor \(t\) by \(T\) and apply Gronwall's inequality.
\end{proof}

\begin{corollary}[Trajectory bound for modal truncation]\label{extended:record:095}\label{extended:cor:modal_trajectory_bound}
Under the assumptions of \cref{extended:thm:modal_memory,extended:thm:volterra_compression}, replacing \(\Kcal\) by \(\Kcal_J\) changes the retained trajectory by at most
\[
M_TBT
\exp(M_TLT)
\sum_{j\notin J}
\sum_{k=0}^{\nu_j-1}
\frac{\norm{B_{j,k}}}
{(-\operatorname{Re}\mu_j)^{k+1}}.
\]
\end{corollary}

\begin{proof}
Insert the \(L^1\) error from \cref{extended:thm:modal_memory} into \cref{extended:thm:volterra_compression}.
\end{proof}

\begin{figure}[t]
\centering
\includegraphics[width=.72\textwidth]{memory_compression.pdf}
\caption{Full and truncated memory kernels with their pointwise difference. The bound uses integrated error.}
\label{extended:fig:memory_compression}
\end{figure}

\begin{remark}\label{extended:record:096}
A finite hidden realization is exact for rational self energy. A general memory kernel may require infinitely many modes or a continuum of modes. The compression theorem gives a stability estimate once an \(L^1\) approximation has been supplied. The largest weight rule is optimal for the certified blockwise bound. It need not minimize the true kernel error when different modal terms cancel.
\end{remark}

\section{Survival, moments, and competing risks}\label{extended:section:16}\label{extended:sec:survival}

Let \(A\) be a transient state subgenerator. Its off diagonal entries are nonnegative and its column sums are nonpositive. Let \(p_0\) be a transient probability vector. Define
\[
p(t)=\e^{tA}p_0,
\qquad
S(t)=\one^Tp(t),
\qquad
a^T=-\one^TA.
\]
The row vector \(a^T\) contains the exit rates.

\begin{theorem}[State dependent killing as survival]\label{extended:record:097}\label{extended:thm:killing_survival}
Let \(Q\) be a finite Markov generator, let \(V\geq0\), and set
\[
A=Q-D_V.
\]
For an initial distribution \(p_0\), define
\[
S_V(t)=\one^T\e^{tA}p_0.
\]
Then
\[
S_V(t)
=
\Exp_{p_0}
\left[
\exp\left(
-\int_0^tV(X_s)\dd s
\right)
\right].
\]
Also,
\[
-S_V'(t)
=
\one^TD_V\e^{tA}p_0.
\]
The matrix \(A=Q-D_V\) has nonnegative off diagonal entries and nonpositive column sums, so it is the subgenerator of the killed, possibly nontransient chain. Thus the potential in the finite Feynman Kac formula is a state dependent killing rate in the forward sub-Markov semigroup. Almost-sure absorption follows under the additional spectral stability hypothesis used below.
\end{theorem}

\begin{proof}
Average \cref{extended:thm:continuous_feynman_kac} with terminal function \(f=\one\) against the initial distribution. Transposition gives
\[
p_0^T\e^{t(Q^T-D_V)}\one
=
\one^T\e^{t(Q-D_V)}p_0.
\]
Since \(\one^TQ=0\), differentiation gives
\[
S_V'(t)
=
-\one^TD_V\e^{tA}p_0.
\]
\end{proof}

\begin{theorem}[Survival and hazard identities]\label{extended:record:098}\label{extended:thm:survival_hazard}
The event density and hazard are
\[
f(t)=a^T\e^{tA}p_0=-S'(t),
\qquad
h(t)=\frac{f(t)}{S(t)}
\]
whenever \(S(t)>0\). The survival function is nonincreasing and
\[
S(t)
=
\exp\left(
-\int_0^th(s)\dd s
\right).
\]
\end{theorem}

\begin{proof}
Because \(A\) is a Metzler matrix, \(\e^{tA}\) is nonnegative. The column sum assumption makes \(a^T\) nonnegative. Differentiate \(S(t)\):
\[
S'(t)=\one^TA\e^{tA}p_0=-a^T\e^{tA}p_0.
\]
The remaining identity follows by integrating \((\log S)'=-h\).
\end{proof}

\begin{theorem}[Laplace transform and moments]\label{extended:record:099}\label{extended:thm:survival_moments}
Assume the spectrum of \(A\) lies in the open left half plane. Then absorption occurs with probability one. For \(s\geq0\),
\[
\Exp[\e^{-s\tau}]
=
a^T(sI-A)^{-1}p_0.
\] For every integer \(m\geq1\),
\[
\Exp[\tau^m]
=
m!\one^T(-A)^{-m}p_0.
\]
\end{theorem}

\begin{proof}
The Laplace transform follows by integrating the density in \cref{extended:thm:survival_hazard}. Stability of \(A\) gives
\[
\int_0^\infty
\e^{-st}\e^{tA}\dd t
=
(sI-A)^{-1}.
\]
At \(s=0\),
\[
a^T(-A)^{-1}p_0
=
\one^Tp_0=1.
\]
For the moments,
\[
\int_0^\infty t^m\e^{tA}\dd t
=
m!(-A)^{-m-1}.
\]
Multiplying by \(a^T=\one^T(-A)\) gives the result. These are the standard matrix formulas for phase type laws \cite{Neuts1981,AsmussenNermanOlsson1996}.
\end{proof}

\begin{proposition}[Residual life as a restarted semigroup]\label{extended:record:100}\label{extended:prop:residual_life}
For \(S(t)>0\), define the conditional transient distribution
\[
p_t
=
\frac{\e^{tA}p_0}{S(t)}.
\]
Then
\[
\frac{S(t+s)}{S(t)}
=
\one^T\e^{sA}p_t.
\]
\end{proposition}

\begin{proof}
Use the semigroup law in the numerator \(S(t+s)=\one^T\e^{sA}\e^{tA}p_0\).
\end{proof}

\begin{theorem}[Competing risk packets]\label{extended:record:101}\label{extended:thm:competing_risks}
Suppose the exit row splits into nonnegative cause rows
\[
a^T
=
\sum_{r=1}^q a_r^T.
\]
Define
\[
f_r(t)=a_r^T\e^{tA}p_0.
\]
Then
\[
f(t)=\sum_{r=1}^qf_r(t).
\]
The cause specific hazard is
\[
h_r(t)=\frac{f_r(t)}{S(t)},
\]
and
\[
h(t)=\sum_{r=1}^qh_r(t).
\]
If the spectrum of \(A\) lies in the open left half-plane, the eventual probability of cause \(r\) is
\[
\int_0^\infty f_r(t)\dd t
=
a_r^T(-A)^{-1}p_0.
\]
\end{theorem}

\begin{proof}
Every identity follows by linearity and the stable resolvent integral.
\end{proof}

\begin{figure}[t]
\centering
\includegraphics[width=.72\textwidth]{survival_causes.pdf}
\caption{Survival and two cause specific event densities from one transient semigroup.}
\label{extended:fig:survival_causes}
\end{figure}

\begin{remark}\label{extended:record:102}
A finite phase type hazard need not be monotone. The semigroup formulas guarantee nonnegative density and decreasing survival, not monotonicity of the hazard.
\end{remark}

\section{Finite illustrations and computational checks}\label{extended:section:17}\label{extended:sec:finite_checks}

The figures in the article use small matrices chosen to make each mechanism visible. They are not evidence for claims outside the proved finite setting. The accompanying verification script recomputes the following identities with independently generated data.

\begin{enumerate}
\item Packet block convolution, the block differential system, and the Laplace resolvent formula.
\item The Duhamel and resolvent defect identities.
\item Exact and perturbed Markov lumping.
\item Conditional expectation, tower laws, martingale packet energy, and Hermite multipliers.
\item Bayesian reconstruction, sufficiency, local distortion, and information decomposition.
\item The evidence variance identity along a Bayesian flow.
\item The Balakrishnan formula and the fractional defect estimate.
\item Nodal Green reconstruction and the Poisson error bound.
\item Discrete Feynman Kac path enumeration.
\item The Schur Feshbach resolvent and memory transform.
\item Jordan memory modes, modal truncation, and the Volterra stability estimate.
\item Feynman Kac killing, survival transforms, moments, and competing risk masses.
\end{enumerate}

\begin{table}[t]
\centering
\begin{tabular}{lc}
\toprule
Identity & Typical residual \\
\midrule
Packet block convolution & below \(3\times10^{-14}\) \\
Duhamel defect & below \(3\times10^{-12}\) \\
Markov lumping & below \(3\times10^{-14}\) \\
Bayesian reconstruction & below \(3\times10^{-15}\) \\
Fractional spectral law & below \(3\times10^{-12}\) \\
Nodal Green identity & below \(3\times10^{-13}\) \\
Feynman Kac enumeration & below \(3\times10^{-14}\) \\
Schur Feshbach resolvent & below \(3\times10^{-12}\) \\
Memory Laplace identity & below \(3\times10^{-11}\) \\
Survival moment formulas & below \(3\times10^{-11}\) \\
\bottomrule
\end{tabular}
\caption{Numerical tolerances in the finite verification suite. The proofs do not depend on these computations.}
\label{extended:tab:verification_tolerances}
\end{table}

\section{Limits and further questions}\label{extended:section:18}\label{extended:sec:limits}

The bounded inverse limit theorem isolates one clean route to continuum evolution. Unbounded generators require graph norms, domain compatibility, and a generation theorem. The range leakage theorem converts the closure quantity into an off diagonal generator estimate for bounded orthogonal refinement models. Concrete discretizations still require a uniform estimate for that off diagonal block.

The finite Bayesian sufficiency theorem is exact, but it concerns one likelihood and one fixed partition. The relative entropy chain rule measures both retained and discarded information for that update. Statistical sufficiency for a parametric family requires a common packet structure that preserves the full parameter dependence. In continuous models the corresponding object is a sufficient sigma algebra or statistic.

The fractional defect theorem uses a common spectral gap. Removing that gap requires weaker norms or a Hölder estimate near the null space. For sectorial operators, the same calculation should be carried out through resolvent bounds on a sector, with domains tracked explicitly. The finite spectral statement does not settle those issues.

The Green and Poisson results exploit a one dimensional second difference matrix. Higher dimensions, irregular meshes, and variable coefficients require a separate stability theory. The packet language remains useful only after coercivity, consistency, and boundary approximation have been established.

Memory compression is exact once a finite hidden realization is known and quantitative once an integrable kernel approximation is supplied. The primary decomposition shows that a finite hidden block always produces polynomial exponential modes, including the effects of Jordan chains. Choosing a short approximation remains a different problem. Strongly nonnormal systems can exhibit transient growth and may require balanced or rational approximation methods. General Volterra kernels need not admit a short finite realization.

Finally, Markov lumping, Bayesian sufficiency, and memory closure answer different questions. Lumpability asks whether coarse probabilities evolve autonomously. Bayesian sufficiency asks whether a fixed coarse likelihood reproduces the posterior packet masses without hiding a change in the conditional distribution inside a packet. Memory closure asks whether discarded dynamics can be represented without history. Similar notation should not obscure these distinctions.

\section{Conclusion}\label{extended:section:19}\label{extended:sec:conclusion}

A bounded semigroup becomes a matrix of packet transfers once a finite packet system is fixed. The matrix obeys convolution in time, a closed differential system, a Laplace resolvent formula, and a directed path expansion. A generator defect propagates through Duhamel's formula and through the resolvent calculus. Exact compatibility gives a semigroup on the bounded inverse limit. Summable adjacent defects give convergent reconstructed evolutions. Range leakage, and in particular an off diagonal generator block, determines whether the limit retains the semigroup law.

In probability, the same structure appears in several forms. Markov transfer is linear and exact lumpability is the vanishing of a canonical kernel obstruction. Conditional expectation is an orthogonal stochastic projector whose differences form martingale packets. The Ornstein Uhlenbeck semigroup resolves exactly into Hermite packets. Bayesian transfer is nonlinear because of evidence normalization. Its exact two stage decomposition separates coarse packet evidence from local conditional change. Relative entropy then records exactly how much information lies at each level. A fixed sufficient coarse description exists for every prior precisely when the likelihood is constant inside each packet. The common level sets of a likelihood family give its unique coarsest fixed sufficient partition.

Positive fractional powers preserve exact intertwining. Under a common spectral gap, the Balakrishnan formula turns an ordinary generator defect into an explicit fractional defect. Green operators and finite differences show how source reconstruction fits the same transfer language. Feynman Kac weights Markov paths by a potential without leaving the semigroup framework.

Hidden packets provide the final structural lesson. Exact elimination does not erase them. It returns them as a rational self energy in the resolvent and a causal memory kernel in time. The primary decomposition of the hidden generator produces every finite memory kernel as a finite sum of polynomial exponential modes. Stable blocks can be truncated with an integrated error certificate, and the Volterra estimate transfers that certificate to the retained trajectory. State dependent killing identifies the Feynman Kac potential with a survival mechanism. Survival and competing risks are then scalar reconstructions of a transient semigroup.

The article therefore separates four operations that are often compressed into one informal idea of coarse dynamics:
\[
\begin{aligned}
\text{intertwining}
&\longrightarrow
\text{refinement control}\\
&\longrightarrow
\text{statistical sufficiency}\\
&\longrightarrow
\text{memory closure}.
\end{aligned}
\]
Each arrow has its own hypotheses and its own obstruction.





\clearpage
\part*{Part III\quad Retained results and operator calculus}\addcontentsline{toc}{part}{Part III: Retained results and operator calculus}
\setcounter{section}{0}\setcounter{figure}{0}\setcounter{table}{0}
\renewcommand{\thesection}{R\arabic{section}}
\renewcommand{\theHsection}{R.\arabic{section}}
\renewcommand{\thefigure}{R\arabic{figure}}
\renewcommand{\theHfigure}{R.\arabic{figure}}
\renewcommand{\thetable}{R\arabic{table}}
\renewcommand{\theHtable}{R.\arabic{table}}
\section{Transfer paths and product-topology limits}\label{earlier:section:01}
Let $T(t)$ be an arbitrary strongly continuous semigroup on a complex Banach space $X$, with generator $A$, and suppose $\norm{T(t)}\leq M e^{\omega t}$ for $M\geq1$ and $\omega\in\mathbb R$. Let $(P_\alpha)_{\alpha\in I_0}$ be a finite family of bounded pairwise orthogonal idempotents summing to $I$, and put $T_{\beta\alpha}(t)=P_\beta T(t)P_\alpha$. The generator need not be bounded. In the following resolvent statement the integral is strong: apply the integrand to each $x\in X$ and take its Bochner integral in $X$. It is not asserted to be a Bochner integral in the operator-norm space $\mathcal B(X)$. The general Laplace formula is \cite[Chapter II, Theorem 1.10]{EngelNagel2000}.
\begin{theorem}[Packet convolution]\label{earlier:record:001}\label{earlier:thm:packet_convolution}
For all \(s,t\geq0\) and all packets \(\alpha,\beta\),
\[
\boxed{
T_{\beta\alpha}(s+t)
=
\sum_{\gamma\in I_0}
T_{\beta\gamma}(s)T_{\gamma\alpha}(t).
}
\]
The original semigroup is reconstructed from its blocks by
\[
T(t)=\sum_{\alpha,\beta\in I_0}T_{\beta\alpha}(t).
\]
\end{theorem}

\begin{proof}
Insert the identity between the two semigroup factors:
\[
\begin{aligned}
P_\beta T(s+t)P_\alpha
&=P_\beta T(s)\left(\sum_\gamma P_\gamma\right)T(t)P_\alpha\\
&=\sum_\gamma
P_\beta T(s)P_\gamma T(t)P_\alpha.
\end{aligned}
\]
The reconstruction formula follows by inserting the packet identity on both sides of \(T(t)\).
\end{proof}

\begin{proposition}[Packet resolvent convolution]\label{earlier:record:002}\label{earlier:prop:packet_resolvent}
Let \(z\) belong to the resolvent set of \(A\). If \(\operatorname{Re}z\) is larger than a growth bound for the semigroup, then
\[
P_\beta(zI-A)^{-1}P_\alpha
=
\int_0^{\infty}\e^{-zt}T_{\beta\alpha}(t)\dd t.
\]
\end{proposition}

\begin{proof}
The standard Laplace representation of the resolvent is
\[
(zI-A)^{-1}
=
\int_0^{\infty}\e^{-zt}T(t)\dd t.
\]
Multiplication by the bounded packet projections may be passed through the Bochner integral.
\end{proof}


For the path-order statement, now assume $A$ is bounded as in Part II. A directed edge from $\alpha$ to $\beta$ means that $P_\beta A P_\alpha$ is nonzero, with $\alpha\ne\beta$; a packet always has the length-zero path to itself.
\begin{corollary}[Minimum transfer order]\label{earlier:record:003}\label{earlier:cor:minimum_transfer_order}
Suppose every directed packet path from \(\alpha\) to \(\beta\) uses at least \(m_0\) off diagonal transfers. Then
\[
P_\beta\e^{tA}P_\alpha
=O(t^{m_0})
\quad\text{as }t\downarrow0.
\]
If there is no directed path from \(\alpha\) to \(\beta\), the block vanishes for all \(t\).
\end{corollary}

\begin{proof}
Every nonzero term in \cref{extended:thm:dyson_paths} contains at least \(m_0\) time integrations. The simplex \(0<r_1<\cdots<r_m<t\) has volume \(t^m/m!\), so the first possible order is \(t^{m_0}\). If there is no path, every packet term vanishes.
\end{proof}


Let $C_j:X_{j+1}\to X_j$ be bounded linear maps, and put $C_{\ell,m}=C_\ell\cdots C_{m-1}$ for $\ell<m$, with $C_{j,j}=I$. No common ambient space is assumed.
\begin{theorem}[Several level telescoping identity]\label{earlier:record:004}\label{earlier:thm:multilevel_telescoping}
Let \(T_j\) be bounded operators on \(X_j\), and set
\[
E_j=C_jT_{j+1}-T_jC_j.
\]
Then
\[
\boxed{
C_{\ell,m}T_m-T_\ell C_{\ell,m}
=
\sum_{j=\ell}^{m-1}
C_{\ell,j}E_jC_{j+1,m}.
}
\]
\end{theorem}

\begin{proof}
Expand each summand as
\[
C_{\ell,j+1}T_{j+1}C_{j+1,m}
-
C_{\ell,j}T_jC_{j,m}.
\]
The intermediate terms cancel. The first surviving term is \(C_{\ell,m}T_m\), and the last is \(-T_\ell C_{\ell,m}\).
\end{proof}

\begin{corollary}[Summable generator defects]\label{earlier:record:005}\label{earlier:cor:summable_defects}
Suppose \(C_j\) are contractions and \((T_j(t))\) are contraction semigroups with bounded generators \(A_j\). Set
\[
\delta_j=\norm{C_jA_{j+1}-A_jC_j}.
\]
Then
\[
\norm{C_{\ell,m}T_m(t)-T_\ell(t)C_{\ell,m}}
\leq
 t\sum_{j=\ell}^{m-1}\delta_j.
\]
If \(\sum_{j\geq0}\delta_j<\infty\), the right side tends to zero as \(\ell\to\infty\), uniformly in \(m>\ell\) and uniformly for \(t\) in a fixed bounded interval.
\end{corollary}

\begin{proof}
By \cref{extended:thm:duhamel_defect}, the interlevel semigroup defect has norm at most \(t\delta_j\). Apply \cref{earlier:thm:multilevel_telescoping} and use the contraction bounds on all coarsening maps.
\end{proof}

\begin{definition}[Projective limit space]\label{earlier:record:006}\label{earlier:def:projective_limit}
Let \((X_j,C_j)\) be a sequence of Banach spaces and bounded linear maps \(C_j:X_{j+1}\to X_j\). Its projective limit is
\[
X_\infty
=
\left\{
(x_j)_{j\geq0}\in\prod_{j\geq0}X_j:
C_jx_{j+1}=x_j\text{ for every }j
\right\}.
\]
We equip \(X_\infty\) with the subspace topology inherited from the product topology.
\end{definition}

\begin{theorem}[Exact projective limit semigroup]\label{earlier:record:007}\label{earlier:thm:projective_limit_semigroup}
Suppose each \(X_j\) carries a strongly continuous semigroup \((T_j(t))_{t\geq0}\), and
\[
C_jT_{j+1}(t)=T_j(t)C_j
\]
for every \(j\) and every \(t\geq0\). Then
\[
T_\infty(t)(x_j)_{j\geq0}
=
(T_j(t)x_j)_{j\geq0}
\]
defines a strongly continuous semigroup on \(X_\infty\).
\end{theorem}

\begin{proof}
Exact intertwining gives
\[
C_jT_{j+1}(t)x_{j+1}
=
T_j(t)C_jx_{j+1}
=
T_j(t)x_j.
\]
Thus the coordinatewise evolution remains in the projective limit. The semigroup law holds in every coordinate. Strong continuity in the product topology follows because every coordinate map
\[
t\longmapsto T_j(t)x_j
\]
is continuous.
\end{proof}

\begin{corollary}[Coordinate generators]\label{earlier:record:008}\label{earlier:cor:projective_generators}
Assume the setting of \cref{earlier:thm:projective_limit_semigroup}. If \(x=(x_j)\) satisfies \(x_j\in\operatorname{Dom}(A_j)\) and the sequence \((A_jx_j)\) belongs to \(X_\infty\), then \(x\) lies in the generator domain of \(T_\infty\) and
\[
A_\infty x=(A_jx_j)_{j\geq0}.
\]
\end{corollary}

\begin{proof}
The difference quotients converge in every coordinate to \(A_jx_j\). The assumed compatibility places the coordinatewise limit in \(X_\infty\), which is exactly convergence in the product topology.
\end{proof}


The product topology used here differs from the bounded inverse-limit norm of Part II. No uniform boundedness of coordinate generators is required by the coordinatewise statement.
\section{Reversible coarse dynamics}\label{earlier:section:02}
Let $Q$ be a finite Markov generator acting on column probability vectors, let $(B_a)$ be a partition into nonempty blocks, and let $C_{ai}=\one_{i\in B_a}$. Strong lumpability means $CQ=\overline Q C$ for a coarse Markov generator. For observables, the backward generator is $Q^T$; weighted $L^2$ inner products use the stated stationary probability.
\begin{theorem}[Reversibility descends]\label{earlier:record:009}\label{earlier:thm:reversibility_descends}
Assume \(Q\) is strongly lumpable and reversible with respect to a positive stationary distribution \(\pi\), so
\[
Q_{ij}\pi_j=Q_{ji}\pi_i.
\]
Define
\[
\overline\pi_a=\sum_{i\in B_a}\pi_i.
\]
Then \(\overline\pi\) is stationary for \(\overline Q\), and
\[
\overline Q_{ab}\overline\pi_b
=
\overline Q_{ba}\overline\pi_a.
\]
Thus the coarse chain is reversible.
\end{theorem}

\begin{proof}
Stationarity follows from
\[
\overline Q\overline\pi
=\overline QC\pi
=CQ\pi=0.
\]
For reversibility, use the aggregate rate formula:
\[
\begin{aligned}
\overline Q_{ab}\overline\pi_b
&=\sum_{j\in B_b}\pi_j\sum_{i\in B_a}Q_{ij}\\
&=\sum_{i\in B_a}\sum_{j\in B_b}\pi_iQ_{ji}\\
&=\overline Q_{ba}\overline\pi_a.
\end{aligned}
\]
\end{proof}

\begin{corollary}[Coarse Dirichlet form]\label{earlier:record:010}\label{earlier:cor:coarse_dirichlet_form}
Under the hypotheses of \cref{earlier:thm:reversibility_descends}, let \(Jg\) be the fine function that is constant with value \(g_a\) on \(B_a\). Then
\[
-\ip{Jg}{Q^TJg}_{L^2(\pi)}
=
-\ip{g}{\overline Q^Tg}_{L^2(\overline\pi)}.
\]
\end{corollary}

\begin{proof}
Expand the left side and group the fine indices by their coarse blocks. The aggregate rate criterion converts the resulting sums into the coarse generator entries.
\end{proof}


\section{Prior-restricted sufficiency and filtering}\label{earlier:section:03}
Let $L$ be a strictly positive finite likelihood and let $\Bcal_L(p)=L\odot p/\ip{L}{p}$ for a probability vector $p$. Fix probability vectors $r^{(a)}$ supported on $B_a$, and restrict the priors to $p_i=q_a r_i^{(a)}$ for $i\in B_a$, with $q$ a coarse probability vector.
\begin{proposition}[Sufficiency relative to a prior disintegration]\label{earlier:record:011}\label{earlier:prop:relative_sufficiency}
Define
\[
\overline L_a
=
\sum_{i\in B_a}L_ir_i^{(a)}.
\]
For every prior in the declared family,
\[
C\Bcal_L(p)
=
\Bcal_{\overline L}(q).
\]
\end{proposition}

\begin{proof}
The evidence in block \(B_a\) is
\[
Z_a
=q_a\sum_{i\in B_a}L_ir_i^{(a)}
=q_a\overline L_a.
\]
Apply \cref{extended:thm:bayesian_packet_reconstruction}.
\end{proof}

\begin{lemma}[Normalization bound]\label{earlier:record:012}\label{earlier:lem:normalization_bound}
Let \(a,b\in[0,\infty)^m\) be nonzero, and set \(A=\norm a_1\) and \(B=\norm b_1\). Then
\[
\boxed{
\norm{a/A-b/B}_1
\leq
\frac{2\norm{a-b}_1}{\max\{A,B\}}.
}
\]
\end{lemma}

\begin{proof}
Using \(A\) in the denominator gives
\[
\norm{a/A-b/B}_1
\leq
\frac{\norm{a-b}_1}{A}
+
\norm b_1\abs{1/A-1/B}.
\]
The second term equals \(\abs{B-A}/A\), which is at most \(\norm{a-b}_1/A\). Interchanging \(a\) and \(b\) gives the analogous bound with \(B\). Taking the smaller of the two bounds gives the result.
\end{proof}


For an arbitrary prior $p$, put $q=Cp$ and $Z(p,L)=\ip{L}{p}$. Define $\ell_a(p)=\sum_{i\in B_a}L_ip_i/q_a$ when $q_a>0$; choose any positive value when $q_a=0$. Only the product $q_a\ell_a(p)$ enters the formulas.
\begin{theorem}[Exact information loss formula]\label{earlier:record:013}\label{earlier:thm:information_loss}
For every prior,
\[
\boxed{
C\Bcal_L(p)
=
\Bcal_{\ell(p)}(q).
}
\]
Let \(\overline L\) be any fixed positive coarse likelihood and set
\[
d(p\mid L,\overline L)
=
\sum_{a=1}^mq_a\abs{\ell_a(p)-\overline L_a}.
\]
Then
\[
\norm{
C\Bcal_L(p)-\Bcal_{\overline L}(q)
}_1
\leq
\frac{2d(p\mid L,\overline L)}
{\max\{Z(p,L),\ip{\overline L}{q}\}}.
\]
\end{theorem}

\begin{proof}
The packet evidence in block \(a\) is \(q_a\ell_a(p)\), which proves the exact identity. Apply \cref{earlier:lem:normalization_bound} to
\[
a=q\odot\ell(p),
\qquad
b=q\odot\overline L.
\]
Their \(\ell^1\) norms are the two evidences, and \(\norm{a-b}_1=d(p\mid L,\overline L)\).
\end{proof}

\begin{corollary}[Oscillation bound]\label{earlier:record:014}\label{earlier:cor:bayes_oscillation}
Suppose \(\overline L_a\) is any representative likelihood for block \(B_a\), and set
\[
\omega_a
=
\max_{i\in B_a}\abs{L_i-\overline L_a}.
\]
Then
\[
d(p\mid L,\overline L)
\leq
\sum_aq_a\omega_a.
\]
Thus small within block likelihood oscillation gives a small coarse posterior defect whenever the evidence is bounded away from zero.
\end{corollary}

\begin{proof}
The conditional average \(\ell_a(p)\) lies in the convex hull of the likelihood values inside \(B_a\). Hence
\[
\abs{\ell_a(p)-\overline L_a}\leq\omega_a.
\]
Multiply by \(q_a\) and sum.
\end{proof}


Along the Bayesian flow $p_i(t)=e^{tg_i}p_i(0)/\sum_j e^{tg_j}p_j(0)$, the following statement holds for every real observable $h$.
\begin{corollary}[Covariance identity]\label{earlier:record:015}\label{earlier:cor:bayesian_covariance}
For any observable \(h\in\RR^n\),
\[
\frac{\dd}{\dd t}\EE_{p(t)}[h]
=
\Cov_{p(t)}(h,g).
\]
In particular,
\[
\frac{\dd}{\dd t}\EE_{p(t)}[g]
=
\Var_{p(t)}(g)\geq0.
\]
\end{corollary}

\begin{proof}
Differentiate \(\sum_ih_ip_i(t)\) and insert the replicator equation.
\end{proof}


For stochastic matrices $K,\overline K$ and positive likelihoods $L_y,\overline L_y$, put $\Phi_y=\Bcal_{L_y}\circ K$ and $\overline\Phi_y=\Bcal_{\overline L_y}\circ\overline K$.
\begin{theorem}[Exact coarse hidden Markov filter]\label{earlier:record:016}\label{earlier:thm:exact_coarse_filter}
Assume
\[
CK=\overline KC
\]
and every likelihood \(L_y\) is constant on each coarse block with block values \(\overline L_y\). Then
\[
C\Phi_y(p)=\overline\Phi_y(Cp)
\]
for every prior. For any observation sequence \(y_1,\ldots,y_r\), the complete filtering recursion commutes with aggregation:
\[
C\Phi_{y_r}\circ\cdots\circ\Phi_{y_1}
=
\overline\Phi_{y_r}\circ\cdots\circ\overline\Phi_{y_1}C.
\]
\end{theorem}

\begin{proof}
Exact Markov lumping gives \(CKp=\overline KCp\). The Bayesian sufficiency criterion gives exact commutation of the update. Their composition proves the one step identity. Iterate it over the observation sequence.
\end{proof}

\begin{theorem}[One step filter defect]\label{earlier:record:017}\label{earlier:thm:filter_defect}
Let \(r=Kp\), \(q=Cr\), and let \(\ell_y(r)\) be the exact coarse likelihood obtained by averaging \(L_y\) inside each block with respect to \(r\). Assume
\[
0<\overline m\leq(\overline L_y)_a\leq\overline M.
\]
Then
\[
\begin{aligned}
\norm{C\Phi_y(p)-\overline\Phi_y(Cp)}_1
\leq{}&
\frac{2d(r\mid L_y,\overline L_y)}
{\max\{\ip{L_y}{r},\ip{\overline L_y}{q}\}}\\
&+
\frac{2\overline M}{\overline m}
\norm{CKp-\overline KCp}_1.
\end{aligned}
\]
\end{theorem}

\begin{proof}
Insert the intermediate posterior \(\Bcal_{\overline L_y}(CKp)\). The first difference is bounded by \cref{earlier:thm:information_loss}. The second is bounded by \cref{extended:thm:prior_stability} applied to the coarse likelihood.
\end{proof}


\section{Fractional modes and killed transfer}\label{earlier:section:04}
For the one-dimensional Dirichlet difference operator, let $h=1/M$ with an integer $M\geq2$, let $1\leq k\leq M-1$, and set $\lambda_{k,h}=4h^{-2}\sin^2(k\pi h/2)$. The retained explicit coefficient below is valid for every discrete mode.
\begin{corollary}[Fractional low mode dispersion]\label{earlier:record:018}\label{earlier:cor:fractional_low_mode_dispersion}
Let \(0<\alpha\leq1\). Then
\[
0
\leq
(k\pi)^{2\alpha}-\lambda_{k,h}^\alpha
\leq
\frac{\alpha(4k^2)^{\alpha-1}(k\pi)^4}{12}h^2.
\]
For every \(t\geq0\),
\[
\abs{
\e^{-t\lambda_{k,h}^\alpha}
-
\e^{-t(k\pi)^{2\alpha}}
}
\leq
\frac{t\alpha(4k^2)^{\alpha-1}(k\pi)^4}{12}h^2.
\]
\end{corollary}

\begin{proof}
The chord bound \(\sin x\geq2x/\pi\) on \([0,\pi/2]\) gives \(\lambda_{k,h}\geq4k^2\). Since \(x\mapsto x^\alpha\) is increasing and has decreasing derivative,
\[
(k\pi)^{2\alpha}-\lambda_{k,h}^\alpha
\leq
\alpha\lambda_{k,h}^{\alpha-1}
\bigl((k\pi)^2-\lambda_{k,h}\bigr).
\]
Use \cref{extended:thm:low_mode} and the lower bound for \(\lambda_{k,h}\). The exponential estimate follows from
\[
\abs{\e^{-ta}-\e^{-tb}}
\leq t\abs{a-b}
\]
for nonnegative \(a\) and \(b\).
\end{proof}



Let $Q$ and $\overline Q$ be finite Markov generators in the column convention. All times in this section satisfy $t\geq0$.
\begin{theorem}[Potential perturbation bound]\label{earlier:record:019}\label{earlier:thm:feynman_kac_potential_bound}
Let \(V\) and \(W\) be nonnegative diagonal matrices. In the supremum norm,
\[
\boxed{
\norm{
\e^{t(Q^T-V)}
-
\e^{t(Q^T-W)}
}_{\infty\to\infty}
\leq
 t\norm{V-W}_{\infty\to\infty}.
}
\]
\end{theorem}

\begin{proof}
Both killed Markov semigroups are contractions in the supremum norm. Apply the Duhamel identity with generator defect \(W-V\).
\end{proof}

\begin{theorem}[Coarse Feynman Kac transfer]\label{earlier:record:020}\label{earlier:thm:coarse_feynman_kac}
Assume \(CQ=\overline QC\), and let \(J=C^T\) lift a coarse function to the corresponding block constant fine function. Suppose the potential is block constant, so
\[
VJ=J\overline V.
\]
Then
\[
\boxed{
\e^{t(Q^T-V)}J
=
J\e^{t(\overline Q^T-\overline V)}
}
\]
for every \(t\geq0\).
\end{theorem}

\begin{proof}
Transpose the Markov intertwining relation to obtain
\[
Q^TJ=J\overline Q^T.
\]
Together with \(VJ=J\overline V\), this gives
\[
(Q^T-V)J=J(\overline Q^T-\overline V).
\]
The exponential intertwining follows from the power series.
\end{proof}


\section{Memory compression, rational realization, and phase-type consequences}\label{earlier:section:05}
Fix $T>0$. For the next estimate let $A$ be a bounded operator on a Banach space $X$, let $\Kcal,\widetilde\Kcal\in L^1((0,T),\mathcal B(X))$, and let $f$ be continuous. Take continuously differentiable solutions of the displayed equations. This direct bounded-generator estimate retains a sharper exponential constant than the general mild-semigroup estimate in Part II.
\begin{theorem}[Stability of memory compression]\label{earlier:record:021}\label{earlier:thm:memory_compression_stability}
Let \(x\) and \(\widetilde x\) solve
\[
\dot x(t)
=Ax(t)+\int_0^t\Kcal(t-s)x(s)\dd s+f(t)
\]
and
\[
\dot{\widetilde x}(t)
=A\widetilde x(t)
+\int_0^t\widetilde\Kcal(t-s)\widetilde x(s)\dd s
+f(t)
\]
with the same initial value. Fix \(T>0\), and assume
\[
\sup_{0\leq t\leq T}\norm{\widetilde x(t)}\leq M.
\]
Set
\[
a=\norm A,
\qquad
k=\norm{\Kcal}_{L^1(0,T)},
\qquad
\delta=\norm{\Kcal-\widetilde\Kcal}_{L^1(0,T)}.
\]
Then
\[
\boxed{
\sup_{0\leq t\leq T}
\norm{x(t)-\widetilde x(t)}
\leq
MT\delta\e^{(a+k)T}.
}
\]
\end{theorem}

\begin{proof}
Let \(e=x-\widetilde x\). Integrating the difference equation and changing the order of integration gives
\[
\norm{e(t)}
\leq
(a+k)\int_0^t\norm{e(r)}\dd r
+Mt\delta.
\]
Since \(t\leq T\), Gronwall's inequality gives
\[
\norm{e(t)}
\leq
MT\delta\e^{(a+k)t}.
\]
Take the supremum over \([0,T]\).
\end{proof}



For the exact closure criterion of Part II, allowing an arbitrary bounded local generator $B$ does not enlarge the class: differentiating $P e^{tA}P=e^{tB}P$ at $t=0$ forces $B=A_{PP}$ on the retained space. Thus its zero-memory criterion excludes every alternative bounded local generator.

A finite auxiliary realization of a matrix kernel means $\Kcal(t)=B e^{tD}C$, with finite matrices of compatible sizes. Laplace transforms are taken in a right half-plane of absolute convergence. The rational characterization concerns continuous kernels of exponential order, so uniqueness of the ordinary Laplace transform applies.
\begin{theorem}[Rational memory realization]\label{earlier:record:022}\label{earlier:thm:rational_memory_realization}
A matrix valued memory transform has a finite auxiliary realization if and only if it is a strictly proper rational matrix function. More precisely,
\[
\boxed{
\widehat\Kcal(z)=B(zI-D)^{-1}C
}
\]
if and only if
\[
\boxed{
\Kcal(t)=B\e^{tD}C.
}
\]
\end{theorem}

\begin{proof}
The forward implication follows by taking the Laplace transform. The resulting matrix is rational and tends to zero as \(z\) tends to infinity.

Conversely, decompose every entry of a strictly proper rational matrix into partial fractions. A term
\[
\frac{c}{(z-\lambda)^r}
\]
is an entry of the resolvent of an \(r\) by \(r\) Jordan block with eigenvalue \(\lambda\). Taking a direct sum of these blocks and choosing the input and output maps entry by entry gives matrices \(B,D,C\) with the required resolvent. Inverse Laplace transformation gives the time domain formula.
\end{proof}

\begin{corollary}[Exponential sum memory]\label{earlier:record:023}\label{earlier:cor:exponential_sum_memory}
If
\[
\Kcal(t)=\sum_{j=1}^{m}B_j\e^{\lambda_jt}C_j,
\]
then \(\Kcal\) has a finite auxiliary realization with one auxiliary block for each exponential. Repeated poles produce polynomial factors in \(t\) and are realized by Jordan blocks.
\end{corollary}

\begin{proof}
If \(B_j\) and \(C_j\) have an intermediate space of dimension \(d_j\), use a block diagonal matrix \(D\) with blocks \(\lambda_j I_{d_j}\), or the required Jordan blocks when poles repeat. Assemble the matrices \(B_j\) and \(C_j\) into the output and input maps.
\end{proof}


For the survival consequences, $A$ is a finite Metzler matrix with nonpositive column sums and spectrum in the open left half-plane, $p_0$ is a probability vector, $S(t)=\one^T e^{tA}p_0$, and $f(t)=-S'(t)$. The one-phase model has $\lambda>0$. The survival Laplace identity of Part II also holds for every complex $s$ with $\operatorname{Re}s\geq0$: stability gives absolute convergence of $\int_0^\infty e^{-st}e^{tA}\,dt=(sI-A)^{-1}$, and multiplying by the exit row and $p_0$ gives the same transform.
\begin{corollary}[Phase type rationality]\label{earlier:record:024}\label{earlier:cor:phase_type_rationality}
The Laplace transform of every finite phase type distribution is rational. If \(A\) is diagonalizable, then \(S(t)\) and \(f(t)\) are finite linear combinations of exponential functions. Jordan blocks add polynomial factors in \(t\).
\end{corollary}

\begin{proof}
The resolvent \((sI-A)^{-1}\) is rational. The time domain statement follows from the Jordan form of \(A\).
\end{proof}

\begin{corollary}[One phase model]\label{earlier:record:025}\label{earlier:cor:one_phase_model}
If \(A=(-\lambda)\) and \(p_0=1\), then
\[
S(t)=\e^{-\lambda t},
\qquad
f(t)=\lambda\e^{-\lambda t},
\qquad
h(t)=\lambda.
\]
\end{corollary}

\begin{proof}
Insert the scalar generator into \cref{extended:thm:survival_hazard}.
\end{proof}


\section{Earlier exact examples}\label{earlier:section:06}
These five short examples are retained from the earlier human edition. Their numerical parameters are independent of the later visual guide's worked models.
\begin{example}[One transfer between two packets]\label{earlier:record:026}\label{earlier:ex:two_packet_transfer}
Let
\[
A=
\begin{pmatrix}
a&0\\
c&b
\end{pmatrix}
\]
and let the two coordinate axes be the packet spaces. If \(a\neq b\), then
\[
\e^{tA}
=
\begin{pmatrix}
\e^{at}&0\\
\displaystyle c\frac{\e^{at}-\e^{bt}}{a-b}&\e^{bt}
\end{pmatrix}.
\]
The lower left block is produced by one transfer from the first packet to the second. Its expansion begins with \(ct\), in agreement with \cref{earlier:cor:minimum_transfer_order}. When \(a=b\), the same block is \(ct\e^{at}\).
\end{example}

\begin{example}[A four state lumpable chain]\label{earlier:record:027}\label{earlier:ex:four_state_lumping}
Partition four states into \(B_1=\{1,2\}\) and \(B_2=\{3,4\}\). For positive rates \(a,b,u,v\), set
\[
Q=
\begin{pmatrix}
-u-a&u&b/2&b/2\\
u&-u-a&b/2&b/2\\
a/2&a/2&-v-b&v\\
a/2&a/2&v&-v-b
\end{pmatrix}.
\]
The rates \(u\) and \(v\) move inside the coarse blocks. Every state in \(B_1\) leaves for \(B_2\) at total rate \(a\), and every state in \(B_2\) leaves for \(B_1\) at total rate \(b\). Hence
\[
\overline Q=
\begin{pmatrix}
-a&b\\
a&-b
\end{pmatrix}
\]
and \(CQ=\overline QC\). The internal rates disappear from the coarse generator even though they remain present in the fine evolution.
\end{example}

\begin{example}[A visible likelihood loss]\label{earlier:record:028}\label{earlier:ex:bayes_information_loss}
Let \(B_1=\{1,2\}\), let \(B_2=\{3\}\), and take
\[
p=(3/10,2/10,5/10),
\qquad
L=(1,3,2).
\]
The coarse prior is \((1/2,1/2)\). The exact block likelihoods are
\[
\ell_1(p)=\frac{(3/10)+3(2/10)}{1/2}=\frac95,
\qquad
\ell_2(p)=2.
\]
Thus the exact coarse posterior is \((9/19,10/19)\). If the first block is replaced by the fixed representative value \(2\), the approximate coarse posterior is \((1/2,1/2)\). The \(\ell^1\) defect is \(1/19\), while \cref{earlier:thm:information_loss} gives the valid upper bound \(1/10\).
\end{example}

\begin{example}[One hidden mode]\label{earlier:record:029}\label{earlier:ex:one_hidden_mode}
Consider
\[
A=
\begin{pmatrix}
-a&c\\
d&-b
\end{pmatrix},
\qquad
a,b>0,
\]
and retain the first coordinate. The memory kernel and self energy are
\[
\Kcal(t)=cd\e^{-bt},
\qquad
\Sigma(z)=\frac{cd}{z+b}.
\]
The retained coordinate has an exact time local closure for every retained initial value precisely when \(cd=0\). When \(cd\neq0\), one auxiliary scalar state realizes the memory exactly.
\end{example}

\begin{example}[Two sequential phases]\label{earlier:record:030}\label{earlier:ex:two_phase_survival}
Let
\[
A=
\begin{pmatrix}
-\lambda&0\\
\lambda&-\mu
\end{pmatrix},
\qquad
p_0=
\begin{pmatrix}1\\0\end{pmatrix},
\]
where \(\lambda\neq\mu\) are positive. The first phase moves to the second at rate \(\lambda\), and the second is absorbed at rate \(\mu\). Direct calculation gives
\[
S(t)=
\frac{\mu\e^{-\lambda t}-\lambda\e^{-\mu t}}{\mu-\lambda},
\]
\[
f(t)=
\frac{\lambda\mu}{\mu-\lambda}
\left(\e^{-\lambda t}-\e^{-\mu t}\right).
\]
The ratio \(f(t)/S(t)\) is generally not constant. A time local Markov model on two transient phases can therefore produce a nonconstant hazard after the phase coordinate is hidden.
\end{example}

\section{Holomorphic operator substitution}\label{migration:section:01}\label{migration:sec:operator_substitution}

This section receives the operator-substitution result deliberately moved out of Paper IV. Fix an integer $N\geq1$ and the entire scalar function
\[
\Dfun_N(z)=\sum_{m=0}^{\infty}\frac{(-1)^m z^{Nm}}{(Nm)!}.
\]
More generally, the result applies to any specified holomorphic scalar function that is nonzero on the declared domain. On a simply connected zero-free domain $U$, choose $g=\log\Dfun_N$ once and define $\Sfun_N^{(r)}(z)=\exp(-r g(z))$ for complex $r$. The bounded Banach-space result below uses classical holomorphic functional calculus; it is separate from the finite canonical-jet calculus formalized in Part I.

\begin{standardfact}[Holomorphic functional calculus]\label{migration:record:001}\label{migration:sf:holomorphic_calculus}
Let $A$ be a bounded operator on a complex Banach space.  If $f$ is holomorphic on a neighborhood of $\Spec(A)$, then
\[
f(A)
=
\frac1{2\pi i}\int_\Gamma f(z)(zI-A)^{-1}\,dz
\]
for an admissible positively oriented contour $\Gamma$ lying in the holomorphy domain and the resolvent set, with winding number one about the spectrum and zero outside the domain. The calculus preserves sums, products, and composition.  We use the standard form from \cite{DunfordSchwartz1958}; an accessible proof is given in \cite[Appendix A]{TaylorBanachAlgebras}.
\end{standardfact}

\begin{proposition}[Operator reciprocal family]\label{migration:record:002}\label{migration:prop:operator_reciprocal}
Let $A$ be bounded and suppose $\Spec(A)$ lies in a simply connected region on which $\Dfun_N$ has no zeros.  Fix a holomorphic branch $g=\log\Dfun_N$ there.  Then
\[
\Sfun_N^{(r)}(A)
=
\exp\bigl(-r g(A)\bigr)
=
\bigl(\Dfun_N^{-r}\bigr)(A)
\]
is well defined and
\[
\Sfun_N^{(r)}(A)\Sfun_N^{(s)}(A)
=
\Sfun_N^{(r+s)}(A).
\]
If $(P_i)$ is a finite complete family of bounded idempotents, then
\[
P_i f(A)P_j
=
\frac1{2\pi i}
\int_\Gamma
f(z)P_i(zI-A)^{-1}P_j\,dz
\]
for every holomorphic $f$ on a neighborhood of the spectrum.
\end{proposition}

\begin{proof}
The spectral mapping theorem makes $\Dfun_N(A)$ invertible.  The product identity follows from multiplicativity of the functional calculus and the fixed logarithm.  Applying $P_i$ and $P_j$ to the Cauchy integral formula gives the packet formula.
\end{proof}


\clearpage\bibliographystyle{alpha}\bibliography{references}
\end{document}
